% AI-written, AI cross-reviewed; no human sign-off; not a claim that RH is proved
% RH Notes blueprint source excerpts for Z4.4
% Verbatim excerpts in the notation of the parent blueprint. The labelled editions and source ranges appear below.
% See the unit YAML for the current proof status. Earlier-round passages retain their original wording.
\section*{Z4.4. Twisted periods on the selected chain}

% Parent source SHA-256: 1af64922b98ed2310744411b1f18b25cc426c97220646321471836e1a3a409d6
% blueprint/ch_P.tex lines 12626 through 12642; r6:review:Z4.4
\paragraph{Z4.4. The period on the selected chain (round-six assessment).}\phantomsection\label{r6:review:Z4.4} \emph{Original argument on p.~\pageref{sec:r5-z44}.} 
The two comparisons were open at round six. Their subsequent assembly is on page~\pageref{r12:assembly}.
The first obligation constructs one comparison $C$ with its standard form and level identifications. Its proof uses
\path{exists_standardFormComparison_symmetricLevel_members},
\path{hermitianColumnChoices_of_formsLevelsClasses},
\path{stepZeroHermitianLevel_uniformiser_kills_values_alt},
\path{exists_hermitian_data_carT_of_std},
\path{hermitianColumns_generators_baseLattices_functionals}
and \path{spaces_in_coordinates}. These identify forms, levels, representatives, stabilizers, generators and functionals across the same columns. The second obligation constructs the twisted record over this $C$ and the given tuple $\xi$. It uses \path{trueColumnChoice_exists_chainData_with_basePoints},
\path{tupleChains_levelStructure_compatible_at_levelPrime} and
\path{tupleLayerCompatibilities_of_X_prime}. The period-to-inner-product identity, Petersson class expansion, point-space and rational identifications, and the coordinate map used when restricting to a class are supplied by that record. See \path{colsTiedSym_std}, \path{colsTwPTiedS_std}, \path{twistedPeriodAdmissible_of_rowInt} (\path{Development/AbelianVarietyTheory/Proofs/IwasawaDecomposition/ExistsStandardColumnComparisonsTrueColumns.lean:20--72}, \path{Development/AbelianVarietyTheory/Proofs/IwasawaDecomposition/TrueColumnChoiceExistsStandardChainData.lean:45--56}, \path{Development/AbelianVarietyTheory/Proofs/GlobalInterpolation/AdmissiblePeriodTwistOfIntegralRows.lean:14--21}).

The transported duality used below is proved on page~\pageref{r13:weyl-dual}.
The standard Weyl and dual Weyl lattices are exact duals for the Fock pairing. In the proof used here, dual frames transport this fact at each finite place and intersection recovers the integral pairing. The standard duality is the input at both localizations in \path{nabDelDual}; \path{pairBDualK} specializes the resulting theorem. See \path{std_dual}, \path{nabDelDual}, \path{pairBDualK} (\path{Development/AbelianVarietyTheory/Proofs/FoundationProofs/StandardWeylLatticesFockDual.lean:237--245}, \path{Development/AbelianVarietyTheory/Proofs/FoundationProofs/TransportedWeylLatticeHermitianDuality.lean:278--309}, \path{Development/AbelianVarietyTheory/Proofs/GlobalInterpolation/IsClassSumPairingFormsLat.lean:181--197}).

With these inputs, the normalized numerator is integral and the single scalar $y=\sum_c b_c/h_c$ gives the period. For each good $w$ above $p$ this scalar is zero or $v_w(y)\ge-e_w v_p(D\Gamma_u)$; only the stabilizer orders $h_c\mid D\Gamma_u$ enter the denominator. The branching calculation concerns the sibling point-integrality theorem. The level-prime and generator-prime comparisons use the period assertion downstream. The verified class summation and bound in Section~\ref{sec:r5-z44} stand conditionally on the two constructions and standard duality above.
See \path{pair_integral}, \path{perDecTw_of_at}, \path{perTwAt_of_perDecTw}, \path{row_aPerTw_of_overCols} (\path{Development/AbelianVarietyTheory/Proofs/GlobalInterpolation/ClassPairingIntegralFormRowIsIntegral.lean:20--69}, \path{Development/AbelianVarietyTheory/Proofs/GlobalInterpolation/TwistedPeriodDecompositionOfClasswiseTwistedPair.lean:141--152}, \path{Development/AbelianVarietyTheory/Proofs/GlobalInterpolation/AdmissiblePeriodTwistsReductionPlaceTuples.lean:12--37}, \path{Development/AbelianVarietyTheory/Proofs/GlobalInterpolation/TwistedPeriodsIntegralAtBundlePoints.lean:11--19}).

% Parent source SHA-256: 509893882390ca545e5542b088ec432ed3c1ed009152eeca67e7fc44894dd956
% blueprint/ch_P_r11.tex lines 434 through 1083; r11:z44
\subsection{Z4.4. Which comparisons supply the twisted period?}\label{r11:z44}

Fix a column system $I$ on $(B,X)$ with $I.\mathrm{Ok}$, pins, laws $l$,
later data $Z$, the zero law $l_0$, and integral rows. Fix also the choice
of primes and the tied tuple $t_p$ with $t_p.\xi=I.\xi$. The comparison
must use this tuple for every divisor datum $Y$. Write $D=I.\mathrm{rec7}(l,Z)$.
The first two calls of \path{twistedPeriodAdmissible_of_rowInt}
(\path{Development/AbelianVarietyTheory/Proofs/GlobalInterpolation/AdmissiblePeriodTwistOfIntegralRows.lean:14--21})
select $C$ first, then $t$ over $C$. Throughout this subsection, a bare
Iwasawa source filename denotes a file in
\path{Development/AbelianVarietyTheory/Proofs/IwasawaDecomposition/}.

\paragraph{The nine calls in the dependency graph.}
The path to each of the first six passes through
\path{exists_standardColumnComparisons_trueColumns} and \path{colsTiedSym_std}
(\path{ExistsStandardColumnComparisonsTrueColumns.lean:20--72}).
The remaining three pass through \path{trueColumnChoice_exists_standardChainData}
and \path{colsTwPTiedS_std} (\path{TrueColumnChoiceExistsStandardChainData.lean:45--56}).
Counts below are distinct source lines in the inclusive core cone of each
supplier, with its theorem-only count in parentheses. They overlap. The
names in rows 1--5 and 7--9 belong to \path{GlobalTheory.T}, row 6 to
\path{GlobalTheory.Mem.Col2.TC}.
\begin{enumerate}
\item \path{exists_standardFormComparison_symmetricLevel_members},
53,609 (43,240). \path{ExistsStandardFormComparisonSymmetricLevelMemberStatement.lean:8--9}.
\item \path{hermitianColumnChoices_of_formsLevelsClasses},
19,178 (10,327). \path{HermitianColumnChoicesOfFormsLevelsClassesStatement.lean:8--9}.
\item \path{stepZeroHermitianLevel_uniformiser_kills_values_alt},
16,725 (7,960). \path{StepZeroHermitianLevelUniformiserKillsValuesStatement.lean:9--10}.
\item \path{exists_hermitian_data_carT_of_std},
21,901 (12,388). \path{HermitianDataAtLevelPrimeStatement.lean:8--9}.
\item \path{hermitianColumns_generators_baseLattices_functionals},
17,252 (8,398). \path{HermitianColumnsGeneratorsBaseLatticesFunctionalStatement.lean:8--9}.
\item \path{spaces_in_coordinates},
14,238 (6,767). \path{TrueColumnSpacesStandardInFrameCoords.lean:90--101}.
\item \path{trueColumnChoice_exists_chainData_with_basePoints},
62,711 (50,996). \path{TrueColumnChoiceExistsChainDataWithBasePointsStatement.lean:10--18}.
\item \path{tupleChains_levelStructure_compatible_at_levelPrime},
39,525 (29,242). \path{TupleChainsLevelStructureCompatibleAtLevelPrimeStatement.lean:8--9}.
\item \path{tupleLayerCompatibilities_of_X_prime},
34,625 (24,494). \path{TupleLayerCompatibilitiesOfXPrimeStatement.lean:9--16}.
\end{enumerate}

\paragraph{Three implications used to construct the column comparison.}
In the following three proofs the initial data are $I$, its witnesses
\path{Ok} and \path{Pins}, the laws $l$, and the later data $Z$. Write
$D=I.\mathrm{rec7}(l,Z)$, $m=D.R.m$. In references with a bare filename
in this subsection the source directory is
\path{Development/AbelianVarietyTheory/Proofs/IwasawaDecomposition/}.

\paragraph{The frame coordinates (supplier 6).}
At every $i\le m$ the Gram matrix is $1$ and the coordinate image of the
standard self-dual lattice is $\mathcal O_E^{N(i)}$. The primed Gram
matrix is also $1$. The embedding of the primed space extends a vector
of length $N(m)-1$ by a zero last coordinate. For each $u,a,i,c$, the
coordinates of $I.\mathrm{latAt}(u,a,i,c)$ give the lattice of the
chosen decorated representative. There is a frame equivalence $e_Y$
with the prescribed \path{IdY} and
$\mathrm{de}_i(z,\tau)=\det((e_Yz)_\tau)$.
Indeed \path{Ok.spaces} identifies the spaces with the positive standard
models. Evaluating their form on coordinate vectors gives the Kronecker
delta; those vectors span their prescribed integral lattice. Expansion
in this basis gives the asserted integrality criterion in both
directions and provides the witnesses to \path{StdTied}. The pin
\path{iX} sends the source vector of index $a$ to the target vector
of the same index. Its sum with coefficients $x_a$ therefore has
exactly the claimed coordinates, including the missing last index.
The selected lattice is defined by pulling the representative back by
the coordinate equivalence, whose forward image recovers it. The two
frame types now coincide definitionally; the identity equivalence has
both required properties. See \path{TC.spaces_in_coordinates} and
\path{Crd.incl_co},
\path{TrueColumnSpacesStandardInFrameCoords.lean:14--101}.

\paragraph{Annihilation at step zero (supplier 3).}
Fix a word $u$, a tied forms comparison $F$ with \path{StdF}, and a
level $L$ with the same carriers as the step-zero forms level.
For every $w\mid T$ and every $a\in S=\mathrm{intV}(E,w)$ whose image
is the image of $\pi=F.\mathrm{atT}.\mathrm{piU}$, one can choose a
single replacement $L'$ such that $aL'.\mathrm{fM}(w)=0$.
The base lattice, local and global level groups, representatives,
stabilizer embeddings, splitting maps and generator frames are those of $L$.
See \path{T.stepZeroHermitianLevel_uniformiser_kills_values_alt},
\path{StepZeroHermitianLevelUniformiserKillsValuesStatement.lean:9--10}.

Here is the local lattice calculation used in this choice. A unitary
local map $y$ carrying the base lattice to the representative lattice
also carries their duals. At residue characteristic $p$, localize
$\mathcal O_E$ at the multiplicative set of rational integers prime to $p$ (that is,
invert those integers). The resulting semilocal Dedekind ring is a PID:
for a nonzero ideal prescribe the valuations of a generator at its
finitely many maximal ideals by the Chinese remainder theorem, choosing
a residue in $\mathfrak q^e\setminus\mathfrak q^{e+1}$ at each.
A full lattice thus has a local basis matrix $B$. The conditions on a
dual vector $x$ are exactly that the entries of $B^*Jx$ be integral,
by pairing with the columns of $B$. This description is the localization
of the global dual: denominators away from $p$ can be cleared
simultaneously on the finitely many generators of the global lattice.
Their images are units at $w$, so after extending to $S$ the columns
of $(B^*J)^{-1}$ still generate precisely the vectors pairing
integrally with the extended lattice. Conjugation preserves $S$ since
it preserves the two generating rings. Now
$h(v,y^{-1}z)=h(yv,z)$ gives the asserted identification in both
directions. See \path{LocB.exists_isLocBasis},
\path{Formalism/LocalToGlobal/GlobalTheory/GL6/PBwLocGen.lean:248},
\path{latS_of_locBasis},
\path{Formalism/SharpRestriction/GrowthBounds/IntegralWeylLatticeSpan/LatSOfLocBasis.lean:117},
\path{isLocBasis_dual}, \path{mem_latS_dualL}, \path{dual_move},
\path{UnitaryCompletionMapsDualLattices.lean:52--188}, and
\path{MovePair.ld_loc}, \path{ClassPairMovedOfMultipliesLocalFunctional.lean:96--106}.

Choose $v$ corresponding to $w$ and a class $c$ above the identity
double coset using surjectivity of the representative map. The class
functional is nonzero. Choose an input where it is nonzero and lift it
by the surjective residue projection to a vector $z$ in the class dual.
The preceding calculation puts $t=y_c^{-1}z_w$ in the local base dual.
The tie between the two functionals, via the given additive isomorphism
of their targets, makes $f_T(t)\ne0$. The residue of $\pi$ is zero;
linearity of the projection with respect to reduction now gives
$f_T(y_c^{-1}(\pi z)_w)=0$. This vector is $at$, since localizing
$\pi$ gives $a$ and scalar multiplication commutes with $y_c^{-1}$.
Thus $a$ kills $f_T(t)$. The additive target has prime order $T$, so
this nonzero value generates it and $a$ kills it all.
Transport the forms-level functional and its value module to the
rank indexing of $L$. The rank, Gram and base-lattice identifications
transport its nonvanishing, vanishing on the base lattice, scalar
annihilation and stabilizer identity; after identifying the indices
these are literally the old identities. Replace only the value module,
functional and their certificates. This gives the promised single $L'$
before quantifying over $w,a$.
See \path{fun_transport}, \path{pi_kills},
\path{stepZeroHermitianLevel_with_formsFunctional},
\path{StepZeroHermitianLevelUniformiserKillsValues.lean:29--54},
\path{StepZeroHermitianLevelUniformiserKillsValues.lean:60--162}.

\paragraph{Generators, base lattices and functionals (supplier 5).}
Assume column data $H$, tied forms $F_u$, prime data $A_u$,
\path{StdF}, \path{StdLvC} for each word, \path{CarH}, \path{CarT},
and that the local images of $\pi=A_u.\mathrm{piU}$ annihilate
$(H.\mathrm{lev0}\,u).\mathrm{fM}(w)$ for every $w\mid T$.
The conclusions are the cofactor identity, the base-lattice identities,
a family of residue-field functionals, and level-zero data $X_0$
with $(X_0u).\mathrm{piG}=X.\mathrm{piG}$ for all $u$.
See \path{T.hermitianColumns_generators_baseLattices_functionals},
\path{HermitianColumnsGeneratorsBaseLatticesFunctionalStatement.lean:8--9}
and \path{TC.herm_rest_tied},
\path{HermitianColumnsGeneratorsBaseLatticesFunctional.lean:16--33}.

First $\pi=X.\mathrm{piT}(u)$ by the two standard comparisons.
Hence the pinned cofactor identity is
$D.\mathrm{cOf}(u)\pi_u=\pi_{u_0}$. In the frame of supplier 6 the
step-zero lattice and each chart lattice have the description
$L=\mathcal O_E^{n-1}\oplus\pi\mathcal O_E$ in positive rank:
at an intermediate index use \path{StdF.mid} and \path{CarH.carC};
at the last index use the unique chart and \path{StdF.ends} with
\path{CarH.carM}; at zero use \path{CarH.car0}. The norm equation
makes $\pi\ne0$. Testing the Hermitian pairing with $e_j$ for
$j<n$ and with $\pi e_n$ for the last coordinate gives
$L^\vee=\mathcal O_E^{n-1}\oplus\bar\pi^{-1}\mathcal O_E$.
Thus $L+\bar\pi L^\vee=\mathcal O_E^n$: the first $n-1$
coordinates belong to $L$ and the last is supplied by the second
summand. The rank-zero equality is immediate. These are precisely
the two parts of \path{SymmetricBaseCompatibility}.
See \path{TC.base_sup_dual_eq}, \path{TC.base_step_zero_tied},
\path{BaseLatticeStepZeroAddConjGeneratorDualSelfDual.lean:12--169}
and \path{ChartBase.sup_dual_eq}, \path{TC.chart_base_sup_dual},
\path{ChartBaseLatticeSupDualSelfDual.lean:22--105},
\path{ChartBaseLatticeSupDualSelfDual.lean:158--211}.

For $1\le i<m$, reindex the local dual by the rank equality, apply the
forms-level functional $f_T$, and follow it by the given additive
identification of its target with the residue field. For a class
$c$ and a dual vector $z$, if $y_c t=z_w$ in coordinates, then the
reindexed $t$ equals the inverse of the reindexed $y_c$ applied to the
reindexed $z_w$. The coordinate pairing is the original form with
permuted indices, so $z$ belongs to the required coordinate dual.
The projection identity \path{CarT.prC} and the given tie
\path{fvC} identify the composite value with
$\mathrm{fv}_{i,v,c}(A_u.\mathrm{pr}_{i,v,c}z)$. At $i=m$ the chart
is unique; the same calculation uses \path{CarT.pr} and the terminal
forms tie \path{fv}. This constructs the full family
\path{SymmetricResidueFieldLevels}.
See \path{FvChart.fvTie_of_car},
\path{ChartLevelFunctionalsIntermediateStepsTrueColumn.lean:47--132},
and \path{TC.fvTie_top},
\path{LevelFunctionalChartLastStepClasses.lean:29--63},
\path{LevelFunctionalChartLastStepClasses.lean:74--122}.

At zero there are two functionals to compare on the same local dual
$M$, one from $H$ and one transported from $F$. Their targets are
additive groups of prime order $T$. Both functionals vanish on the
local base lattice $L_w$. For $n>0$ the coordinate calculation above
gives $M=L_w+S e$ with $e=e_n/\bar\pi$: test a global dual
vector on $\pi e_n$, subtract its resulting last coefficient times
$e$, and extend by $S$-linearity. In rank zero take $e=0$.
Both targets are killed by $\pi$, the first by hypothesis and the
second by the nonzero-value argument of supplier 3. A common scalar
action is needed here, since two modules of the same order alone do
not identify their actions. For a scalar from the completed base
integers the actions on these cyclic groups are integers $k_1,k_2$.
The differences between the scalar and each integer annihilate a
nonzero module and so are nonunits in that local ring. Therefore
$k_1-k_2$ lies in its maximal ideal, hence is zero modulo $T$.
For $b\in\mathcal O_E$ set
$d=\pi+\bar\pi$ and $c=b\bar\pi+\bar b\pi$; both descend to
base integers by fixedness and integrality. The residue of $d$ is
nonzero, since the residue kills $\pi$ and does not kill $\bar\pi$.
On each module $db$ acts as $c$, since $\pi$ acts as zero. The
common residue of $d$ is invertible, so $b$ acts on both by the
same integer modulo $T$. The two rings just considered generate $S$,
giving the same conclusion for every $s\in S$.
See \path{localDual_eq_add_span}, \path{LocalDualEqAddSpan.lean:9--144},
\path{same_int_of_local}, \path{acts_by_same_integer},
\path{ActsBySameInteger.lean:71--127}, \path{ActsBySameInteger.lean:180--253},
and \path{stepZeroLevel_functional_residueField},
\path{StepZeroLevelTypePrimeOrientationFramesFunctiona.lean:187--252}.

For each $t\in M$, write $t=se+v$ with $v\in L_w$.
Its two values are therefore the same integer multiple of the two
values at $e$. Both values at $e$ are nonzero since the two maps are
nonzero. They generate their respective targets; the multiple vanishes
on either side precisely when $T$ divides it. The two kernels coincide.
The quotient by this kernel supplies an additive isomorphism carrying
the first functional to the second. Composing with the forms-side
identification gives the required level-zero residue identification
and its value formula for each class and every permitted lift. Choosing
these identifications at all $w$, take $\mathrm{piG}=X.\mathrm{piG}$;
its field-membership, norm and split orientation are those of
the standard forms. The generators' frame-lattice identities reindex
by \path{CarH.car0} and \path{CarH.fr0}; the $\pi$-annihilation is the
given hypothesis. Together these are the fields of $X_0$.
See \path{TC.exists_equiv_of_cyclic}, \path{TC.colLevel_functional},
\path{TC.intV_coe_eq_inclE_of_T_mem},
\path{StepZeroLevelTypePrimeOrientationFramesFunctiona.lean:36--81},
\path{StepZeroLevelTypePrimeOrientationFramesFunctiona.lean:118--184},
\path{StepZeroLevelTypePrimeOrientationFramesFunctiona.lean:255--302}.

\paragraph{Supplier 2. Recovering Hermitian columns from a supplied family of forms.}
Suppose the finite chart sets $\mathrm{Ch}_{u,i}$, the maps from classes,
and the forms $F_u$ with $I.\mathrm{StdF}(u,F_u)$ have been supplied,
with a unique chart at the top for each $u$. We construct the initial
levels, higher levels with class identifications, places, and residue-field
maps giving $\mathrm{CarH}(H,F)$. In this paragraph $E/F_0$ denotes the
quadratic CM extension, and $p=F_u.\mathrm{atT}.\pi_U$, $q=\bar p$,
$T=pq$ (a rational prime). \path{hermColsCh_of_forms},
\path{HermitianColumnChoicesOfFormsLevelsClasses.lean:68--82};
\path{AtTTied}, \path{Formalism/LocalToGlobal/HermitianForms/PrkSmul.lean:25--48}.

First identify the fields at $T$. For a selected place $v$, let $r_v$ be
the given surjection from $\mathcal O_E$; $r_v(p)=0$, $r_v(q)\ne0$.
Every $r_v(y)$ is the image of an element of $\mathcal O_{F_0}$. Indeed $p+q$ and
$yq+\bar y p$ are fixed integral elements, and the kernel of $r_v$ on
$\mathcal O_{F_0}$ is $v$. Multiply the second by an inverse of the first
modulo $v$. Its image is $r_v(y)$. Consequently the source residue field
in the column identifies with $\mathcal O_{F_0}/v$ in the required
direction, inverse to reduction from $\mathcal O_{F_0}$.
\path{exists_intF}, \path{res_algebraMap_surjective}, \path{exists_fv_equiv},
\path{ResidueFieldEquivRealSubfieldQuotientTrueColumns.lean:19--57}.

For the chart level, reindex its rank by the given frame. In that frame the
Gram matrix is $1$, and the base lattice has coordinates
$\mathcal O_E^{n-1}\oplus p\mathcal O_E$. This follows by testing each
coordinate and dividing the last by $p\ne0$. At a prime over either
generator, a unitary lattice stabilizer acting identically modulo that
generator belongs to the chart level: lift each column of the identity
using the chart's split frame. The difference after applying the
stabilizer is a multiple of the generator in that lattice. The resulting
matrix reduces to the identity, which belongs to the given residue
subgroup. Away from generators and $T$ the chart already carries the
prescribed hyperspecial subgroup.
\path{level_of_chart}, \path{baseL_eq},
\path{HermitianStepLevelOfChartFormsLevel.lean:98--130,136--169};
\path{mem_Uloc_of_cong}, \path{ValNatCastLtOne.lean:41--83};
\path{spaces_in_coordinates}, \path{TrueColumnSpacesStandardInFrameCoords.lean:90--101}.

Here is the discriminant calculation at $v\mid T$. The kernels $P_v$ of
the $f$ given residue maps from $\mathcal O_E$ are distinct by their
restrictions to the real field and have norm $T$ by surjectivity and
the given field sizes. All contain $p$ and exclude $q$.
The conjugate primes contain $q$ and exclude $p$. The norm of $(T)$ is
$T^{2f}$ by the degree of $E$, so these $2f$ primes account for $(T)$
with exponent one each. In particular the completed quadratic algebra
over $F_{0,v}$ is split, its integral algebra is $R\times R$ with
$R=\mathcal O_{F_{0,v}}$, its involution interchanges the factors, and
$T$ generates the maximal ideal of $R$. Equivalently, localization and
completion in the factorization just obtained split the integral algebra
by the Chinese remainder theorem; each local field extension has degree
one. \path{notMem_sq}, \path{finrank_E}, \path{unram},
\path{LevelPrimeUnramifiedUniformizer.lean:35--74,75--110};
\path{T_notMem_sq}, \path{dvd_of_val_lt_one},
\path{LevelPrimeDvdOfAdicValuationLtOne.lean:21--79}.

Put $A=R\times R$ and $L=A^{n-1}\oplus pA$.
The dual is $L^\#=A^{n-1}\oplus q^{-1}A$, by testing on the coordinate
generators. With $w=q^{-1}e_n$, every $y\in L^\#$ is $aw+y_0$ with
$a\in A$, $y_0\in L$. The chart functional kills $L$ and is nonzero;
thus its value $m$ on $w$ is nonzero. For $x$ in the chart level,
$x$ preserves both lattices, preserves the functional and is unitary.
Writing $d=x_{nn}$ gives $d\in A$ and $(d-1)m=0$ by evaluating on $xw$.
For $j<n$ one has $x_{jn}/q\in A$. The last diagonal entry of
$\bar x^{\,t}x=1$ then gives $\bar d d\equiv1\pmod{TA}$, since
$q\bar q=T$. The additive target has prime order $T$, so $m$ generates
it. Its scalar action is a unital character
$R\times R\longrightarrow\mathbf F_T$ and factors through one residue
projection: the two orthogonal idempotents sum to one, and the
homomorphism from the retained local factor kills $T$. Hence the
corresponding component of $d$ is $1$ modulo $T$. Their product is
$1$ modulo $T$, so both components are $1$. Now $(x-1)w\in L$, including
the last coordinate $(d-1)/q\in pA$. By the decomposition of $y$ this
gives $(x-1)L^\#\subset L$. Conversely, for a unitary stabilizer of
$L$ with this containment, $x$ preserves $L^\#$ in both directions
by duality, and the functional on $L^\#/L$ is fixed. These are the chart's
level conditions at $T$. This argument uses the particular $F_u$ at hand.
\path{uloc_T_iff}, \path{HermitianStepLevelOfChartFormsLevel.lean:78--95};
\path{uloc_acts_trivially_on_dual_quotient},
\path{LevelActsTriviallyOnDiscriminantGroup.lean:50--186};
\path{eq_one_add_of_fixes}, \path{UnitaryScalarEqOneModOfFixesSimpleModule.lean:13--68};
\path{uloc_of_triv_disc}, \path{UlocOfTrivDisc.lean:10--70}.

Use these same lattices, local groups and global group for the Hermitian
step level. Its dual inclusion extends from generators to their span by
linearity. Its class representatives and stabilizer injections are those
of the chart, transported by the equality of ranks. After this
identification the Gram matrices agree entrywise; the class bijection,
adelic lattice criterion and stabilizer image criterion hold by composing
their old versions with the class-index bijection. For $1\le i<m$
use the fiber over each chart label. At $i=m$ the top chart is unique,
so its fiber maps bijectively onto all classes.
\path{exists_herm_level}, \path{HermitianStepLevelOfChartFormsLevel.lean:35--69};
\path{classTie_of_reindex}, \path{classTie_of_chart},
\path{HermitianClassConditionOfChartClasses.lean:29--62,66--81};
\path{level_classes_of}, \path{HermitianColumnChoicesOfFormsLevelsClasses.lean:13--65}.

At step zero reindex the existing endpoint data by $N(0)=k$.
At each generator prime the lattice is self-dual: $T L^\#\subset L$
and $T$ is a unit there. In the split frame the two local factors are
dual, and a unitary element acts by inverse transpose on the second;
preserving the lattice is therefore equivalent to integrality and
invertibility in the first lattice frame. At this endpoint the blocks
in the residue test are singletons and cover all indices. Its equations
say precisely that the reduction is upper unipotent, giving the required
Iwahori-type subgroup. A place above the first generator exists by
lying over. If a rational stabilizer is $y_c k y_c^{-1}$ adelically,
$k$ preserves the completed base lattice everywhere, and the rational
stabilizer preserves the class lattice: a full finite lattice is the
intersection in its rational space of its localizations (test the
annihilator of the class of a vector modulo the lattice).
This proves the extra lattice-stability field. The split maps and frames
come from the endpoint record itself. With the places $F_u.\mathrm{atT}.eV$
and the residue identifications above, all seven fields of
$\mathrm{CarH}$ follow. This proves supplier 2 from its stated premises.
\path{exists_reindex}, \path{inJ_iff}, \path{mem_Uloc_iff_iwU},
\path{stab_class_lattice}, \path{exists_colLevel},
\path{ColumnLevelOfStepZeroFormsLevelTrueColumns.lean:32--67,74--139,143--178,201--256};
\path{hermColsCh_of_forms}, \path{HermitianColumnChoicesOfFormsLevelsClasses.lean:68--82}.

\paragraph{Supplier 1. Shared reference charts (partial).}
For each $u,i$ we need finite chart labels $\mathrm{Ch}_{u,i}$ and a map
from classes, unique label at $i=m$, and for every $u$ a form record
which is simultaneously $\mathrm{StdF}$ and $\mathrm{StdLvC}$. Its generator
data $G$ must be identical for every word and its character at the
reference word must satisfy \path{CharacterCompatible}.
For $i\le m$, the reference chart label is the image of
$c\mapsto (c.1,\mathrm{cosetIn}(\mathrm{rpI}(c)))$; this is finite as
a surjective image of the finite orbit-class set. Two classes in a fiber
have the same index and reduction coset. At $i=0,m$ the index is unique
and the coset independent of the class ($i+1\ne m$, using $m\ge2$);
a base decorated point gives an element. For $i>m$ take a singleton,
since no level condition is requested there.
\path{exists_charts}, \path{cos_of_ne},
\path{ExistsMemberStepCharts.lean:393--404,521--566};
\path{formsTiedSym_std}, \path{ExistsStandardFormComparisonSymmetricLevelMember.lean:52--79}.

The missing assertion already at \path{exists_charts} is that each fiber
admits a \path{LevelChart} on its specified stabilizers and lattices,
with $I.\mathrm{ReadsK}$ and representative
$(K_{\mathrm{chOf}(c)}.y(c))=I.y_{\rm std}(u,i,c)$.
The construction of a chart on a fiber requires that compatible
decoration and reduction data come from a point, and that those data
determine the relevant point. Neither follows from the choice of a
representative of the fiber. The cone of \path{exists_charts} has
27,859 core lines (19,468 theorem lines).
Even given those charts the full supplier also requires compatible
values, inner products, diamond and up/down operations, the
penultimate-level group, and the period/evaluation laws on the same
forms. The assembly sets every $F_u.gl=G$, and at an endpoint replaces
the family by its unique chart. In an orthonormal frame its terminal
vector has only the last entry $p$, so the Hermitian pairing with $x$
is $\bar p x_n$, as required for the functional comparison. This
calculation does not give the missing simultaneous chart and form data.
\emph{Statement only for the full supplier; its source proof is not reproduced.}
\path{chart_of}, \path{exists_charts}, \path{ExistsMemberStepCharts.lean:458--486,521--566};
\path{hform_last}, \path{formsTiedSym_std},
\path{ExistsStandardFormComparisonSymmetricLevelMember.lean:33--45,80--145,151--243}.

\paragraph{Supplier 4. Stabilizer comparison and the remaining moves (partial).}
Given the same $H,F$ with $\mathrm{CarH}$ and standard forms and levels,
one needs $A$ at all words with $\mathrm{CarT}(H,F,A)$ and both the
adelic and stabilizer compatibilities. We can prove the latter for
\emph{any} $A$ satisfying $\mathrm{CarT}$. In coordinates write $g$
for the matrix of $\mathrm{tilG}_c$ and $M$ for the matrix of the
chain move carrying the point to its representative. Test evaluation
at the all-identity frame $y$ with the indicator of that single frame.
The two evaluations are multiplied by units. If $M^{-1}gy\ne y$,
one evaluation is zero and the other nonzero. Cancel the invertible
frame at one embedding; injectivity of the embedding gives $g=M$.
The class stabilizer image is defined by conjugation by the move:
if $e$ is the extended stabilizer, its image at the chosen representative
has matrix $M^{-1}eM$. Conjugating back gives $e$. The initial-coordinate
and last-coordinate laws for $e$ say that on $(x,0)$ it acts by the
hyperplane stabilizer and on $(0,\ldots,0,p)$ it is the identity.
The endpoint and hyperplane class comparisons in $\mathrm{CarH}$ and
$\mathrm{CarT}$ identify those matrices with the two actual linear
actions. Injectivity of frame coordinates proves both required identities.
\path{ga_eq_gam_mv}, \path{eq_of_frMul}, \path{crdN_stabM},
\path{crdP_stabP}, \path{gam_ext}, \path{stabMap_tied},
\path{StabiliserMapHyperplaneEqConjClassMap.lean:27--40,50--156}.

The existence of $A$ itself also uses the seen-prime condition, the
discriminant quotient, middle projections, three adelic moves and the
comparison of hyperplane classes and lattices. For example the precise
downward assertion asks, for every $u,i,v,c$ and residue functional
$g\in(\mathrm{lk}_{i,v,c})^*$, for one $\gamma$ such that if $i\le m$ and
$g\ne0$ then both $\mathrm{DnCl}(F_u.\mathrm{atT},F_u.\mathrm{lv},i,v,c,g,\gamma)$
and $\mathrm{AdOf}(H,u,i,c,\mathrm{dn}(i,v,c,g),\gamma,v)$ hold.
Outside this case take $1$. At an endpoint use the move in $F_u.\mathrm{dnN}$;
at an interior step the move in $F_u.\mathrm{dnC}$ is indexed by the
chart of $c$ and requires the chart-to-global projection identities.
In both cases one must compare the moved lattices away from the two
primes over $v$, compare the localized dual at the selected prime,
and identify the functionals at $v$ and at the other places.
The required downward supplier has 18,740 core lines (9,681 theorem
lines); the stabilizer argument above does not supply these comparisons.
\emph{Statement only for the full existence assertion.}
\path{dnCl_of_tie}, \path{dn_move_exists},
\path{ExistsHermitianDownMoveAtLevelPrimePlace.lean:15--63,69--86};
\path{hermitian_data_at_level_prime}, \path{HermitianDataAtLevelPrime.lean:18--62}.



\paragraph{The chain over a fixed comparison (supplier 7).}
Use the columns $H$, member forms $F$, and the hypotheses
$\mathrm{StdF}$, $\mathrm{StdLvC}$, $\mathrm{CarH}$ of
\nolinkurl{trueColumnChoice_exists_chainData_with_basePoints}
(\path{TrueColumnChoiceExistsChainDataWithBasePointsStatement.lean:10--18}).
The finite set $Q_a$ is the image of the primes of $I.Q(a)$ under
contraction. Hence membership is exactly $\mathrm{UnderQ}(a,v)$.
At $i=0$ use the class map transported by $e\mathrm{Cls}_0$.
For $0<i<m$ the reference tuple uses the given member chart with an enlarged
bad set. A nonreference tuple requires a chart trivialized at $Q_a$.
At $i=m$ one must first choose a chart with the top-point property, then
overwrite the $m$th entry of the family of stage choices by this chart.
Thus the same top chart witnesses both the reading of $U,\chi$, lattice,
representatives and stabilizers and the following condition for all $d,g',g$:
if $g'$ represents $d$ and $g$ extends $g'$ fixing the specified $z_0$, then
$g k_\beta$ represents $\widetilde d$.
At $i>m$ only the class map is needed. At $m-1$, choose a coupling level and
a character agreeing with every chart and agreeing with the reference
character on elements whose components at $Q_a$ are trivial. The common
base lattice follows by taking a class at $m-1$ and comparing the two
membership tests on each vector. With this data the chain structures
$T,J,W$ use the same stage family, and $\mathrm{StdW}.\mathrm{topk}$
supplies the displayed $k$ without a second choice.
See \nolinkurl{twStep_mid}, \nolinkurl{twStep_top}, \nolinkurl{twPack_ex},
\nolinkurl{tw_chains} (\path{TrueColumnChoiceTwistedChainsFormsClasses.lean:153--245},
\path{TrueColumnChoiceTwistedChainsFormsClasses.lean:272--319},
\path{TrueColumnChoiceTwistedChainsFormsClasses.lean:464--489}) and
\nolinkurl{tw_ties_with_points}
(\path{TrueColumnChoiceExistsChainDataWithBasePoints.lean:13--23}).

\emph{Open in this construction.} In particular
\nolinkurl{chart_lastStep_twistBasePoint}
(\path{ChartLastStepTwistBasePointStatement.lean:8--9}, 27,250 full core lines)
must produce the top chart and this single $k$, preserve $U_v$ outside
$Q_a$, identify characters on elements trivial at $Q_a$, and give a
finite exponent for the new character and the top-base condition.
The source takes $k=k_{\rm ref}t$ with a unitary adelic diagonal $t$,
constructs a framed chart, and uses block extension to identify the lifted
point. Once the character is exhibited as the product of the reference and
twist characters of exponents $m_1,m_2>0$, its exponent divides $m_1m_2$
by raising both factors to that power. This last argument gives no
construction of $t$ or of the chart. See \nolinkurl{chart_twm} and
\nolinkurl{root_top} (\path{ChartLastStepTwistBasePoint.lean:195--209},
\path{ChartLastStepTwistBasePoint.lean:223--336}).
The intermediate charts \nolinkurl{charts_tuple_trivQ}
(\path{ChartsTupleTrivQ.lean:86--101}, 40,240 lines) and the penultimate
cosets \nolinkurl{exists_penultimateLevel_couplingCosets}
(\path{ExistsPenultimateLevelCouplingCosetsStatement.lean:8--9},
32,279 lines) also remain inputs to this argument.

\paragraph{Level structures (supplier 8).}
There is a direct proof of the reference-tuple comparison at every stage.
Assume the identification of primes is the identity and that the inverse
residue-field identification fixes the residue of each integer.
Each residue has an integral representative. Applying the field
identification to the assumed equality on representatives shows it fixes
every residue. Take the identity on both additive quotients. Scalar
multiplication, functional evaluation and the up maps then agree. The two
down maps agree because their linear functionals agree on every vector.
For dilation the two units have equal residues and hence are equal as
units. This supplies also the transported maps on frames and the
basis condition. This is \nolinkurl{lvSameMu0}
(\path{TupleChainsLevelStructureCompatibleAtLevelPrime.lean:44--90}).

For a nonreference tuple and $1\le i\le m$, given a chart reading, identical
lattice and residue maps, compatible frame and prime equivalences, and
functionals killed by the chosen uniformizer, the missing assertion produces
a level $L$ whose projections on $N$ and $\mathrm{dn}N_v$ are the
respective quotient maps. The same $L$ has a functional comparison and
both up and down naturality. The assembly applies the comparisons to a
single choice of projections. In addition it requires the basis comparison
and the stage-zero identification for all tuples.
See \nolinkurl{exists_lvTied_of_std} and \nolinkurl{tw_lv_tied_and_same}
(\path{TupleChainsLevelStructureCompatibleAtLevelPrime.lean:17--40},
\path{TupleChainsLevelStructureCompatibleAtLevelPrime.lean:94--134}).
The two naturality propositions alone use \nolinkurl{upNatTie}
(\path{InvSmulMem.lean:281--329}, 30,447 full core lines) and
\nolinkurl{dnNatTieG_cols}
(\path{ChartTiedLevelChainHost.lean:48--58}, 20,761).
The reference-tuple argument does not settle this existence assertion.

\paragraph{The point and its period (supplier 9).}
For a fixed $u,a$ write $S$ for the primed step, $r_c$ for its chosen
representative, $\lambda$ for the shifted weight, and $E^*(P)$ for
evaluation of the coefficientwise conjugate of $P$ on frames, including
the prescribed determinant factor. Evaluation is injective. Indeed that
factor is nonzero on invertible tuples. If a polynomial vanishes there,
multiplying it by the product of the determinants gives a polynomial zero
everywhere over $\mathbb C$. The determinant product is a nonzero
polynomial, so the original polynomial is zero. An admissible value at
a class has an $E^*$-preimage in the complex weight space: its generators
do, addition preserves preimages, and scalar multiplication by $r$ acts on
preimages by $\bar r$. See \nolinkurl{evFr_injective},
\nolinkurl{fnOf_injective} (\path{Formalism/LocalToGlobal/Columns/PolyFn.lean:259--291})
and \nolinkurl{exists_cev} (\path{CevInjective.lean:17--39}).

Define $(j_M G)_c=E^*(G(r_c))$. This is conjugate linear. For a group
element $g$, the polynomial action is $\sigma(\chi(g))^{-1}\rho(g)$.
Its character has finite order, so conjugating its inverse gives
$\sigma(\chi(g))$. The unitary matrix of $g$ has inverse its conjugate
transpose. The embedding identity \path{fields_ok.conjE} gives
$\mathrm{emb}_\tau(\bar x)=\overline{\mathrm{emb}_\tau(x)}$ for each
coefficient of this matrix (\path{CevInjective.lean:219--225}).
Substitution in the frame action proves that $E^*$ intertwines
the two actions. In particular the evaluated values at $r_c$ are fixed
under its stabilizer. Conversely, lift each stabilizer-invariant point
value to $P_c$ by the preceding paragraph. Equivariance and injectivity
of $E^*$ put $P_c$ in the fixed subspace of the stabilizer. Extend to
$g r_c$ by the polynomial action of $g$ on $P_c$. Two choices of $g$
differ on the right by a stabilizer, so this definition is independent
of the choice. It gives the required form, and the same construction
shows $j_M$ is injective since the source uses the full group at each
class. See \nolinkurl{cev_act}, \nolinkurl{jM_injective},
\nolinkurl{jM_mem}, \nolinkurl{jM_onto}
(\path{CevInjective.lean:74--97}, \path{CevInjective.lean:106--181},
\path{CevInjective.lean:219--225}).

To compute its period, recover a polynomial from evaluation in the
identity frame: for $E^*(P)$ the recovered polynomial is $\bar P$.
In the normalized coefficient pairing the conjugation on that polynomial
cancels the coefficientwise conjugation, leaving the bilinear contraction
with $P$ and the same normalization. Thus for each class the Hermitian
contribution with first entry $E^*(G(r_c))$ equals the corresponding
bilinear contribution to the period. Both sums have weight
$1/|\operatorname{Stab}(r_c)|$ in that class, proving
$\operatorname{period}(w,G)=\operatorname{ipP}(j_M G,w)$ for every $w,G$.
For a nonzero admissible point at least one class polynomial is nonzero.
Its self-pairing has positive numerator (a sum of squared absolute
coefficients times positive monomial factorials) and positive denominator
(the corresponding highest-weight norms). All remaining class
contributions are nonnegative. See \nolinkurl{polOf_cev},
\nolinkurl{nfAll_conj}, \nolinkurl{ipVal_cev}, \nolinkurl{ipVal_pos},
\nolinkurl{period_eq} and \nolinkurl{re_petersson_self_pos}
(\path{CevInjective.lean:41--72}, \path{CevInjective.lean:184--211});
\nolinkurl{ps_fockF_pos}
(\path{Formalism/LocalToGlobal/Columns/CoefficientsModP/PairingNondegenerateHolds.lean:51--70}).
This proves the period identity for the fixed chain without appealing to
the existence of a standard point package.

The rationality test also descends class by class. A rational embedded
point gives an admissible complex value and each component is the
evaluation in the identity frame of a polynomial in the rational span.
Conversely write such a polynomial at $c$ as $\sigma(P_c)$ with $P_c$
in the weight space over the coefficient field. This description of the
rational span follows by expressing its generators as products of integral
highest-weight words, mapping their coefficients into $\mathbb C$, and
taking sums and coefficient-field multiples in both directions.
Stabilizer invariance of the complex values and injectivity of evaluation
imply that each $P_c$ is stabilizer invariant over the coefficient field.
These $P_c$ therefore form a rational point mapping to the given row.
The weight conventions match since $t=2k+1$ gives
$m(t+1)/2=m(k+1)$, and the common determinant shift cancels in the
normalized weight. See \nolinkurl{ptEmb_iff_rat}
(\path{TupleChainFormsRationalIffFrameValues.lean:16--30}),
\nolinkurl{exists_embF_iff}, \nolinkurl{rat_iff}, \nolinkurl{VV_shW}
(\path{ExistsEmbFIff.lean:29--41}, \path{ExistsEmbFIff.lean:101--112},
\path{ExistsEmbFIff.lean:157--167}), and \nolinkurl{mem_RatSpanG_iff}
(\path{Formalism/LocalToGlobal/Columns/TiedValRat.lean:154--156}).

\emph{Open for supplier 9.} For each $(u,a)$ obtain one primed class
chart $K$ with the reference base lattice, finite-exponent character,
$\mathrm{ReadsP}$ for the chosen tuple, the reference $U_v$ at generator
places and character agreement on common elements trivial at $Q_a$
(\nolinkurl{tuplePrimedSpaceClasses_level_and_comparison},
\path{TuplePrimedSpaceClassesLevelAndComparisonStatement.lean:9--10},
24,746 full core lines). On this $K$, at every allowed $w$, point
membership in $\mathrm{SP}(w)$ must be equivalent to having an integral
polynomial in \emph{every} local basis at \emph{every} class, evaluated in
the inverse basis (\nolinkurl{tuple_integralForms_row_place_by_place},
\path{TupleIntegralFormsRowPlaceByPlaceStatement.lean:11--12},
19,117 lines). Finally, with the top $k$ above one must construct the
tilted datum with the specified $z_0$ and prove for each pair of
transported frames that its evaluation is
$\chi(k_a(c))^{-1}\phi(\widetilde c,s')$
(\nolinkurl{til_tw}, \path{ExtMMulVecFix.lean:253--274}, 17,886 lines).
These inputs enter \nolinkurl{forms_tw}
(\path{FormsXTupleChainOfLevelReading.lean:81--91}) and
\nolinkurl{tupleLayer_compatibility}
(\path{TupleLayerCompatibilitiesOfXPrime.lean:24--46}). They still
prevent the simultaneous construction of the standard point and
evaluation data over the fixed level $L$.

\paragraph{Assembly over the chosen column.}
Given suppliers 1--6, choose $F$, replace the level zero of supplier 2
using supplier 3, and keep its base lattice, level, class representatives,
stabilizers, splittings and frames by the seven supplied equalities.
Its vanishing condition therefore applies to the same $\pi_U$ as in
supplier 4 after rewriting the uniformizer equality in $\mathrm{CarT}$.
Supplier 5 now gives the cofactor, base, residue data, and $X_0$ with
$X_0(u).\pi_G=X.\pi_G$ at every word. Use the first component of
supplier 6 for the standard spaces. Constancy of the initial and
hyperplane frames follows by writing each frame as the specified
standard frame. These are exactly the fields of $C$, including
$\mathrm{StdC}$ and $\mathrm{StdLvC}$.
\path{colsTiedSym_std}, \path{ExistsStandardColumnComparisonsTrueColumns.lean:20--72}.

Now $t_p.\xi=I.\xi$ identifies $t_p.\mathrm{InQ}'(a,v)$ with
$I.\mathrm{UnderQ}(a,v)$: both say that some place above $v$ ramifies for
that character. In one direction take the place above the index in
the selected set; in the other take the designated place over $v$,
using correctness of the selection. Supplier 7 is called on $C$'s
own $H,F$; supplier 8 fixes the level on its $T,J,W$, and supplier 9
takes this very level and its base-point witness. The resulting $t$
therefore has compatible period, coordinate and evaluation fields.
Only after those existence inputs and lattice duality can the already
checked class sum be applied to the selected period.
\path{underQ_iff}, \path{carH_of_cols}, \path{colsTwPTiedS_std},
\path{TrueColumnChoiceExistsStandardChainData.lean:19--56};
\path{twistedPeriodAdmissible_of_rowInt},
\path{Development/AbelianVarietyTheory/Proofs/GlobalInterpolation/AdmissiblePeriodTwistOfIntegralRows.lean:14--21}.

% blueprint/ch_P_r11.tex lines 1085 through 1178; r11:fock
\subsection{The integral contraction needed for one class}\label{r11:fock}
The step from two integral values to an integral numerator uses
\path{GlobalTheory.PairAtClass.pair_integral} (\path{ClassPairingIntegralFormRowIsIntegral.lean}:49--55) through
\path{GlobalTheory.PeriodClassSum.pairBDual} and \path{pairBDualK}
(\path{IsClassSumPairingFormsLat.lean}:181--197).
Here is the integrality direction, including the normalization factor.
It explains which part of exact Weyl duality the class numerator consumes.

Fix a single dominant weight with natural partition $\nu$. Write $h$ for its
highest weight polynomial and $N=(h,h)_F>0$ for its rational Fock norm
squared. A generator of the integral Weyl lattice is $Ah$, with $A$ a word
in divided row operators. If $F$ belongs to the integral dual Weyl lattice,
put $H=A^t F$, where transposition swaps row indices and reverses word order.
Each divided operator takes the coefficient of a power of $t$ after an
integral row substitution. Thus $H$ has integer coefficients and stays
in the dual Weyl space: row substitution preserves the column weight and
commutes with the column unipotent substitutions, coefficient by coefficient.
With $I$ the identity frame (rectangular if necessary), the formula is
\[
             (Ah,F)_F/N=(h,H)_F/N=H(I)\in\mathbb Z.
\]
For the first equality expand in monomials. For different rows, moving
$k_j$ copies of variable $x$ to $y$ in column $j$ gives the coefficient
$\binom{a_j}{k_j}$ if their initial exponents are $a_j,b_j$.
Its product with the final monomial's factorial norm agrees with that
for the reverse move, since
$\binom{a_j}{k_j}(a_j-k_j)!(b_j+k_j)!
= a_j!b_j!\binom{b_j+k_j}{k_j}$.
Multiply over columns and sum over the transfers of the prescribed total
degree. For equal rows the substitution acts diagonally in the monomial
basis. Iteration gives the reversed word.

For the second equality one may test $H$ on the spanning products of
initial-column minors of shape $\nu$. Each product is row homogeneous.
The Fock pairing with $h$ kills every row weight other than that of $h$;
evaluation at $I$ also kills such products. For a nonzero product with
the same row counts as $h$, row one must occur in every nonempty minor:
its required count equals the number of such minors and no minor can
repeat that row. Inductively row $s$ must occur in each minor of size
at least $s$: all shorter minors are already filled by the previous rows.
The product is therefore $h$ up to the signs of the row permutations.
Since $h(I)=1$, the two functionals agree there as well. This proves
the formula and integral linear combinations preserve it.
Source comparisons: \path{GlobalTheory.WeylDuality.nablaEqDualWeylOf},
\path{ApplyWordMemDualWeyl.lean}:110--149 (the forward direction);
\path{standardMonomialSpanDomain}, \path{Formalism/GradedWeylDivisibility/FinalC.lean}:37--38 (the
spanning input); \path{GlobalTheory.FockHighestWeight.fockQ_hw},
\path{LowerTriOfInterlaced.lean}:105--169; \path{fock_space_row_divided_adjoint},
\path{FockQRowDivAdj.lean}:169--173; \path{fockQ_applyWord},
\path{FockQWordRowActAdj.lean}:8--18.

Over $\mathcal O_K$ the standard lattices are spanned by the scalar
extensions of the bideterminants and words (\path{dualWeyl_toK},
\path{weyl_toK}, \path{StandardWeylLatticesFockDual.lean}:129--139).
On pure tensors the normalized pairing is the product of the single-weight
pairings. The denominator is exactly $\prod_\tau N_\tau$; star on the
first scalar preserves $\mathcal O_K$. Additivity now gives the desired
integrality on the two standard lattices (\path{std_dual}, same file,
237--307). For transported lattices choose dual local frames at $\bar w,w$
with $B^*GB'=1$. In these coordinates the pairing reduces to the standard
one, denominators from the first argument are conjugated and are units at
$w$, and denominators from the second are units there as well. An element
of $K$ integral at every finite $w$ belongs to $\mathcal O_K$. Specialize
to trivial star and $G=1$ and use symmetry to interchange the two arguments
of the bilinear pairing at the class.
This last transport uses the localization and dual-frame identifications
of \path{GlobalTheory.TransDual.dual_of_local} and \path{nabDelDual}
(\path{TransportedWeylLatticeHermitianDuality.lean}:218--309).

\emph{Boundary of this argument.} The stronger source assertion characterizes
both integral lattices by testing against the other, for all $x\in VV$.
The converse over $\mathbb Z$ needs the integral straightening step:
a rational dual Weyl polynomial whose lower-triangular monomial
coefficients are integers has an integral polynomial lift. In
\path{GlobalTheory.WeylInt.lowerTriIntegral}
(\path{LowerTriIntegral.lean}:41--82) one subtracts an integral multiple
of a monic bideterminant with the same leading exponent. The new polynomial
has smaller leading exponent and retains the lower-triangular integrality
test. Existence of the requisite tableau for an interlaced leading
exponent and the identification of the localized transported lattices are
inputs here. The full duality declaration, \path{standard_weyl_lattices_fock_dual}
(\path{StandardWeylLatticesFockDualStatement.lean}:14--15), has an
inclusive cone of 5,730 core lines. The one-direction argument above
supplies no additional expert-unit completion.

\paragraph{Independent review queue at the end of round 11.}
For Z1.2 the review candidate is the finite-count implication
\path{UpperTheory.V.Z.J.rankTwo} (\path{RankTwo.lean}:177--229).
For Z4.4 the candidates are suppliers 2, 3, 5 and 6 in the numbered list,
each with its supplied form or chart hypotheses. These five candidates
were submitted for review. The round 11 submission left the existence
and compatibility statements in suppliers 1, 4, 7, 8 and 9 open.
The resulting corrections and the subsequent assembly are on
pp.~\pageref{r12:trust}, \pageref{r12:assembly}.

% Parent source SHA-256: 3c960212a7e63fd06c5daea8f4f37c62338eb8a3a6dfd9389bb3bf4387e70253
% blueprint/ch_P_r12.tex lines 224 through 584; r12:weyl, r13:weyl-dual
\subsection{The two standard Weyl lattices are exact duals}\label{r12:weyl}

Let $K$ be the number field with involution, $R=\mathcal O_K$,
$\iota$ a finite index set, $\lambda_\tau$ dominant weights, and
$x\in VV(K,\lambda)$. For the normalized pairing $P$ at the identity
Gram matrices the two claims are
\[
  P(\Delta,x)\subseteq R\ \Longleftrightarrow\ x\in\nabla,
  \qquad P(\nabla,x)\subseteq R\ \Longleftrightarrow\ x\in\Delta .
\]
Here $\Delta=\mathrm{delStd}$ and $\nabla=\mathrm{nabStd}$.
The condition $x\in VV$ is essential for coordinate testing.
The target is \path{GlobalTheory.T.standard_weyl_lattices_fock_dual}
(\path{StandardWeylLatticesFockDualStatement.lean:14--15}). Its inclusive
cone has 5,730 core lines. References to \path{StatementA.lean},
\path{StatementC.lean}, \path{StatementE.lean}, and the tableau files
in this subsection are in \path{Formalism/GradedWeylDivisibility/}
and its \path{FockSpaceTableaux/} subdirectory. Other bare filenames refer
to \path{Development/AbelianVarietyTheory/Proofs/FoundationProofs/}.

\paragraph{Recovering a polynomial coefficient by testing a Weyl word.}
Start over $\mathbb Q$ with one partition $\nu$, $n\le r$, and the
space $V_\nu$ of polynomials of column multidegree $\nu$ invariant
under adding an earlier column to a later column, with an indeterminate
scalar. Let $h$ be the product of the prescribed powers of leading
principal minors and $N=(h,h)_F>0$. Write
$B=(\cdot,\cdot)_F/N$. The lattice $\nabla_{\mathbb Z}$ consists of
the integral polynomials in $V_\nu$ (the defining coefficient identities
descend along $\mathbb Z\hookrightarrow\mathbb Q$);
$\Delta_{\mathbb Z}$ is the integer span of the divided row words $Ah$.
See \path{dualWeyl}, \path{rowGen}, \path{hw}, \path{applyWord},
\path{weyl} (\path{StatementA.lean:24--73}) and \path{bDual}
(\path{StatementE.lean:21--45}).

Row substitutions preserve column totals and commute with column substitutions. Operators on a polynomial occur in the reverse order to successive row substitutions on its frame. A divided row operation selects the coefficient of $t^k$ in a row
addition. Its adjoint for the factorial monomial pairing interchanges
the two rows. On a monomial this is the factorial identity in
\S\ref{r11:fock}; reversing the order of a product gives $A^t$.
The evaluation identity $B(h,H)=H(I)$ on $V_\nu$ follows by checking
products of initial-column minors. Each such product is row homogeneous.
A product with the row weight of $h$ must, successively, use row $i$ in
every minor of size at least $i$ (one-based rows). Such a product is
$\pm h$ or zero. All other row weights vanish both in the pairing against
$h$ and on evaluation at $I$. The spanning assertion for these products
also follows by the leading-term elimination below, over $\mathbb Q$.
This proves $B(A^t h,F)=(AF)(I)$; for integral $F$ the right side is
integral. See \path{fockQ_hw} (\path{LowerTriOfInterlaced.lean:135--169}),
\path{fockQ_applyWord} (\path{FockQWordRowActAdj.lean:8--18}) and
\path{nablaEqDualWeylOf} (\path{ApplyWordMemDualWeyl.lean:110--180}).

Suppose conversely that $F\in V_\nu$ and $B(\Delta_{\mathbb Z},F)$
is integral. For a lower triangular exponent $e=(e_{as})$ with the
prescribed column totals, set the diagonal variables to one, the
strictly upper variables to zero and leave all strict lower variables
independent. The coefficient at $e$ below the diagonal in the resulting
polynomial is $[X^e]F$: the column totals recover the missing diagonal
powers. This same coefficient is $(A_e F)(I)$. Indeed start at $I$
and install below-diagonal variables column by column from right to left.
For position $(a,s)$ add an indeterminate multiple of row $s$ to row
$a$ and extract its power $e_{as}$ in the substitution on $F$.
At this stage row $s$ is still the $s$th coordinate row: all previously
installed entries have column at least $s$ and row strictly larger
than their column. The substitution therefore installs precisely
the variable at $(a,s)$; repeated coefficient extractions recover the
coefficient claimed. This gives the source's ordered word under the
fold-right convention. Exponents of the wrong column totals have zero
coefficient. Thus every lower triangular coefficient is an integer,
by applying the preceding identity to $A_e^th$.
See \path{coeff_single}, \path{map_rowGen}, \path{eval_applyWord},
\path{wordOf_eq}, \path{coeffWord}
(\path{UnitriMonomialOfLowerTri.lean:36--58,60--88,93--131},
\path{UnitriMonomialOfLowerTri.lean:141--156,221--236}).

\paragraph{Leading exponents and integral straightening.}
Take the lexicographic order comparing first by column, then by row.
The leading exponent $m$ of a nonzero $F\in V_\nu$ satisfies, at
each pair of consecutive columns $l,j$ and each row $a$,
\[
       \sum_{b\le a}m_{bj}\ \le\ \sum_{b<a}m_{bl}.                 \tag{*}
\]
Here is the two-column calculation explaining the strict endpoint.
Keep the component with the same exponents outside those columns as
$m$ and with the same combined degree in the two entries of each row.
This retains the leading term, and substitution
$X_{bj}\mapsto X_{bj}+tX_{bl}$ preserves this grading. The transformed
component has $t$-degree zero since the original polynomial is invariant.

More generally, give this transformed polynomial a degree allowance
$s$. At the first remaining row write its variables as $X,Y$ and its
degree there as $d$. Expand $\sum_{k=0}^K X^kY^{d-k}w_k$,
where $w_K\ne0$ and the $w_k$ are independent of $X,Y$.
The column order selects $k=K$ for the leading term. Put $e=d-K$.
After substitution split off $(Y+tX)^e$; the residual factor is
\[
       L=\sum_{k=0}^{K} X^k(Y+tX)^{K-k}\,w_k(t),
\]
where $w_k(t)$ denotes substitution in all other rows.
Substitution is injective (put $t=0$); the polynomial ring is a
domain, so $e\le s$ and $\deg_t L\le s-e$. Putting $X=1$, $Y=-t$
recovers $w_K(t)$, whose $t$-degree is at most $s-e+K$ since $L$
has total $X,Y$-degree $K$. Now remove that row from $w_K$.
At the first row $e\le s$; at every subsequent cut, the remaining
right-column total is at most $s-e+K$ plus the preceding left-column
total in the remaining rows, by induction. Restoring the first row
adds $e$ on the left of the desired inequality and $K$ on the right.
This proves (*) with allowance $s$, with the empty row set as base
case. Apply it to the initial allowance zero. Sources:
\path{InterlacedAt}, \path{pairW} (\path{StatementC.lean:34--46}),
\path{leadInterlaced_of} (\path{LeadInterlaced.lean:30--62}),
\path{lctwo_main} (\path{TwoColumn.lean:47--208}) and
\path{peelDegBound_of} (\path{PeelDegBound.lean:119--181}).

In column $s$ make a sorted list $L_s$ containing row $a$ with
multiplicity $m_{as}$. Its length is $\nu_s$. By (*) the $k$th entry
of $L_s$, whenever present, is strictly larger than the $k$th entry
of $L_{s-1}$: apply (*) with $a$ that $k$th entry. For each $k$ take
the initial-column minor with rows the $k$th entries of these lists,
in column order. Their columns form an initial segment since $\nu$ is
a partition. The row indices strictly increase. The diagonal
permutation gives the unique leading monomial of this minor, with
coefficient one: the first column where another permutation differs
uses a larger row index. The product $D_m$ of the minors is therefore
monic with leading exponent $m$, column totals $\nu$, and integer
coefficients. The minors are invariant under the permitted column
additions. This constructs the required tableau, including the empty
product for the zero exponent. In particular $m$ is lower triangular,
since the $s$th row in any strictly increasing list of row indices
is at least $s$.
See \path{sSYTExists_of} (\path{SSYTExists.lean:11--62}),
\path{colLexInitMinor_of} (\path{ColLexInitMinor.lean:22--52}) and
\path{columnLexBideterminant_of_initialMinor}
(\path{ColumnLexBideterminant.lean:11--27}).

Subtract $([X^m]F)D_m$. The leading exponent falls and every lower
triangular coefficient stays integral. There are finitely many
exponents at the fixed column totals. Iteration ends at zero and
expresses $F$ as an integer sum of the $D_m$. With arbitrary coefficients
the same elimination over a field proves the spanning assertion used
above. Hence $\Delta_{\mathbb Z}^{\#}=\nabla_{\mathbb Z}$ in
$V_\nu$, where $\#$ denotes the dual for $B$.
This proves the converse in \path{nablaEqDualWeylOf}; see also
\path{GlobalTheory.WeylInt.lowerTriIntegral}
(\path{LowerTriIntegral.lean:41--82}).

\paragraph{The other dual and the number-field coefficients.}
The lattice $\Delta_{\mathbb Z}$ is contained in the finite free
coefficient module at the fixed multidegree. It therefore has a finite
integer basis by \path{Submodule.basisOfPid}
(\path{Mathlib/LinearAlgebra/FreeModule/PID.lean}:300--302).
Its rational span is all of $V_\nu$: a vector
orthogonal to it has every lower triangular coefficient zero by word
extraction, whereas a nonzero vector has a lower triangular leading
exponent. The restricted Fock form is positive definite over $\mathbb Q$.
Let $d_i$ be an integer basis of $\Delta_{\mathbb Z}$. Its
$B$-dual basis $b_i$ exists over $\mathbb Q$ by inverting the Gram
matrix. Testing a rational vector against the $d_i$ reads off precisely
its coordinates on the $b_i$. Thus $\nabla_{\mathbb Z}$ has basis
$b_i$. Testing on the $b_i$ gives back $\bigoplus_i\mathbb Z d_i$.
This proves both dualities over $\mathbb Z$ without any unimodularity
assumption on an individual Gram matrix. See \path{dual_dual_of_isLattice},
\path{normB_nondeg}, \path{weylEqDualNablaOf}
(\path{ApplyWordMemDualWeyl.lean:182--217},
\path{ApplyWordMemDualWeyl.lean:298--324}).

Apply this construction at each $\tau$ to $\nu=\lambda_\tau.\mathrm{natW}$.
The divided words commute with coefficient extension, so the $d_i$
span the Weyl lattice over $R$. The leading-term argument over the
domain $R$ expresses each integral dual Weyl polynomial as an $R$-sum
of the same integral bideterminants. They belong to the integral span
of the $b_i$ already over $\mathbb Z$, and the $b_i$ retain their
defining invariance over $R$. Thus these two rational bases also generate
the two lattices defined over $R$. Tensor over the finite set of weights.
On products of polynomials in the disjoint blocks, the factorial monomial
form factors: the monomial factorials and the coefficients both multiply.
The denominator of $P$ is the product of the single-block highest-weight
norms, so the two product bases pair by $\delta_{IJ}$, in either order.
They span the prescribed coefficient space over $K$ and their integral
spans are the two standard lattices. Indeed the product polynomials
give the pure tensors generating those spaces, and duality of the
displayed families gives independence. The involution fixes the rational
basis coefficients and preserves $R$ because it preserves integrality
over $\mathbb Z$. Testing $x$ on either first-argument basis recovers
its second-argument coordinates on the opposite basis; integral sums
of basis elements pair integrally. This proves both biconditionals.
See \path{dualWeyl_toK}, \path{weyl_toK}, \path{exists_dual_bases}, and
\path{std_dual} (\path{StandardWeylLatticesFockDual.lean:129--139},
\path{StandardWeylLatticesFockDual.lean:145--235},
\path{StandardWeylLatticesFockDual.lean:237--307}).

\paragraph{Localization of the transported lattices.}
The preceding argument concerns the standard models. At a finite
prime $w$ write $R_w$ for the localization of $R$ inside $K$.
Localize either standard lattice $T$. Its defining integrality test is
independent of the local frame: two bases of the same full lattice
differ by a matrix in $\mathrm{GL}_n(R_w)$.
For $\nabla$ an integral basis substitution with integral inverse
preserves integral polynomials, column weights and right column
invariance. For $\Delta$, an elementary row addition acts by the sum
of its divided row operations multiplied by the corresponding powers
of its integral parameter, a finite sum on each polynomial. Each
divided operator prepends a letter to a Weyl generator. Diagonal units
act on a word by a unit since $h$ is row homogeneous and every divided
operation shifts the row weight by its prescribed amount. These
matrices generate $\mathrm{GL}_n(R_w)$ by elimination with unit
pivots and row interchanges (the latter are products of elementary
matrices and diagonal units). The determinant twist is a unit even
at a negative exponent. This gives frame independence for both
lattices, including products of blocks.
See \path{isBasisAt_inv_mul}, \path{indep_of_stable},
\path{nabIndep_and_delIndep}
(\path{NabIndepAndDelIndepLocalConditions.lean}:21--43) and the two
stability inputs \path{GlobalTheory.T.polyAction_mem_locP_nabStd_of_isBasisAt}
and \path{GlobalTheory.T.polyAction_mem_locP_delStd_of_isBasisAt}
(\path{ActPMemLocPNabStdOfIsBasisAtStatement.lean}:14--15;
\path{ActPMemLocPDelStdOfIsBasisAtStatement.lean}:14--15).

Given a full lattice $M_\tau$ in each block, choose a local basis
matrix $B_{w,\tau}$ over $K$ at $w$: $R_w$ is a DVR and the finitely
generated torsion-free localization is free of full rank
(\path{GlobalTheory.LatBasisAt.basisAtEx},
\path{LatticeHasBasisAtEveryFinitePlace.lean}:66--105). Write
$T_w(B_w)=\mathrm{act}(B_w)(T\otimes_R R_w)$ and
$P_T=\bigcap_w T_w(B_w)$ in the rational coefficient space.
The equality $(P_T)_w=T_w(B_w)$ merits a denominator argument.
The inclusion from the left follows by localizing the defining test.
For $x\in T_w(B_w)$, away from finitely many primes all $M_\tau$
have identity frame (clear denominators in both inclusions between
two full finitely generated lattices); also $x\in T_v(1)$ away from
finitely many primes since $T$ spans the coefficient space. At each
exceptional $v\ne w$, choose $k_v$ so $v^{k_v}x\subseteq T_v(B_v)$
by clearing local denominators. Choose $s_v\in v^{k_v}\setminus w$,
which exists because distinct maximal ideals are comaximal.
The product of the finitely many $s_v$ belongs to $R\setminus w$
and sends $x$ into every $T_v(B_v)$. Thus $x\in(P_T)_w$. With a
single global frame the same intersection gives precisely the transported
global lattice, since if $x\notin T$, the proper ideal
$\{a\in R:ax\in T\}$ lies in some maximal ideal; there $x\notin T_w$.
Finally, a single local test places $x$ in the $K$-span of the
transported lattice and hence in $VV$, which is stable under changing
frame (row substitution commutes with the defining right-column identities).
These prove the three assertions in
\path{GlobalTheory.T.transportedLattice_localization_free_le_coefficientSpace}
(\path{TransportedLatticeLocalizationFreeLeCoefficientSStatement.lean:14--15});
see \path{exists_isBasisAt_one}, \path{transLat_colLat}, \path{locP_transLat}
and \path{nabDelLoc_free_leVV}
(\path{TransportedLatticeLocalizationFreeLeCoefficientS.lean:91--125},
\path{TransportedLatticeLocalizationFreeLeCoefficientS.lean:139--234}).

\paragraph{Hermitian duality after transport.}\phantomsection\label{r13:weyl-dual}
Let $R=\mathcal O_K$ for a number field $K$ with involution. Fix a
finite set $\iota$, a rank $n$, dominant weights $\lambda_\tau$,
Hermitian matrices $G_\tau$ with nonzero determinant and full
$R$-lattices $M_\tau\subset K^n$. Write $V=VV(K,\lambda)$,
$P_G=\mathrm{pairK}(\lambda,G)$ and
$M_\tau^\#=\{x:\bar y^tG_\tau x\in R\text{ for all }y\in M_\tau\}$.
Use $\nabla(M),\Delta(M)$ for the two transported lattices restricted
to $V$. The two conclusions in \path{NabDelDual}
(\path{CoefModLemmas.lean}:143--148) are equalities of $R$-submodules of $V$:
\[
 \nabla(M)^{\vee,P_G}=\Delta(M^\#),\qquad
 \Delta(M)^{\vee,P_G}=\nabla(M^\#).
\]
The dual on the left always tests the \emph{first} argument against
the specified lattice. We prove both equalities at each finite place and
then intersect.

For a sesquilinear $h$ and a finitely generated $R$-module $T$ in its
first argument, suppose $x\in V$ and the module $D$ is characterized
on $V$ by $x\in D\iff h(T,x)\subset R$. For each finite prime $w$,
\[
 x\in D_w\quad\Longleftrightarrow\quad h(T_{\bar w},x)\subset R_w .
\]
If $s x\in D$ for $s\notin w$, and $t y\in T$ for $t\notin\bar w$,
then $h(ty,sx)=\bar t s h(y,x)\in R$, with $\bar t s\notin w$.
Conversely the local pairing condition holds on generators $y_j$ of
$T$. Choose $s_j\notin w$ clearing $h(y_j,x)$ and put $s=\prod_j s_j$.
Then $h(T,sx)\subset R$ by conjugate linearity and $sx\in V$,
so $sx\in D$. The empty product also works. For every $R$-submodule
$A$ inside a $K$-space one has $A=\bigcap_w A_w$: if $x\notin A$,
the proper ideal $\{s:sx\in A\}$ is contained in a maximal ideal
$w$, at which $x\notin A_w$. This applies also to $R\subset K$.
Hence $h(T,x)\subset R$ is equivalent to all the tests
$h(T_{\bar w},x)\subset R_w$; this last equivalence follows simply
by localizing each scalar and also holds without finite generation.
See \path{mem_of_forall_locP}, \path{pair_denom},
\path{mem_localization_dual} and \path{pair_integral_iff}
(\path{TransportedWeylLatticeHermitianDuality.lean}:29--44,70--99,101--123).

Fix $w$. Each $M_\tau^\#$ is full and finitely generated:
$y\mapsto\bar y^t G_\tau$ takes generators of $M_\tau$ to a
finite spanning family of row vectors over $K$. Its $R$-span is
therefore between $dR^n$ and $d^{-1}R^n$ for some nonzero rational
integer $d$ (clear the coordinates and a spanning inverse).
Taking integral dot-product duals gives the same sandwich for
$M_\tau^\#$, which is finitely generated since $R$ is Noetherian.
This supplies the lattice clause of \path{dualFr_lattice_and_hdual}
(\path{DualLatticeIsLatticeAndHermitianDual.lean}:92--147,173--175).
Choose $B_\tau$ with columns an $R_{\bar w}$-basis of
$(M_\tau)_{\bar w}$ and put $B'_\tau=(B_\tau^*G_\tau)^{-1}$.
Since $B_\tau^*G_\tau B'_\tau=1$, the tests against the columns of
$B_\tau$ read off the coordinates $(B'_\tau)^{-1}x$. By the
localized dual test these coordinates belong to $R_w$ precisely when
$x\in (M_\tau^\#)_w$. Thus $B'_\tau$ is a basis there, including
at primes fixed by the involution. This proves the frame step in
\path{mem_colLat_dual}, \path{isBasisAt_hdual} and \path{exists_dual_frame}
(\path{TransportedWeylLatticeHermitianDuality.lean}:162--201,204--210).

Take in turn $(T,T')=(\nabla,\Delta)$ and $(\Delta,\nabla)$ in the
standard polynomial model. Write $A=T(M)$ and $A'=T'(M^\#)$ for
the transported polynomial modules. The standard $T$ is finitely
generated by the integral bases above
(\path{std_fg_span}, \path{StdWeylLatticesFgAndSpanCoeffSpace.lean}:108--121).
Every element of $A$ belongs to $V$: the preceding localization
argument gives this for $\nabla(M)$, and
$\Delta_{\rm std}\subset\nabla_{\rm std}$ since its divided
words are integral polynomials in the same coefficient space.
Transport preserves inclusions by the defining local tests
(\path{transLat_mono}, \path{delStd_le_nabStd},
\path{CoefModLemmas.lean}:248--257).
Thus restricting to $V$ loses no test vectors.

For $x\in V$ set $x_0=\mathrm{act}(B'^{-1})x\in V$.
The localization formulas just established give
$A_{\bar w}=\mathrm{act}(B)T_{\bar w}$ and
$A'_w=\mathrm{act}(B')T'_w$. The covariance identity is
\[
 P_G(\mathrm{act}(B)y,x)
   =P_{B^*GB'}(y,x_0)=P_1(y,x_0).
\]
Indeed $x=\mathrm{act}(B')x_0$; adjointness of row substitution for
the factorial monomial pairing moves the first block substitution
to the second, and the Gram determinant twist transforms by
$\overline{\det\mathrm{Tw}(B)}\det\mathrm{Tw}(B')$, precisely
the two scalar factors in the actions. This is \path{pairTransport}
and its restriction \path{pairK_apply}
(\path{CoefModLemmas.lean}:477--504).
Therefore $P_G(A_{\bar w},x)\subset R_w$ is equivalent to
$P_1(T_{\bar w},x_0)\subset R_w$, hence to $x_0\in T'_w$ by
the localized dual test applied to standard Weyl duality, and
finally to $x\in A'_w$. Intersecting over all $w$ gives
$P_G(A,x)\subset R\iff x\in A'$ in $V$. The two choices of
$(T,T')$ prove both equalities. See \path{dual_of_local} and
\path{GlobalTheory.TransDual.nabDelDual}
(\path{TransportedWeylLatticeHermitianDuality.lean}:218--274,278--309).

For the Z4.4 numerator, \path{GlobalTheory.PeriodClassSum.pairBDualK}
sets the star on $K$ to the identity and $G_\tau=1$. The vector
dual then becomes the ordinary dot-product dual. Symmetry of the
bilinear normalized pairing identifies $P_1(y,x)$ with
$\mathrm{pairBK}(x,y)$; the two transported equalities give
the two integral tests. \path{PeriodClassSum.pairBDual} applies
this to the scalar extensions of the class lattice
(\path{IsClassSumPairingFormsLat.lean}:181--197).
\path{PairAtClass.pair_integral} uses the first test on the integral
class component and the supplied opposite row to get its coefficient
in $R$ (\path{ClassPairingIntegralFormRowIsIntegral.lean}:49--57).
This is the route from the registered standard Weyl and transported
localization theorems into the Z4.4 root. The public alias
\path{GlobalTheory.T.transported_weylLattice_hermitianDuality}
(\path{TransportedWeylLatticeHermitianDualityStatement.lean}:14--15)
states the same two equalities.

% blueprint/ch_P_r12.tex lines 586 through 795; r12:point
\subsection{One adelic point realizing the lattice, decorations and frames}\label{r12:point}

\paragraph{The genus and the residual decorations in the common point lift.}
Let $E/F$ have the CM datum $hc$, $n>0$, $L=\mathcal O_E^n$ and
$h(x,y)=\sum_i\bar x_i y_i$. For a full self-dual lattice $N$
we construct $x_0\in U(h)(\mathbb A_{F,f})$ whose associated global
lattice is $N$. The Lean proof of \path{PtExist.unimodGenus} applies
\path{Mass.Genus.exists_globLat_eq} to \path{SelfDualPoint.mem_genus}.
The latter finds a rational unitary using isometry over semilocal rings
inside $E$, by \path{Mass.UnimodLoc.exists_isom_of_isUnimodularAt}.
The hand proof below uses a different construction: it classifies the
completed lattices, approximates a local unitary by a rational one, and
then glues. See \path{ValuedLtOfMem.lean}:102--123,
\path{ExistsFiniteAdelicPointLatticeEqSelfDual.lean}:14--27 and
\path{ExistsIsometryOfUnimodularLatticesAtPlace.lean}:229--245.
The form is Hermitian, testing against each coordinate vector proves
nondegeneracy and gives $L^\vee=L$. Conjugation has order two with
fixed field $F$ by $hc.\mathrm{fixed}$. At each maximal $P$ the datum
$hc.\mathrm{cmUnram}$ supplies an integral $\theta$ with
$\theta-\bar\theta\notin P$ (\path{StatementMass.lean}:44--48).
At a ramified prime in a quadratic extension conjugation acts trivially
on the residue field, contradicting this inequality. Thus at every
finite $v$ we have a split pair or an unramified quadratic extension.
Duality commutes with completion over $v$: in a local basis the
dual is obtained by applying the inverse Gram matrix to the dual
coordinate lattice, including the interchange of factors at a split
place. More explicitly, self-duality gives a bijection from the
conjugate lattice to its linear dual. A finite projective lattice is
a direct summand of a finite free module, so its dual after completion
is obtained by completing the dual of this summand; on a free module
this is just taking a dual basis. Completing the bijection therefore
makes the new pairing perfect. This is the completed version of the
test in \path{mem_dualAt}
(\path{ExistsIsometryOfUnimodularLatticesAtPlace.lean}:155--184).

At a split pair, choose a basis on the first factor. The perfect pairing
with the second factor makes its dual basis a basis there and
identifies the paired lattice with the standard one. In the unramified
field case the residue Hermitian form is nondegenerate. In positive
rank it represents a nonzero element: if every diagonal vanishes,
evaluating the diagonal on $u+a v$ gives
$\mathrm{Tr}(a h(u,v))=0$ for all residue scalars $a$, and the
nondegenerate trace pairing forces $h=0$. Lift such a vector; its
length $d$ is a unit in the base ring. Every such unit is a norm.
Indeed the norm of the finite quadratic residue field is onto by
cyclicity. Starting with a lift of a solution modulo the uniformizer
$\pi$ of $F_v$, multiplication by $1+\pi^j t$ changes the norm by
the factor $1+\pi^j\mathrm{Tr}(t)$ modulo $\pi^{j+1}$ for $j\ge1$.
The residue trace is onto also in characteristic two, and the corrections
converge. Scale by a unit of norm $d^{-1}$. Subtracting
$e h(e,z)$ splits off the resulting norm-one vector $e$ from $z$;
its orthogonal complement is free over the DVR and self-dual by the
orthogonal splitting. Induction supplies an orthonormal basis and a
local unitary carrying $L_v$ to $N_v$.

The rational semilocal formulation in \path{SelfDualPoint.mem_genus}
uses a rational isometry for each individual $v$. To see this stronger
point, take a local unitary $u$ just constructed. Multiply by a rational
norm-one scalar $\lambda$ for which $1+\lambda u$ is invertible:
there are finitely many excluded scalars and infinitely many choices
$(t+\delta)/(t+\bar\delta)$ for $t\in F$ and
$\delta\ne\bar\delta$. The Cayley transform
$(1-\lambda u)(1+\lambda u)^{-1}$ belongs to the skew-adjoint
$F_v$-vector space. Approximate it by rational skew-adjoint matrices,
using density of $F$ in $F_v$ in an $F$-basis of that space. The inverse
Cayley transform is continuous here and gives rational unitaries;
multiplying back by $\lambda^{-1}$ approximates $u$. Lattice stabilizers
are open, so a sufficiently close approximation has the same local
lattice image. Intersect the completed lattices with $E^n$ to recover
the localizations at the prime or conjugate pair.
See \path{ExistsIsometryOfUnimodularLatticesAtPlace.lean}:186--245 and
\path{ExistsFiniteAdelicPointLatticeEqSelfDual.lean}:14--27.
Clearing denominators for generators in the two directions gives
nonzero integers $a,b$ with $aL\subset N$ and $bN\subset L$. Choose
the local isometries over divisors of $ab$, and the identity elsewhere,
where both lattices coincide. The resulting adelic unitary $x_0$
has global lattice $N$. Indeed if $z$ misses one of two full lattices,
its denominator ideal $\{a\in\mathcal O_E:az\text{ belongs to that lattice}\}$
is proper; a maximal ideal containing it witnesses failure of membership
in the localization. Testing both inclusions at the finite primes
therefore suffices (\path{Mass.Genus.exists_globLat_eq},
\path{UnitaryMassLatticeEqGenusMass.lean}:154--212).

Three declarations under \path{GlobalTheory.Published} are recorded as
literature leaves in this cone. Their bodies give the following
elementary proofs. For a CM field $E$ with conjugation $c$, expanding
$(P^{\mathsf t}G\,c(P))_{ab}$ gives
$\sum_{i,j}P_{ia}G_{ij}c(P_{jb})$, the pairing of columns $a,b$
in the convention of \path{herm_col}. For \path{exists_cc_eq_neg},
nontriviality of $c$ gives $x$ with $c(x)\ne x$; then
$\theta=x-c(x)\ne0$ and $c(\theta)=-\theta$. For any such $\theta$,
a relation $s+t\theta=0$ with $s,t$ in the maximal real subfield
also gives $s-t\theta=0$ by conjugation. Characteristic zero gives
$s=0$, and $\theta\ne0$ gives $t=0$. Since $E$ has degree two over
that subfield, $(1,\theta)$ is the required ordered basis in
\path{exists_basis_one_imaginary}. The three graph leaves are thus
accounted for within the hand argument (\path{AlgebraMapNormOfBasisOneSqrt.lean}:33--36,252--257,259--277).

\paragraph{Decorated-point surjectivity on the standard lattice.}
The implication \path{Formalism.L2.Col.decoPtSurj_of} takes
\path{UnimodGenus}, \path{RedSurj}, splitness of $S_T\cup S_G$ and
disjointness of $S_T,S_G$, at $n>0$. The first hypothesis supplies
$x_0$ for the given $N$. The second will be used on the single standard
Gram matrix and standard frame.
For disjoint $S_T,S_G$ with $S_T\cup S_G$ split in the precise sense
that no two members, including a repeated member, are conjugate,
set $r_P(z)=\mathrm{red}_P(x_0^{-1}z)$. It induces an isomorphism
$N/PN\simeq k(P)^n$: localization and completion at $P$ preserve
the quotient modulo $P$, and $x_0^{-1}$ gives a basis identification
of the completed lattices. This proves surjectivity and identifies
the kernel on $N$ (\path{redN_surj}, \path{reduction_kernel_equality},
\path{AdelicLemmas.lean}:538,1032). For $P\in S_G$,
the prescribed frame map is additive, residue-linear, onto, with kernel
$PN$ on $N$, so its composition with $r_P^{-1}$ on the quotient gives
$Y_P\in\mathrm{GL}_n(k(P))$. Explicitly the reverse factorization
through the prescribed frame exists by the same kernel equality.
Composing the two factorizations fixes the images of every vector of
$N$ under the respective surjections, so the two matrices are inverses.
For $P\in S_T$ the specified functional
factors in the same coordinates as a nonzero row; extend it to a basis
of the dual with that row last, giving $Y_P$ also in this case.
This recovers the two uses of \path{lin_factor} and the application
of \path{exists_GL_lastRow} (\path{AdelicLemmas.lean}:1276--1290,1293--1315).
Disjointness now defines one family on the union. In Lean, \path{RedSurj}
applied to this family at the identity Gram matrix supplies $k$
(\path{AdelicLemmas.lean}:1318--1337);
\path{PtExist.exists_isPtOf_cm} obtains that hypothesis from
\path{reduction_surjective_bis} (\path{ValuedLtOfMem.lean}:309).
Here is a direct proof of this specialization. Lift each
matrix entrywise to the local integers above $P$; its determinant is a
unit by reduction. On $\bar P$ prescribe the inverse conjugate transpose
of this lift. The split-set and disjointness hypotheses prevent conflicting
prescriptions. Use identity on the remaining factors. This gives one
integral unitary $k$ stabilizing $L$ with all the prescribed reductions,
the identity-Gram, standard-frame instance of \path{reduction_surjective_bis}
(\path{DefGroupLemmas.lean}:26--27,1475--1481).
Now $x=x_0 k^{-1}$ still carries $L$ to $N$. The reductions of
$x^{-1}z=kx_0^{-1}z$ give the frame maps or their specified last
coordinates on $N$. On $N_P$ clear a denominator $s\notin P$, apply
this equality on $N$ and cancel its nonzero residue. Off $N_P$ the
decoration axioms prescribe zero. This realizes all decorations by one
point and proves the implication \path{decoPtSurj_of}
(\path{AdelicLemmas.lean}:1245--1401) in the application at
\path{ValuedLtOfMem.lean}:309.

\paragraph{Prescribing the frames of an existing adelic representative.}
This is the point-lifting step used by several of the column constructions.
Fix a unitary finite adele $x$ for the standard Hermitian space and
$N=\mathrm{globLat}(x,\mathcal O_E^n)$. Let $Q$ be a finite set of primes
disjoint from its conjugate and suppose $B_P$ is the matrix of the prescribed
rational $E$-linear frame for $N$ at $P\in Q$, evaluated in the completion. We construct a unitary $g$ carrying the
standard lattice to the same $N$, with these precise matrix entries at
$Q$, whose pullbacks off $Q\cup\bar Q$ agree with those for $x$.
This is \path{GlobalTheory.Mem.Col2.PtExist.exists_modify}
(\path{ValuedLtOfMem.lean}:219--295).

Replace $x_P$ by $B_P$ and $x_{\bar P}$ by the conjugate of
$(B_P^{-1})^{\mathsf t}$ and use $x$ at every other place. The component
idempotents give an adele entry by entry: the two finite sums supported on
$Q$ and $\bar Q$ have disjoint support. The proposed inverse uses
$B_P^{-1}$ at $P$, $\overline{B_P}^{\mathsf t}$ at $\bar P$ and
$x^{-1}$ elsewhere. At all places the products are the identity, and the
involution transposes these inverse entries. Thus $g$ is unitary. At $P$,
$x_P^{-1}B_P$ and its inverse are integral: both matrices express one
basis of the same completed lattice in the other. Here the localized
global lattice carried by $x_P$ is the completion of $N$:
\path{carries_globLat} obtains this by integral adelic approximation.
Explicitly, clear a common denominator to sandwich the adelic lattice
between $d\widehat{\mathcal O}_E^n$ and
$d^{-1}\widehat{\mathcal O}_E^n$. In their finite quotient, the Chinese
remainder map on global vectors is onto. Thus the global intersection
surjects onto the same subquotient and spans the original lattice
after completion. At $\bar P$ the change of basis is the conjugate
inverse transpose of the one at $P$, hence integral in both directions.
At every remaining place it is the identity. Therefore $x^{-1}g$ and its
inverse preserve the standard integral adelic lattice, giving the
same transported lattice, and off those paired places even the full
pullbacks are identical. See \path{adAt_glue}, \path{inv_apply},
\path{int_of_uni}, \path{valued_pullV_le}, \path{valued_symm_le},
\path{carries_globLat} (\path{ValuedLtOfMem.lean}:92--100,133--150,172--178,180--210).

If $x$ already realizes the prescribed residue decorations at $S_T,S_G$,
disjointness from $Q\cup\bar Q$ now preserves their actual pullbacks:
one last coordinate at $S_T$, every coordinate at $S_G$. Rank zero
uses the identity and empty coordinate tests. In positive rank the
earlier existence of $x$ requires unimodular genus realization and
surjectivity for the specified reductions; changing frames by itself
does not prove that existence. These are the two parts of
\path{PtExist.exists_isPtOf_cm} (\path{ValuedLtOfMem.lean}:299--337)
that have to be checked separately.

For the application over all twists let $S_T$ be the primes dividing
$\pi_u$, $S_G$ the generator primes, and $Q=I.SQ$. At $P\mid\pi_u$,
$T=\pi_u\bar\pi_u$ belongs to $P$. No generator prime lies over $T$,
and the rational primes under $Q$ avoid the bad set containing $T$ and
the generator primes. Thus the three sets are disjoint; this property
persists for the mixed pairs after conjugation since the rational prime
under a place is unchanged. No $P,\bar P$ both divide $\pi_u$,
as $(\pi_u,\bar\pi_u)=1$. The chosen split orientations at generator
places exclude conjugate pairs there. For $Q$ the selected prime above
a given prime of the real field is unique, and the selected place is
split. Consequently $S_T\cup S_G\cup Q$ satisfies the conjugate-exclusion
hypothesis. Apply the decorated-point construction to obtain one matrix
for the entire decorated lattice, with exactly the marked reductions and
specified frames. The genus, residue and frame calculations above thus give
\path{GlobalTheory.Mem.Col2.PtExist.exists_isPtOf_cm} and its application
\path{GlobalTheory.T.exists_adelicPoint_decoratedLattice_allTwists} on their stated hypotheses. See \path{good_of_mem}, \path{plc_mem}, \path{levOk_twists}
(\path{LevOkTwists.lean}:24--29,31--46,48--85), \path{exists_ptOf_twists}
(\path{ExistsAdelicPointDecoratedLatticeAllTwists.lean}:11--60), and
\path{ExistsAdelicPointDecoratedLatticeAllTwistsStatement.lean}:7--8.

% blueprint/ch_P_r12.tex lines 797 through 1038; r12:s4
\subsection{Supplier 4. A single Hermitian datum at the level prime}
\label{r12:s4}
Fix $I$, its witnesses \path{Ok} and \path{Pins}, laws $l$, later data
$Z$, Hermitian columns $H$, and forms $F_u$ on the chart family of $H$.
Assume $\mathrm{StdF}(F_u)\land\mathrm{StdLvC}(F_u)$ at every word
and $\mathrm{CarH}(H,F)$. We construct $A$ with $\mathrm{CarT}(H,F,A)$,
symmetric adelic compatibility and symmetric stabilizer-map compatibility.
All three conclusions concern the same $A$. The statement is
\path{GlobalTheory.T.exists_hermitian_data_carT_of_std}
(\path{HermitianDataAtLevelPrimeStatement.lean}:8--9); its eight calls
and their outputs are displayed in \path{TC.hermitian_data_at_level_prime}
(\path{HermitianDataAtLevelPrime.lean}:18--62). Bare filenames for these calls
are in \path{Development/AbelianVarietyTheory/Proofs/IwasawaDecomposition/}.

\paragraph{Choosing the prime and computing the discriminant quotient.}
Fix $u$, put $\pi=F_u.\mathrm{atT}.\pi_U$, and let $P$ be the kernel of
the surjective residue map at $v$. It kills $\pi$ and is nonzero on $\bar\pi$.
Pull the valuation at a test place back through its given embedding.
Its positive elements on $\mathcal O_F$ constitute the specified prime
below $P$. The field \path{AtTTied.seenP} supplies $\nu(\pi)>0$ in the true branch and $\nu(\bar\pi)>0$ in the false branch. In the true branch the fixed integer $\pi+\bar\pi$ has
nonzero reduction, so $\nu(\bar\pi)=0$. For $y\in P$ the fixed
integer $y\bar\pi+\bar y\pi$ has zero reduction. Its valuation and
that of $\bar y\pi$ are positive. Their difference is $y\bar\pi$,
so $\nu(y)>0$. The positive elements form a proper ideal on
$\mathcal O_E$ by nonnegativity of the valuation there; it contains the
maximal ideal $P$ and hence equals it. In the other branch apply this
argument to $\nu\circ\mathrm{conj}$ and conjugate the ideal back.
This gives precisely the two alternatives of \path{SeenCl}.
See \path{pos_iff_mem_Pv}, \path{seenCl_of},
\path{HermitianSeenPrimeConditionHermitianSide.lean}:15--105.

In the standard coordinates of supplier 6 the class lattice and its dual
are the decorated $L,L^\sharp$. For clarity, the local calculation behind
the discriminant identification is as follows. In the self-dual lattice
$N$ the kernel defining $L$ imposes one surjective residue functional
at each selected $P$. Locally at $P$ lift an element on which that
functional is $1$ and complete it to a basis after choosing the other
residue vectors in its kernel. Thus the elementary divisors of
$L_P\subset N_P$ are $1,\ldots,1,\varpi_P$.
All selected primes are distinct and none is conjugate to a selected
prime. At other places $L=N$ locally, since denominators away from the
place can kill all those finite residue conditions. The perfect
Hermitian pairing exchanges the localization at $Q$ with that at
$\bar Q$: in dual bases one inverts the elementary divisors. Consequently
$L^\sharp/L$ has a single residue line at each $P$ and $\bar P$
and vanishes elsewhere. Its support is finite, so CRT on its annihilator
splits it into the direct sum of precisely those lines. The selected
prime product is $(\pi)$: the datum gives square-free factors, or
equivalently divisibility by $\pi$ for every element lying at all
selected primes. Its conjugate is $(\bar\pi)$ and the two ideals are
coprime, with product $(T)$. CRT therefore gives a linear isomorphism
$L^\sharp/L\simeq\mathcal O_E/(T)$. Only an existence assertion is
needed; no residue generator is fixed. Pulling the quotient back along
the common coordinate map gives \path{DiscCl}. See \path{Deco.Ok.Lsh_quot_L},
\path{Formalism/LocalToGlobal/Columns/DecoFin.lean}:2199--2394 and
\path{quotMapEquiv}, \path{span_pi_eq}, \path{quotT_equiv}, \path{TC.disc_cl},
\path{DualModClassLatticeEquivQuotient.lean}:18--32,40--46,71--92,98--110.

\paragraph{Middle projections.}
Both the chart and one-level projections evaluate on the same underlying
vector. To compare their different lattice presentations, multiply that
vector by $s=\bar\pi$, whose residue at $P$ is a unit. In coordinates
$L^\sharp$ is the decorated dual and the kernel lattice is
$L+P L^\sharp$. For the latter identity, choose by CRT a multiplier
$b$ with $sb=1$ modulo $P$ and $b=0$ at all other selected primes.
For a vector $x$ in the decorated kernel, $sbx$ is in $N$ and has all
its residue functionals zero. Hence $sbx\in L$ and
$x=sbx+(1-sb)x\in L+P L^\sharp$; the reverse containment follows
by applying the residue functional. Multiplication by $s$ places a
dual vector in $N$, and a kernel vector in the lower chart lattice:
it clears the possible simple poles at the conjugate primes and
preserves the zero residue condition at $P$.
On these smaller lattices the two standard reading laws
$\mathrm{StdF}$ and $\mathrm{StdLvC}$ identify both projections with
the same decorated value. Both maps scale by $\operatorname{res}_v(s)$,
so cancellation gives equality before multiplying, on both domains.
See \path{pr_lvC_eq_lv},
\path{MiddleStepOneLevelProjectionEqChartProjection.lean}:12--67;
\path{lv_LshKer_eq}, \path{lv_smul_LshKer_mem},
\path{Formalism/LocalToGlobal/Columns/TiedLv.lean}:97--152.

\paragraph{Transport shared by the three moves.}
The shared transport and the up and down calculations below are
outlines of the Lean sublemmas \path{sameoff_gen},
\path{upCl_of_crd}, \path{dnCl_of_tie} and \path{dn_move_exists}. They specify
the lattice and functional identities used in the supplier assembly;
the local transport sublemmas carry their Lean proofs in the cited
ranges.
Write $c'$ for the new class and $y,y'$ for the representatives on
one and the same chart. A coordinate move is a rational unitary
$\gamma$ with
\[
   \gamma y'=y\eta\kappa,\qquad \kappa\in U,\quad \eta_w=1
                      \quad(w\ne v_0),\qquad v_0=e_V(v).
\]
Here the left factor is extended to the adeles. At $w\ne v_0$ the
identity gives $y_w^{-1}\gamma y'_w=\kappa_w$. The lattice reading
of the representatives and preservation of the base lattice by
$\kappa_w$ now give $(\gamma L_{c'})_Q=(L_c)_Q$ at every $Q\mid w$.
Indeed completing a finite lattice and then projecting to a factor above
$w$ preserves its generators; an inclusion between two rational
lattices can be tested there by clearing a common denominator and
testing the elementary divisors over the DVR at $Q$. Completion preserves
those valuations. The involution in a quadratic Galois extension
permutes the primes over $v_0$ transitively, so the only excluded
$Q$ are $P,\bar P$. For the functionals at another $v'$, if
$z=\gamma z'$, the local input for $z$ is $\kappa_{e_V(v')}$
times the one for $z'$. Both belong to the local dual lattice, because
unitarity transports the defining integral pairing. The level condition
makes $\kappa_{e_V(v')}$ fix the local functional. Applying the chart's
additive identification of value groups gives equality in the original
residue field. In an interior chart this final identification uses the
middle projection comparison just proved. Finally $\mathrm{CarH}$
transports the displayed adelic equation: reindex $\eta,\kappa$ by
the common rank, replace representatives by the identified ones, and
replace $\gamma$ by its conjugate under the Hermitian frame. The result
is precisely $\mathrm{AdOf}$; use chart zero, an interior chart, or
the unique top chart according to $i$.
See \path{eq_Pv_or_conj}, \path{HostPrimesOverLevelPlace.lean}:25--36;
\path{locLat_of_adStep}, \path{sameoff_gen},
\path{AdelicConjugateCharacterStandardForms.lean}:133--158,213--235;
\path{fT_move}, \path{fv_eq_of_adC},
\path{LevelFunctionalEqOfClassesRelatedAtPlace.lean}:16--58,91--138;
\path{adOf_of_adC_gen}, \path{AdelicRelationClassesCoordinates.lean}:51--109.

\paragraph{Up and down.}
For $i=0,m$ take the $\gamma,\eta,\kappa$ from the endpoint fields
$F_u.\mathrm{upN}$ or $F_u.\mathrm{dnN}$. At $0<i<m$ take them from
$F_u.\mathrm{upC}$ or $F_u.\mathrm{dnC}$ at the label of $c$.
These interior fields also certify that $c'$ has that same label. This
is where the assumption that the forms on these fixed columns are given
enters. In the up case the input is $y_0\in L_c^\sharp$ of projected
functional value $1$. At $P$ the given two lattice equalities read
\[
 (\gamma L_{c'})_P=(L_c^\sharp)_P,\qquad
 (\gamma L_{c'}^\sharp)_P=
      (L_c^\sharp+\mathcal O_E T^{-1} y_0)_P .
\]
For $x\in L_{c'}^\sharp$ with $T\gamma x\in L_c^\sharp$ the
selected functional on $x$ equals that on $T\gamma x$. In the down
case the input is a nonzero residue linear functional $g$ on the lower
space. Here $(\gamma L_{c'}^\sharp)_P=(L_c)_P$; for
$\gamma x\in L_{k,c}$, the selected functional on the projection
of $x$ equals $g$ on the lower projection of $\gamma x$. These are
the selected-place fields of the respective coordinate moves. Replace
chart projections on both arguments by the common one-level projections
at interior stages. Transporting back by the frame preserves the dual
(pair against all lattice vectors), sums, scalar multiplication and
localization (clear the same denominator). The other-prime equalities
and functional identities follow from the preceding paragraph. Thus
the very same $\gamma$ gives both the move law and its adelic reading.
For a false guard ($i>m$, incorrect up value, or zero down functional)
choose $1$; the conclusion for this index is conditional on the guard.
See \path{upCl_of_crd}, \path{crd_up}, \path{up_move},
\path{HermitianUpMoveAtLevelPrimePlace.lean}:19--72,82--116,125--143;
\path{dnCl_of_tie}, \path{dn_move_exists},
\path{ExistsHermitianDownMoveAtLevelPrimePlace.lean}:15--86.

\paragraph{Diamond.}
For a residue unit $a$ choose an integer $k$ mapping to it; the field
has prime order $T$, so $T\nmid k$. The choice works for every
relevant chart. At $v_0$, $\alpha=\pi-\bar\pi$ is anti-invariant,
$\alpha^2\in\mathcal O_F$. Lift the nonzero residue of $\alpha$ at the selected prime to an integer $b$ of $F$ using the given residue identification. Squaring gives $b^2\equiv\alpha^2\pmod {v_0}$. Since
\path{D.T_ne_two u} follows from $D.\mathrm{T\_mod\_four}(u)$,
$2b$ is a unit (\path{T_ne_two}, \path{IntCastDvdOfDvdIntegers.lean}:118--120;
\path{splitAtT_of_tied}, \path{Formalism/LocalToGlobal/GlobalTheory/GL6/PSplitT.lean}:263--271). Newton iteration for $X^2-\alpha^2$
starting at $b$ doubles the valuation of the error at each step and
preserves the unit derivative. Completeness gives a unit square root
$d$ of $\alpha^2$ in the completed base ring. Now
$e=(1+\alpha/d)/2$ belongs to the tensor order,
$e^2=e$ and $\bar e=1-e$. On the order-$T$ additive value module an
idempotent has image either zero or the entire group. Replace $e$ by
$1-e$ if necessary so that it acts as identity. The scalar
$\mu=k^{-1}e+k(1-e)$ and its conjugate inverse are integral,
$\mu\bar\mu=1$, and $k\mu$ acts as the identity on the value module.
Take $\delta_{v_0}=\mu I$ and $\delta_w=1$ elsewhere. This finite
support change is an adelic unitary, preserves the base lattice at all
places, and satisfies the defining diamond multiplier on the chosen
functional. The endpoint or chart diamond law in $F_u$ gives
$[y']=[y\delta]$, where $c'=\mathrm{dil}(i,v,a,c)$. Lift this
double-coset equality to $\gamma y'=y\delta\kappa$. In an interior
chart its law supplies equality of chart labels too. Every local
lattice is carried to the old lattice since $\delta\kappa$ stabilizes
the base; equality at all localizations gives equality of global lattices.
At $v_0$ the local functional fixes $\kappa$ and satisfies
$f(z')=k f(\gamma z')$ after putting both inputs in representative
coordinates. The additive residue equivalence commutes with integer
multiplication, giving $a$ on the right. Use the common transport
above for all other places. Again $i>m$ may use the identity.
See \path{splitAtT_of_tied}, \path{Formalism/LocalToGlobal/GlobalTheory/GL6/PSplitT.lean}:134--201,217--271;
\path{local_diamond_of_split_idempotent},
\path{AdmissibleOfSplitAndRoundTrips.lean}:35--46,90--149,151--178;
\path{exists_int_isDia}, \path{ExistsDiamondElementsStepLevels.lean}:10--25;
\path{move_end}, \path{move_mid}, \path{diamond_on_hermitian_side},
\path{DiamondOperatorHermitianSideLevelPrime.lean}:24--120,131--157,162--192,219--249.

\paragraph{Hyperplane classes, stabilizers, and assembly.}
In the primed orthonormal frame pull each supplied class lattice
back by the inverse coordinate map. Its base lattice is exactly the
standard lattice: the coordinate test of $\mathrm{StdF}$ requires
integral coordinates and divisibility by $1$. Take the same level
subgroups, representatives and stabilizer homomorphisms; their range
and double-coset conditions survive literal coordinate transport. At
a generator prime a principal congruence in the base lattice, applied
to each coordinate vector, gives columns of $X=1+g_iW$ with
integral split coordinates for $W$. The standard generator frame
is $1$, hence reduction gives $1$ and the specified generator subgroup
contains it. Away from the generator primes the supplied hyperspecial
level law applies since the primed level parameter is $1$. Mapping
each pulled-back lattice forward gives the original lattice. This constructs
the whole primed class and level tie.
See \path{mem_Uloc_of_princ}, \path{tie_of_forms}, \path{clsP_hermitian},
\path{HermitianSideClassesEqFormClasses.lean}:25--120,143--235.

Build $A_u$ with these primed objects, the given place and field
identifications, the seen-prime and quotient laws, and the three families
of moves. In \path{CarT} the two uniformizer/place identifications
are identities, the two primed comparisons were just proved, the two
global projections are their coordinate restrictions, and the middle
one is the cancellation proved above. For the final vector and move
matrix the identities are inverse frame evaluation and matrix reindexing.
The six adelic conditions are the adelic readings from the up, down and
diamond witnesses, respectively at zero and at a positive stage.
For the stabilizer comparison write $M$ for the move to the chosen
top representative. Comparing the two values on the characteristic
function of the identity frame forces the matrix of the tilted
isometry to equal $M$: a different frame gives zero on exactly one
side, and the character multipliers are nonzero. More explicitly, equality at the identity frame gives
$M^{-1}g y=y$ at any embedding, where $g$ is the tilted isometry.
Right cancellation of the invertible frame and injectivity of the
embedding give $M^{-1}g=1$. Write $e$ for the matrix of the extended
primed stabilizer. The class stabilizer acts at the representative by
$M^{-1}eM$; transporting by $g=M$ gives $e$. By the extension laws,
$e$ acts on $(x,0)$ by the prescribed primed stabilizer and fixes
$z_0$. The carrier comparisons identify the coordinate actions with the
actual two actions, and injectivity of the frame then gives both
identities (\path{eq_of_frMul}, \path{crdN_stabM}, \path{crdP_stabP}).
Sources: \path{ga_eq_gam_mv}, \path{gam_ext}, \path{stabMap_tied},
\path{StabiliserMapHyperplaneEqConjClassMap.lean}:27--40,50--156;
\path{TC.hermitian_data_at_level_prime}, \path{HermitianDataAtLevelPrime.lean}:18--62.
The supplier-4 implication, with $F_u$ and $\mathrm{CarH}$ as hypotheses,
is confirmed under these hypotheses.

% blueprint/ch_P_r12.tex lines 1040 through 1579; r12:s9, r13:s9-values
\subsection{Supplier 9. Points and evaluations over the fixed level}
\label{r12:s9}
Write $D=I.\mathrm{rec7}(l,Z)$ and fix $H,F,T,J,W,L_v$ as in
\path{GlobalTheory.T.tupleLayerCompatibilities_of_X_prime}
(\path{TupleLayerCompatibilitiesOfXPrimeStatement.lean}:9--16).
Here $L_v$ denotes the whole record \path{Lv}, including its comparisons at
all places. Retain $I.\mathrm{Ok}$, the pins, $l:I.\mathrm{Laws}$,
the additional law $l_0:I.\mathrm{Law0}(l)$, the
equivalence defining $\mathrm{UnderQ}$, the standardness assumptions
on $F,W,L_v$, and the supplied top reading and top-point property for
one $k$ for each $(u,a)$ with $a\ne\mu_0$ whenever $1\le m$.
At \path{colsTwPTiedS_std} supplier 7 provides this reading and its
block-extension property, while supplier 8 provides the given $L_v$
and its standardness (\path{TrueColumnChoiceExistsStandardChainData.lean}:53--55).
The conclusion asks for a
point record $P$ and an evaluation record over this same $P$ satisfying
both the generator-place conditions and $I.\mathrm{StdP}$. The full
source cone is 34,625 core lines.

\paragraph{The tilted datum on a supplied primed chart.}
Fix $a\ne\mu_0$, $1\le m$, a primed chart $K$ with its
$\mathrm{ReadsP}$, and the actual $k$ and its two certificates at
$(u,a)$. We construct $\mathrm{Tl}$ and its evaluation equality
(\path{GlobalTheory.Mem.Col2.TwP.til_tw},
\path{ExtMMulVecFix.lean}:253--377).
In the coordinates of confirmed supplier 6 the inclusion is $x\mapsto
(x,0)$, the Gram matrix is $1$, and $z_0=\pi e_{n+1}$ with
$\pi\bar\pi=T\ne0$. Hence $z_0$ spans the orthogonal complement
of the inclusion and has norm squared $T$. An extension $g$ of $g'$
agreeing on the inclusion and fixing $z_0$ is uniquely
$\operatorname{diag}(g',1)$: the first $n$ columns are prescribed,
and the last is determined by dividing by the nonzero field scalar
$\pi$. Its image stays invertible over the adele algebra. This block
extension itself satisfies both equations. See
\path{lastCoordinateUniformizerVector_spans_orthogonal_line},
\path{hermitianForm_lastCoordinateUniformizerVector_self},
\path{eq_extM_of_ext}, \path{HermitianFormLastCoordinateUniformizerVectorSelf.lean}:692--701,766--781,953--999.

For a primed class $c$, put $d=\tilde c$ in the top system. Let $m_c$
be the specified transporter from its selected representative to the
tilted point, $\gamma$ its rational unitary component, $y$ the top
adelic representative, and $y'=K.y(c)$. For $b=\operatorname{diag}(y',1)$,
the top-point hypothesis says $bk_\beta$ represents the tilted point,
where $\beta=\mathrm{chOf}(d)$. Write $\gamma_f$ for the image at the
finite adeles. Set $x=y^{-1}\gamma_f^{-1}b$. Then $yxk_\beta$
represents the point acted on by $(1,m_{c,G},m_{c,Q})$ from the
original representative: this follows by cancelling $\gamma^{-1}$
against the rational component of $m_c$. Rational change of frame
changes the lattice and its marked residues by the same transformation
(\path{TilX.isPtOf_rat}, \path{HermitianFormLastCoordinateUniformizerVectorSelf.lean}:123--152).
The last-stage auxiliary coset is unchanged because its decoration
there is empty. The level and character clauses of the given top
reading therefore put $x\in U_\beta$ and give its character as the
embedded character of $m_c$. Taking $\kappa_c=x^{-1}$ gives both
$\gamma_f y=b\kappa_c$ and
$\chi_\beta(\kappa_c)^{-1}=\sigma(\mathrm{chr}(m_c))$.
The same witness works for every test function. This proves
\path{TwP.exists_ka_tw} (\path{ExtMMulVecFix.lean}:80--144).

Transport the top lattice by $\gamma$. Its decorated-point law gives
the tilted lattice. A vector in the latter has first $n$ coordinates
in the primed representative lattice $N_c$ and integral last
coordinate $r$, with $\bar\pi r=0$ modulo every designated prime over
$T$. At each such prime $\bar\pi$ is a unit. The given divisibility
property of $\pi$ says exactly that membership of $r$ in all these primes
implies $\pi\mid r$. Conversely a multiple of $\pi$ passes every test.
Thus the transported lattice is $(N_c,0)+\mathcal O_E z_0$.
See \path{Deco.Ok.til_L}, \path{Formalism/LocalToGlobal/Columns/DecoFin.lean}:1800--1845,
\path{TwP.til_lat_eq_tw}, \path{ExtMMulVecFix.lean}:146--165.

For the functional clause fix a designated prime $v$ above $T$ and
$x'$ in the dual of the selected top lattice. Put $s=\bar\pi$.
The Gram matrix is $1$, so the dual condition puts $x'$ in the
decorated dual $L^{\sharp}$. The lattice's defining divisibility
property puts $sx'$ in the representative lattice $N$. On $N$ the
standard level projection is the quotient map. Its nonzero induced
functional agrees there with the original functional (it factors
through that quotient by its definition). Moving the representative
by $m_c$ allows application of the given tilted functional law on $N$.
Writing $r=h(z_0,\gamma x')\in\mathcal O_E$ as in the hypothesis,
one obtains
\[
 f_v(\mathrm{pr}_v(sx'))=\mathrm{red}_v(sr).
\]
Indeed $r=\bar\pi(\gamma x')_{n+1}$, and $\gamma$ is linear.
The projection is linear for the action through $\mathrm{red}_v$;
$f_v$ is linear on the residue space. Since
$\mathrm{red}_v(s)\ne0$, cancellation gives the required
$f_v(\mathrm{pr}_v x')=\mathrm{red}_v(r)$. This is where the
choice of the conjugate uniformizer matters.
\path{TwP.til_fv_eq_tw}, \path{ExtMMulVecFix.lean}:168--249;
\path{dualL_le_Lsh}, \path{TDat.fv_mk},
\path{HermitianFormLastCoordinateUniformizerVectorSelf.lean}:809--831.

The stabilizer of the tilted representative is obtained by conjugating
the extended stabilizer by $m_c^{-1}$. The supplied period law on
initial and last coordinates forces the rational matrix of that
extension to be $\operatorname{diag}(\mathrm{gam}(s_c),1)$: apply the
two coordinate laws to each basis vector. Conjugation by $\gamma$
therefore gives exactly the required stabilizer comparison for $K$.
Set $\mathrm{Tl}.z_0=z_0$, $\mathrm{Tl}.ga(c)=\gamma$ and
$\mathrm{Tl}.ka(c)=\kappa_c$. The block uniqueness above yields
its adele identity for every extension, not only for $b$.
For two frames in the prescribed transport relation, standardness
of $W$ and the end-frame certificate on $F$ identify their underlying
matrices, and the relation gives $s'=\gamma_\infty^{-1}s$ in
these coordinates. Evaluating the moved function now gives
$\sigma(\mathrm{chr}(m_c))\phi(d,s')=
\chi_\beta(\kappa_c)^{-1}\phi(d,s')$. This proves the last field for
arbitrary $\phi$ without choosing new witnesses. See
\path{eq_extM_of_laws} and \path{til_tw},
\path{ExtMMulVecFix.lean}:63--76,253--377.
The candidate here is the exact implication \path{TwP.til_tw} under
its chart, level and top-point assumptions.

\paragraph{One denominator for all classes.}
A separate deduction applies to any $\mathrm{Ok}$ step, its finite
class set, its compatible coefficient family and the given valuation
realization at $w$. The integral point lattice imposes membership in
$\Delta_c=\mathrm{delAtr}(\mathrm{emb},\lambda,L^{\sharp}_c)$
at each representative. After localization at $w$, its membership
test is exactly $x_c\in(\Delta_c)_w$ for every $c$.
A common denominator $s\notin w$ proves one direction. Conversely
choose $t_c\notin w$ with $t_c x_c\in\Delta_c$. The product $s=\prod_c t_c$
is outside $w$ by primality. Multiplying $t_c x_c$ by the other
factors gives $sx_c\in\Delta_c$ for every $c$; the scalar preserves
the stabilizer equations since $x$ was already a point. The empty
product is $1$. This proves \path{GlobalTheory.Mem.Col2.mem_locW_hlat_iff}
(\path{Development/AbelianVarietyTheory/Proofs/GlobalInterpolation/MemLocWHlatIff.lean}:14--30),
with \path{vl_mem_span_intW_iff} (\path{Formalism/LocalToGlobal/Columns/TiedValBwG.lean}:20--47) for the place dictionary.

In a local basis $B$, the dual frame is the inverse transpose of its
base change. Transposing again in the evaluation formula explains
why the value is evaluated in $B^{-1}$, not in $B$. The fixed
place test on $\Delta_c$ is equivalent to the assertion that the row
value comes from a polynomial in the localized standard lattice in
that frame. This equivalence uses the transported-lattice localization
identity \path{GlobalTheory.LvT.Built.ValLat.mem_localization_delAtr}
(\path{MemClosureIff.lean}:141--150, inclusive core 5,887); its converse
uses injectivity of polynomial evaluation on frames, as in
\path{BwLat.mem_localization_delAtr_fn}
(\path{IsBasisAtOfIsLocBasis.lean}:59--78).
For the quantifier over \emph{every} local basis, the forward argument
applies to any $B$. For the reverse implication choose one at each
class, possible since the representative lattice is full, and use the
common-denominator argument. The coefficient subgroup is the image
under $\sigma$ of the localized standard lattice: generators have
integral coefficients, localized scalars have a denominator outside
$w$, and such denominators act invertibly in the valuation ring.
The determinant shift preserves the natural partition and hence the
integral polynomial generators. See \path{mem_closure_iff},
\path{MemClosureIff.lean}:78--125, and \path{locW_lat_iff},
\path{tw_lat_at}, \path{TupleIntegralFormsRowPlaceByPlace.lean}:34--65,77--105.
This separates the transported-lattice identity from either direction
of duality for the \emph{standard} lattice.

\paragraph{Geometric inputs to the primed construction.}
The primed chart used in both deductions must have base lattice and
local subgroups at the generator primes matching those of the reference
chart. Given its level, its character on the restricted level is the
reference character times the additional character at $Q_a$.
An element trivial at every place of $Q_a$ has all added local
characters equal to $1$, so the two characters agree there. If their
exponents are $m_1,m_2>0$, the product has exponent dividing $m_1m_2$.
The input \path{ReadsP} on the reference chart is used first to
identify its level and character from the degree-zero representative.
Generators lie outside $Q_a$ by the given separation hypothesis.
See \path{memberLevel_read} and \path{tuplePrimedSpace_classes},
\path{TuplePrimedSpaceClassesLevelAndComparison.lean}:84--120,131--267.

Two construction steps in that argument require the following geometric input.
For each decorated representative of the primed step, produce one
unitary finite adele giving its lattice, marked residue data and
specified frames simultaneously. The Lean construction first realizes
the unimodular lattice in the genus of the standard one, lifts the
marked reductions, and modifies the frames at $Q_a$ while preserving
the previously fixed components. This is \path{PtExist.exists_isPtOf_cm}
(\path{ValuedLtOfMem.lean}:299--337; 8,918 inclusive core lines), with
\path{unimodGenus} at lines 102--123 (4,078).
Then these representatives must identify class indices bijectively
with double cosets of the restricted level and identify stabilizers
with its rational conjugate intersections. The new level must retain
the old lattice and local groups outside $Q_a$; the level test must be
equivalent to movement by the two prescribed residue groups.
See \path{TopCh.exists_clsG_top},
\path{ExistsFramedChartFromAdelicDictionary.lean}:98--158 (7,471),
and \path{decoratedPointClasses_bijective_doubleCosets},
\path{DecoratedPointClassesBijectiveDoubleCosetsStatement.lean}:8--9 (2,412).
The resulting primed-chart supplier has 24,746 inclusive core lines.

Once one has $K_{u,a}$ with those properties, \path{forms_tw} uses it
for the complex and rational point tests, the integral frame criterion,
the orthonormal class sum and the metric norm. The determinant shift
in both metric identities agrees since $t=2k+1$ gives
$m(t+1)/2=m(k+1)$. For each nonreference pair call the tilted-datum
implication on this very chart and the prescribed top $k$. Use the
identity for the row equivalence and the conjugate period map of
the preceding round for $j_M$. Now the class, valuation, and tilt
components belong to the same point record. See \path{weightedInner_twisted},
\path{norm_tw}, \path{forms_tw}, \path{FormsXTupleChainOfLevelReading.lean}:18--91,
\path{period_conj}, \path{CevInjective.lean}:219--225, and
\path{tupleLayer_compatibility}, \path{TupleLayerCompatibilitiesOfXPrime.lean}:17--46.
The following deductions supply the inputs to this assembly.

\paragraph{The localization test used by suppliers 8 and 9.}\phantomsection\label{r13:nabdelloc}
In \path{ValLat.mem_localization_delAtr}
(\path{MemClosureIff.lean}:141--150), take a full lattice $L$, a place
$w$ and bases for the duals of all $\tau$-extensions of $L$ at $w$.
\path{tauLatLattice} gives full scalar-extended lattices;
\path{du_dualFr_lattice} with \path{basisAtEx} makes their dot-product
duals full (\path{CoefficientsModP.lean}:21--26;
\path{TauLatticeLatticeHolds.lean}:230--247).
Unfolding \path{delAtr} now identifies its localized membership on
the coefficient space with membership of the polynomial in localized
\path{delP}. Lean substitutes the $\Delta$ equality of
\path{Formalism.L2.DIV.nabDelLoc} (\path{CoefficientsModP.lean}:24),
which calls \path{nablaDeltaLocal_holds}
(\path{NablaIndependenceHolds.lean}:578--580). This is the same
\path{NabDelLoc} proposition proved in Section~\ref{r12:weyl}
for both lattices. The registered three-part localization theorem
enters the root separately through \hyperref[r13:weyl-dual]{the proof of transported Hermitian duality}.
For each embedding $\tau$ a basis at the rational residue characteristic
$p$ gives a basis at the contraction of $w$: denominators prime to $p$
are units there. Conversely the finitely many coordinates in that basis
allow multiplication by the product of the local denominators outside the
contracted prime. Applying the given semilocal basis criterion supplies
the additional integer denominator prime to $p$, so equality holds after
localizing. Under scalar extension the columns still generate the localized
lattice. A vector pairs integrally with every column precisely when the
transpose of the basis matrix carries it into the standard localized
integer module. The dual basis is consequently the inverse transpose.
Applying $\Delta$ localization in these dual bases identifies $x_c$ with
the polynomial action on a localized standard polynomial $z$. The
transpose in the frame action gives the announced $B^{-1}$.
Equality of the resulting frame values recovers the polynomial identity:
the determinant factors are invertible there. Given arbitrary matrices
$(M_\tau)_\tau$, all $M_\tau+tI$ are invertible except for finitely
many $t\in\mathbb C$ since each characteristic polynomial is monic.
The difference of the two polynomials on this line therefore vanishes
identically, also at $t=0$. Varying the matrices and using injectivity
of the coefficient embedding proves the assertion.
See \path{isBasisAt_of_isLocBasis} and \path{mem_localization_delAtr_fn}
(\path{IsBasisAtOfIsLocBasis.lean}:26--78),
\path{pb_isBasisAt_dualFr} (\path{CoefficientsModP/PairBDual.lean}:48--85),
\path{fnOf_injective} and \path{evFr_injective}
(\path{Formalism/LocalToGlobal/Columns/PolyFn.lean}:265--287).
Combining these two implications with the common-denominator and
weight-shift arguments above yields the exact conditional candidate
\path{GlobalTheory.T.tuple_integralForms_row_place_by_place}
(\path{TupleIntegralFormsRowPlaceByPlaceStatement.lean}:11--12;
19,117 core lines) for a supplied primed chart with its reading, all
the original place and avoidance assumptions, and every admissible $w$.
Its localization input is the $\Delta$ projection of
\path{Formalism.L2.DIV.nabDelLoc} in the displayed route.

\paragraph{From decorated classes to double cosets.}\phantomsection\label{r13:s9-cosets}
Write $G_f$ for the adelic unitary group and $\Gamma$ for the image of the
rational unitary group in it. Fix the source hypotheses $U\le G_f$,
$k'\in G_f$, an adelic lift $y_c$ of each chosen decoration $p_c$
in the convention $y_ck'\leadsto p_c$, and the two following properties.
When $gk'$ and $gxk'$ represent $p,p'$, respectively, one has
$x\in U$ if and only if $p'$ is obtained from $p$ by the two prescribed
residue groups $(h,\theta)$ with trivial rational component.
For every $z\in G_f$ some $u\in U$ makes $zuk'$ represent a decoration.
Assume also that the $p_c$ meet every joint orbit of the rational and
residue groups and that any two in the same joint orbit have identical
indices. The map $c\mapsto\Gamma y_c U$ is bijective under these
hypotheses.

If $y_{c'}=\gamma_f y_c u$, rational transport shows that
$y_cuk'$ represents $\gamma^{-1}p_{c'}$. The level test gives
$\gamma^{-1}p_{c'}=(1,h,\theta)p_c$, so joint-orbit separation gives
$c=c'$. Conversely fix a representative $z$ and choose $u,p$ by
the stated density property. Write $p=(\gamma,h,\theta)p_c$.
Then $\gamma_f^{-1}zuk'$ represents $(1,h,\theta)p_c$, hence
$v=y_c^{-1}\gamma_f^{-1}zu\in U$ by the reverse level test.
Since $z=\gamma_f y_c vu^{-1}$, its double coset is represented.
This proves the exact \path{GlobalTheory.T.decoratedPointClasses_bijective_doubleCosets}
(\path{DecoratedPointClassesBijectiveDoubleCosetsStatement.lean}:8--9,
\path{ClsBij.classes_bijective}, \path{DecoratedPointClassesBijectiveDoubleCosets.lean}:34--83).
The same level test applied to $x=y_c^{-1}\gamma_f y_c$ gives the
stabilizer criterion: a residue pair moves $p_c$ to $\gamma p_c$
if and only if its inverse together with $\gamma$ stabilizes $p_c$.
See \path{TopCh.stab_iff_mem_conj} (\path{StabIffMemConj.lean}:9--40).
These deductions require exactly the indicated inputs, in particular
the density property for every adelic $z$.

For the primed chart $k=t=1$ in the generic chart construction.
The restricted level is $U'=U_0\cap\bigcap_{P\mid Q_a}
\rho_P^{-1}(\mathrm{locLev}_P)$, with the intersection running over
the designated primes $P$ in the source tuple. It has the required local
membership test: if $x\in U_0$ agrees at each place below $Q_a$ with
some element of $U'$, its $\rho_P$ image at that place satisfies the
local condition for every such $P$. The new character is the restriction
of $\chi_0$ times the inverse product of the added local characters.
Its auxiliary factor is trivial on elements trivial at all those places.
The paragraph uses the generic construction
\path{TopCh.exists_clsG_top} under the stated reading, level and
orbit hypotheses. Its proof builds the smaller group from its local
tests, applies the double-coset bijection and uses the rational
stabilizer intersection. We give the point and orbit tests here;
the construction of the abstract smaller step level inside this
lemma is used as a Lean input. See \path{TopCh.exists_clsG_top}
(\path{ExistsFramedChartFromAdelicDictionary.lean}:98--158);
\path{topChi_finite_order_and_trivial}
(\path{LevelCharacterAtQFiniteOrderAndTrivial.lean}:17--37).
For its density hypothesis the source first finds $z_1$ in the
principal level with the specified congruences at each marked prime
such that $g(k t)z_1$ represents a decorated point.
The correcting element is $(kt)z_1(kt)^{-1}$: conjugation by $k^{-1}$
lies in the reference level, and the marked components belong to
the required local levels by the principal congruences. Multiplication
on the right by $kt$ gives the same representing adele.
The lifting input used here is
\path{GlobalTheory.T.exists_K1_point_congruent}
(\path{ExistsK1PointCongruentStatement.lean}:8--9),
as used by \path{topU_dense}
(\path{ExistsFramedChartFromAdelicDictionary.lean}:63--92).
\paragraph{Producing a point within an arbitrary prescribed congruence.}\phantomsection\label{r13:s9-congruence}
Fix a unitary $g$, an integer $c\ge 0$, and a marked $P\in Q$. Put
$L=\operatorname{globLat}(g)$, and let $B$ be a rational frame of
$L_P$. Such a frame comes from a basis of the lattice over the discrete
valuation ring $\mathcal O_{E,P}$.
Both $A=g_P^{-1}B$ and $A^{-1}$ are integral since the completed
frames generate the same lattice. Write $\delta_c=|q_P|_P^c$,
where $q_P=|\mathcal O_E/P|$ is the absolute norm of $P$. Approximate $A^{-1}$
entrywise by a rational matrix $M$ to error at most $\delta_{c+1}$,
using density of $E$ in $E_P$. Multiplication by the integral $A$
gives $|AM-1|_P\le\delta_{c+1}<1$ entrywise. Thus $AM$ reduces
to the identity, its determinant is a unit, and $M,M^{-1}$ are integral
at $P$. In particular $M$ is also invertible over $E$, and $B M$
is a rational frame of $L_P$ satisfying the requested bound
$|g_P^{-1}B M-1|_P\le\delta_c$. This works even for $c=0$.
Choose one such frame separately for each $P$ and keep the original
decoration of $g$. Replace $g_P$ by this frame and its conjugate
component by the conjugate inverse transpose. The avoidance assumptions
make the pairs disjoint and leave $S_T\cup S_G$ untouched. The replacement
has unitary inverse by construction and is adelic because its support
is finite. At the marked primes the frame change and its inverse are
integral. At their partners this follows by conjugate transpose.
Consequently $z=g^{-1}g'$ preserves the standard lattice, and its
components at $S_T,S_G$ are identities. In particular it satisfies
the mirabolic condition and has trivial prescribed reductions, i.e.
$z\in K_1$. Now $g z=g'$ has the original lattice and residues and
the chosen frames exactly. This gives the full
\path{GlobalTheory.T.exists_K1_point_congruent} (5,105 core lines).
See \path{MemPrinCongOfSm.lean}:144--222,
\path{ValuedLtOfMem.lean}:219--295 and \path{ExistsK1PointCongruent.lean}:11--52.

\paragraph{Testing the level on two actual points.}\phantomsection\label{r13:s9-twopoints}
Write $g=ykt$, $z=(kt)^{-1}x(kt)$, so the two representing adeles are
$g,gz$. Their unitarity and that of $k,t$ imply that $y,x,z$ are
unitary. The old-level hypothesis and $t\in K_1\subseteq\mathrm{levU}$
give $x\in U\iff z\in\mathrm{levU}$. For $h$ in the old residue group,
the unframed decorations obey
\[
 p'.\mathrm{deco}=(1,h)p.\mathrm{deco}
 \quad\Longleftrightarrow\quad z\in J_0,\quad\mathrm{red}_G(z)=h^{-1}.
\]
Indeed right multiplication by an integral unitary element changes
the pullback coordinates by its inverse. Conversely the equality of
decorations gives equality of the underlying lattices. Their completed
spans are $g\widehat{\mathcal O}_E^n$ and
$gz\widehat{\mathcal O}_E^n$, hence $z$ preserves the standard
completed lattice. Recovery of the completed span follows by
approximating each local basis column by rational lattice vectors
modulo the maximal ideal and applying Nakayama. One imposes integrality
at all other primes by clearing denominators outside the given prime.
The pullback reduction on the lattice surjects onto the residue
space. Evaluating its last coordinate on lifts of each standard
vector gives the prescribed mirabolic last row at $S_T$. Evaluating
all its coordinates gives the matrix $h^{-1}$ at $S_G$.
This proves both implications with their inverse convention
(\path{AdelicLemmas.lean}:607--696, \path{IsPtOfMulJ0Decoration.lean}:13--22,24--50).

If the lattices agree, their two rational frames give
$D_P=(p.\mathrm{fr}_P)^{-1}p'.\mathrm{fr}_P\in
\mathrm{GL}_n(\mathcal O_{E,P})$. Their exact frame equations give
$\rho_P(z)=\iota_P D_P=\rho_{T,P}(t,x)$, the last equality using
$\rho_P(k)=1$. For rational integral $D_P$, its image belongs to
$\mathrm{locLev}_P$ precisely when $D_P\in\Theta_P.H$.
Indeed an admissible lift $D_0\in\Theta_P.H$ differs from $D_P$ by
an element congruent to $1$ modulo $d_P$, thus of depth at least $N_P$;
the assumed principal subgroup lies in $\Theta_P.H$ with trivial
character. Conversely take $D_0=D_P$.
Therefore $x$ in the restricted level gives successively
$h=\mathrm{red}_G(z)^{-1}$ and $\theta_P=D_P^{-1}$. These give
$p'=(1,h,\theta)p$, since the frames act on the right by
$\theta_P^{-1}$. In reverse this equality gives the old-level
criterion and every required local test. This proves the full
\path{GlobalTheory.T.mem_framedLevel_iff_two_adelic_points} (5,394),
including both directions for arbitrary supplied $k,t$.
See \path{TwQLevU.lean}:124--134 and
\path{MemFramedLevelIffTwoAdelicPoints.lean}:14--124.

\paragraph{Realizing one primed family and retaining all its readings.}\phantomsection\label{r13:s9-shrink}
Apply the shared decorated-point construction
\path{PtExist.exists_isPtOf_cm} to each representative of the
\emph{given} step, with $S_T=\varnothing$, $S_G=v_{Gs}$ and $Q=I.Q(a)$.
The old split condition separates the old primes from their partners.
The twist avoidance condition separates the new ones from both old
sets and their own partners. The common construction starts from the
unimodular genus realization, lifts the given split residue matrices
and then splices the prescribed frames. Hence it gives
$y_c\leadsto p_c$ for these particular $p_c$. Choose this family once.
Set $d_P=q_P$ and $|q_P|_P=\exp(-N_P)$. A principal element of depth
$N_P$ reduces to $1$, so lies in the prescribed row group and has
trivial row character. For the level $U'$ above, sufficiently small
identity neighborhoods satisfy $x_P-1\in d_PM_n(\mathcal O_{E_P})$,
with $x_P,x_P^{-1}$ integral, simultaneously at the finitely many
marked primes. Thus $U'$ is open, and closed inside the old compact
group since it is an open subgroup. Use the image of $U'$ at $v\mid Q$
for its new local group and the old local group elsewhere. The local
membership test above proves equality with the resulting global group.
The old base lattice, its index and all marked functionals are retained.
The old-level hypotheses away from $Q$ are inherited, since $T=1$.
These are the fields of the smaller step level
(\path{ExistsStepLevelSmallerAtFinitePlaces.lean};
\path{QbKLeLocLev.lean}:60--129).
The primed step's additional set coordinate is always empty. Its
orbit representatives therefore meet every joint orbit and distinguish
them by equality in the orbit quotient. Its stabilizer consists precisely
of triples fixing $p_c$. Projection to the rational component is
injective by the supplied step stabilizer law applied to $s^{-1}s'$.
The two-point equivalence, congruence correction and double-coset argument
above now identify the required quotient and rational stabilizer
intersection. Since $k=t=1$, the adelic span of the lattice at $c$ is
$y_c\widehat{\mathcal O}_E^n$, from its point property; $S_T$ empty
identifies the derived lattice with this representative lattice.
See \path{TuplePrimedSpaceClassesLevelAndComparison.lean}:25--82,131--209.

At the reference index $Q=\varnothing$. Choose a representative
in the identity double coset. Multiplying it by $x$ leaves no extra
frame to match. The old reading and the unframed test above identify
membership in its level, and on a move by $(1,h,1)$ its character
is $\sigma(\Lambda.\mathrm{chr}(h))$. On a move by $(1,h,\theta)$
in the primed family the additional completed components are
$\iota_P(\theta_P^{-1})$. Their local character is
$\lambda_P(\theta_P)^{-1}$ by the principal-kernel test.
The inverse product defining the additional character gives
$\sigma(\Theta.\mathrm{chr}(\theta))$. Thus the level character
takes the required value on this actual move. The old character factors
through the finite old residue character and all the new local factors
have bounded finite orders, so their product has a common positive
exponent. The added factor is trivial when the components over $Q$
are identities; generator places are outside $Q$. This proves all six
reading clauses, the base-lattice equality, the finite exponent and the
two generator comparisons on the same $y_c$
(\path{LevelFrameCharacterMulMovedPoint.lean}:25--36,38--60,62--72,74--86,88--93,
\path{TuplePrimedSpaceClassesLevelAndComparison.lean}:84--120,212--267).
In conjunction with the shared genus and lift construction, this proves
\path{GlobalTheory.T.tuplePrimedSpaceClasses_level_and_comparison}
(24,746 core lines).

\paragraph{Putting the values and their metrics on that family.}\phantomsection\label{r13:s9-values}
For the row weight, the first $k-1$ entries use the initial parameter.
For the remaining entries the quotient of $i-(k-1)$ by $t+1$ is less
than $m$, since $i<N_m-1$ and $N_m=k+m+m t$. Both parameter lists
use the same block and within-block entry. Their $v,b,e$ coordinates
are independent of $u$, and subtracting the same $\rho$ gives the row
weight. The shifts agree by the oddness calculation above. In these
coordinates the frame equivalence is the identity.
See \path{RowFormsWeightEqParameterSubRho.lean}:13--43 and
\path{EWtApply.lean}:14--23.

Evaluation at the chosen $y_c$ identifies complex automorphic forms
with the allowed row frame functions obeying the stabilizer equation.
In fact extend from $y_c$ to $\gamma_f y_c\kappa$ using rational
transport and the level character. The double-coset bijection gives
the unique class. Two decompositions within it differ by a rational
element of the stabilizer intersection. Its preimage $s$ has the same
rational matrix by the stabilizer reading, and the two-point character
reading gives $\chi(\kappa)=\sigma(\mathrm{chr}(s))$ when
$y_c=\gamma_f y_c\kappa$. Thus the stabilizer condition is precisely
the independence of the extension. Conversely each such relation applied
to a form gives the condition at $c$. For the rational test, each
class polynomial must be obtained by embedding a $K$-polynomial in the
given Weyl subspace. This is exactly membership in the defined rational
span, since it is the image of that subspace. Conversely choose one
preimage per class. The equivariance equations already hold after
evaluation and hence over $K$ by injectivity of evaluation on all
invertible frames; the chosen coefficients give a single rational point.
See \path{ExistsEmbFIff.lean}:29--41,59--97,101--112 and
\path{TupleChainComplexFormsEqAlgebraicModularForms.lean}:20--83.
The integral test on this very chart was proved above.

For the inner product, the Gram matrix is $1$, so every orthonormal
frame $o$ is a tuple of unitary matrices. The determinant shifts
contribute factors of absolute value $1$. The normalized Fock form
is invariant under each unitary change: in each homogeneous tensor
degree substitution acts by a tensor power of a unitary matrix, and
on symmetric tensors its induced pairing is the factorial-weighted
monomial pairing. The fixed highest-weight normalizing constants
are unchanged because the shifted weights have the same partition.
Thus the frame pairing for class polynomials $P_c,Q_c$ is
$\mathrm{nfAll}(Q_c,P_c)$ in the source order. Sum with the given
class weights. For the norm, the rational Gram form is the orthogonal
sum of the class pairings divided by their stabilizer orders. On the
rational structure its image under the embedding equals the normalized
Fock form, by coefficientwise conjugation there. The positivity
of this class sum identifies the archimedean majorant with it.
The two choices of embedding for the same infinite place give
conjugate diagonal values and hence equal real parts. Injectivity on
frames identifies the supplied $P_c$ with the actual coefficient at $c$.
For the unitary invariance, decompose the polynomial into homogeneous
degrees. In degree $d$ the symmetrized tensor for a monomial of
multiplicities $\alpha$ has squared norm $\prod_j\alpha_j!/d!$ for the
product Hermitian form. Distinct multiplicities are orthogonal.
Multiplying this form by $d!$ gives the factorial monomial pairing.
The tensor power of a unitary acts unitarily and commutes with
symmetrization, hence preserves this pairing. Products over the
embedding blocks and division by the fixed highest-weight norm
preserve the identity. See \path{nfAll_blockAct}
(\path{Formalism/LocalToGlobal/Columns/PolyFnLemmas.lean}:628--646).
The same unitary invariance and determinant calculation therefore
give the prescribed norm at every embedding. See
\path{UnitaryOfIsON.lean}:91--170 and
\path{FormsXTupleChainOfLevelReading.lean}:18--91.

Choose the resulting valuation records $V_{u,a}$ on $K_{u,a}$. For each
nonreference index and each required rank bound apply the tilted proof
above to exactly that $K_{u,a}$ and the supplied top reading. Take the
identity row equivalence; equality of underlying frame matrices gives
the valuation evaluation equation. Take $j_M$ from the conjugate period
argument in the preceding round (\path{ch_P_r11.tex}, period comparison;
\path{CevInjective.lean}:128--225). Its surjectivity onto the point
space, image containment and period identity are on the very same
row space. These choices give respectively the classes, base lattice,
finite character order, valuations, norm and tilted datum of the point
record. The associated evaluation record uses their tilted evaluation,
the identity and this $j_M$. Its generator conditions are the two
comparisons just proved; the four standardness fields are the primed
reading, the frame evaluation, the identified tilted vector and the
identity row map. This proves the original supplier
\path{GlobalTheory.T.tupleLayerCompatibilities_of_X_prime}
(34,625 core lines) on its
stated hypotheses, including its already supplied top-point witness.
It uses the shared decorated-point construction and the
\path{Formalism.L2.DIV.nabDelLoc} calculation of
Section~\ref{r13:nabdelloc}. See \path{TupleLayerCompatibilitiesOfXPrime.lean}:17--46.

% blueprint/ch_P_r12.tex lines 1581 through 2223; r12:s8
\subsection{Supplier 8. Reduction maps and moves at a level prime}
\label{r12:s8}
Keep $I$, its correctness and pins, the laws $l,l_0$, the later data
$Z$, the family $H,F$, and the fixed twisted charts $T,J,W$.
Assume $\mathrm{StdF}$ and $\mathrm{StdLvC}$ for every word and
$\mathrm{StdW}$ for this $W$. The conclusion is a single family of
level data satisfying $\mathrm{StdLv}$ together with a witness to
$\mathrm{TwLvSameS}(W)$.
The projections, functionals and moves below all use the same classes
and representatives of $W$. Sources: \path{tw_lv_tied_and_same}
(\path{TupleChainsLevelStructureCompatibleAtLevelPrime.lean}:94--134),
\path{exists_lvTied_of_std} (same file, 17--40).
Bare filenames in this subsection refer to the Iwasawa decomposition
directory unless a different directory is given.

\paragraph{The stage law used in upward naturality.}\phantomsection\label{r13:stage-unit}
In \path{UpChart.upNatTie} (\path{InvSmulMem.lean}:307),
the matrix of the rational part $\gamma$ of a stage move must belong
to the unitary group for the displayed Gram matrix.
\path{step_gam_mem_UJ} (\path{StepGamMemUJ.lean}:10--16)
uses the identity Gram comparison of supplier~6, then projects the
single field \path{unit} from \path{ColIn.Ok.stp_ok}.
For the matrix $M$ of $\gamma\in U(E^n)$, its $(a,b)$ entry in
$M^*M$ is the Hermitian product of $\gamma e_a$ and $\gamma e_b$.
The defining isometry gives $\langle e_a,e_b\rangle=\delta_{ab}$
for every $a,b$, so $M^*M=1$. This is \path{ufGL_conjTranspose_mul}
(\path{FiniteOrbitsOfFiniteFibres.lean}:28--45,89). Changing the
preferred representative preserves this field verbatim in
\path{GStep.Ok.reRep} (\path{TrueColumnChainsSatisfyConditions.lean}:18--20).

For the source ledger, \path{ColIn.Ok.stp_ok} rebuilds the record
from \path{unrepresented_step_conditions}
(\path{TrueColumnChainsSatisfyConditions.lean}:35--50).
\path{StepOk.stepY_ok} calls \path{FramedForms.main} for its finite
orbit and stabilizer fields
(\path{FiniteOrbitsOfFiniteFibres.lean}:65--103;
\path{ExistsSmulOfDecoEq.lean}:83--152).
In the Z4.4 cone the sole caller of \path{stp_ok} is
\path{step_gam_mem_UJ}. Those additional record fields are proved
in Lean. The hand proof of the move uses the matrix computation above.
The field and prime laws also enter elsewhere as facts of the base
data: \path{FieldsLaws.fields_ok} gives the field and embedding laws,
\path{BasePrimes.primes_ok} gives the two generator and level-prime
laws, and \path{base_fOk} and \path{base_fields_primes_ok} package them
(\path{BaseFieldsDefiniteLaws.lean}:114--165;
\path{ProdAbsNormDivsDvd.lean}:148--182; \path{BaseFieldsPrimesOk.lean}:15--22).
The union of these eight named source entries has 12,289 inclusive
core lines. Removing them from the Z4.4 root traversal removes 4,485
core lines; removing \path{stp_ok} alone removes 3,832. The difference
is 653 additional lines in the combined cut. The figures count source
dependencies and do not count extra hypotheses in the upward move.

\paragraph{Two residue projections.}
Fix a word, a tuple, a stage $i\le m$, an oriented prime $P$ above
the rational level prime $q$, and a representative with decoration $d$.
Write $R=\mathcal O_E$, $k=R/P$, $N=d.N$, $D=d.L$ and $D^\sharp=d.Lsh$.
The down lattice is $N_-=d.\mathrm{dnN}(P)$.
For $A_P=M_P$ inside the ambient vector space define
\[
 p_{A,M}:A\longrightarrow M/PM,\qquad
 x\longmapsto \bar s^{-1}[sx],\quad s\notin P,\quad sx\in M.
\]
Cross multiplication gives independence of $s$, additivity and
compatibility with $R\to k$. Given $y\in M$, choose $s\notin P$ with
$sy\in A$, then $t\in R$ with $ts\equiv1\pmod P$. The element $tsy$
lifts $[y]$. The kernel is $PA$: reduction zero means $x\in P A_P$.
Clear denominators to get $sx\in PA$, then
$x=t(sx)+(1-ts)x\in PA$. The other inclusion follows by reduction.
On $A\cap M$ the map is the ordinary quotient map.

Apply this twice, to $(D^\sharp,N)$ and $(D+PD^\sharp,N_-)$.
Here are the local identifications. The decoration includes $N\subset
D^\sharp$, and the separating element, outside $P$ and divisible at
the conjugates of the marked primes, sends $D^\sharp$ inside $N$.
Furthermore $D_P$ is the kernel of the residue functional $f_P$ on
$N_P$. Indeed membership before localization means membership in $N$
and vanishing at every marked prime; after clearing a denominator at
$P$ one can multiply by factors in each other marked prime and outside
$P$. This proves the reverse kernel inclusion. The down construction
takes the kernel lattice at $P$ and adds the correction at $\bar P$ for
unimodularity. More explicitly, if $K_P^0=N\cap\ker f_P$, then
$PN\subset K_P^0\subset N$, $K_P^0\subset N_-$ and
$\bar P N_-\subset K_P^0$. The last containment follows on the added
summand $P(K_P^0)^\sharp$ by dualizing $PN\subset K_P^0$ and using
$N^\sharp=N$. Since $\bar P\ne P$, $N_{-,P}=D_P$.
Finally $PD_P^\sharp\subset\ker f_P$ by residue linearity. Thus
$(D+PD^\sharp)_P=D_P$. The same denominator argument yields, for
$z\in D^\sharp$,
\[
           f_P(z)=0\quad\Longleftrightarrow\quad z\in D+PD^\sharp .
\]
in the forward direction choose $sz\in D$, $s\notin P$ and write
$z=t(sz)+(1-ts)z$ as above.

In the coordinates of confirmed supplier 6 the two source lattices of
the required projections are exactly $D^\sharp$ and $D+PD^\sharp$.
The specified residual functional agrees with $f_P$ on $N$, hence
on $D^\sharp$ after clearing a denominator. For its value-one test,
choose the distinguished lift $e\in N$ with $f_P(e)=1$ and
$qe\in N_-$ representing the prescribed target $t_e$ (the level law).
If $f_P(z)=1$, then $z-e\in D_P=N_{-,P}$ and $q\in P$, so
$p_{D+PD^\sharp,N_-}(qz)=[qe]=t_e$. This gives both surjections,
their kernels and scalar actions, the functional-zero criterion, the
value-one test, and the standard representatives required by
$\mathrm{IdL}$. This is the implication \path{TC.exists_level_projections}
(\path{ReductionModPrimeLevelPrimeSmulEq.lean}:197--231), with the
residue calculation in \path{proj_surj}, \path{proj_ker},
\path{f_eq_zero_iff}, \path{reductionModPrime_levelPrime_smul_eq}
(same file, 43--81, 137--159). For the two lattice identities see
\path{Deco.Ok.dnN_spec}, \path{Deco.Ok.mem_locP_L}
(\path{Formalism/LocalToGlobal/Columns/DecoFin.lean}:653--690,
1192--1229) and \path{Deco.Ok.locP_Lsh}
(\path{Formalism/LocalToGlobal/Columns/BtExistsSmulLeOfFg.lean}:120--127).

\paragraph{The quantifiers and the witnesses for the two moves.}\phantomsection\label{r13:move-fields}
Fix $I$ with $I.\mathrm{Ok}$ and the pins, $l,Z,u,a$, $i\le m$, a bad set,
a label set and class map, a level chart $K_\beta$ for each label, and the
maps $e_Y,e_V$ and compatible lattice record $L$. Assume
$\mathrm{ReadsK}(K)$, $\mathrm{IdY}(e_Y)$, $\mathrm{IdV}(e_V)$ and
$\mathrm{IdL}(L)$. The chart reading supplies one right factor $b_\beta$
for all its representatives $y_c$. Write $P$ for a marked level prime
and $v=e_V(P)$ for its restriction. The kernel $P_v$ of the residue
map is the ideal of $P$. For $c$ of label $\beta$, put
\[
  p_c=L.\mathrm{pr}(P,c),\qquad k_c=L.\mathrm{prk}(P,c),\qquad
  f_c=\mathrm{fv}(i,P,c),\qquad D_c=\mathrm{lat}(c).
\]
Use the coordinate identifications of supplier 6, so $D_c$ is the
decorated kernel lattice and $D_c^\sharp$ is its dual for Gram matrix $1$.

In the up implication take $z\in D_c^\sharp$ with $f_c(p_c z)=1$;
in the down implication take any nonzero residue-linear
$g:N_{-,c}/PN_{-,c}\to k$. Write $c'$ for the resulting up or down
class. In either case one must produce the equality of labels
$\mathrm{chOf}(c')=\beta$, a rational unitary $\gamma$, an adelic
unitary $\eta$ with $\eta_{v'}=1$ for $v'\ne v$, and
$\kappa\in U_\beta$ such that
\[
   \gamma y_{c'}=y_c\eta\kappa,\qquad
   M^{\mathrm{move}}_T(h)(s)=\chi_\beta(\kappa)^{-1}h(s')
   \quad\text{if } e_Y(s)=\gamma_\infty e_Y(s').
\]
The second equality must hold for every test function and every such
frame pair.  Up additionally requires the two localized lattice
identities below, and, for each $x\in D_{c'}^\sharp$ with
$q\gamma x\in D_c^\sharp$, the equality
$f_{c'}(p_{c'}x)=f_c(p_c(q\gamma x))$. Down requires one localized
lattice identity and, for each $x\in D_{c'}^\sharp$ with
$\gamma x\in D_c+P D_c^\sharp$, the equality
$f_{c'}(p_{c'} x)=g(k_c(\gamma x))$.
These are \path{UpNatTieG} and \path{DnNatTieG}
(\path{Formalism/LocalToGlobal/HermitianForms/TiedTwLv.lean}:68--87,91--106).
In both proofs we take the rational part of the actual move
from the chosen representative of $c'$ to the moved point. It is
unitary by \path{step_gam_mem_UJ} (\path{StepGamMemUJ.lean}:10--16):
with the Gram matrix $1$ its defining equation is
$\gamma^*\gamma=1$, the stage law's \path{unit} field.

\paragraph{Downward naturality on the two lattices.}
Write $c'$ for the class reached from $c$ by a nonzero $k$-linear
functional $g$ on $N_-/PN_-$. Let $\gamma$ be any specified rational
unitary matrix equal to the rational part of the chosen move from the
representative of $c'$ to the down point. The required equality is
\[
 (\gamma D_{c'}^\sharp)_P=(D_c)_P .
\]
The new decoration uses $N_-$ and the extension of $g$ by reduction:
on $y$ in its localization put $\bar s^{-1}g([sy])$, and put zero
off that localization. This is additive and residue-linear on its domain;
lifting an element on which $g$ is nonzero proves nonvanishing. Thus the
down point satisfies the decoration conditions (the lattice is the
unimodular neighbor, the other marked components keep their data).
At $P$ its dual kernel lattice localizes to its underlying lattice
$N_{-,P}$ by the first identification above applied to the new
decoration. Transport by the specified $\gamma$ takes $D_{c'}^\sharp$
onto that dual kernel lattice: the representative is carried to the
point, its lattice is equivariant, and unitarity commutes with taking
the Hermitian dual; the inverse inclusion uses $\gamma^{-1}$.
Now $N_{-,P}=(D_c)_P$ by the second identification. Compare
\path{dn_locP_lat}, \path{ok_dn_extF}, \path{locP_Lsh_dn} and
\path{mulLat_Lsh_rep} (\path{DownMoveLocalDualLatticeEqLevelPrime.lean}:
39--81, 88--120).

There is a separate functional equation, requiring
$x\in D_{c'}^\sharp$ and $\gamma x\in D_c+PD_c^\sharp$:
\[
 \mathrm{fv}_{c'}(p_{D_{c'}^\sharp,N_{c'}}x)
     =g(p_{D_c+PD_c^\sharp,N_-}(\gamma x)).
\]
Choose $s\notin P$ with $sx\in N_{c'}$. Equivariance of the underlying
lattice puts $s\gamma x$ in $N_-$; equivariance of the decorated
functional identifies its value on $sx$ with $g([s\gamma x])$ by the
definition of the down decoration. At these two integral lifts both
projections are ordinary reduction by $\mathrm{IdL}$. Cancel the
nonzero factor $\bar s$ using scalar compatibility. This proves
\path{TC.dn_level_functional} with the same $\gamma$ as the lattice
equality (\path{MoveDownLevelFunctionalEq.lean}:24--82, 84--130).

\paragraph{Upward naturality and the adelic factor.}
For an upward move take $z\in D_c^\sharp$ of residual functional
value one and lift its projection by $x_0\in N_c$.
Then $z-x_0\in PD_c^\sharp$. In $R_P$ the element $q$ generates
$P R_P$ (valuation one), so $q^{-1}(z-x_0)\in D_{c,P}^\sharp$.
The local upper neighbor adjoins the line $R_Pq^{-1}x_0$ to $N_{c,P}$.
Its functional is $y\mapsto f_P(qy)$, whose kernel there is
$N_{c,P}$ since $f_P(x_0)=1$. Its dual kernel lattice at $P$
is the localized underlying neighbor. Consequently for the rational
unitary move $\gamma$ the two equalities are
\[
 (\gamma D_{c'} )_P = D_{c,P}^\sharp,\qquad
 (\gamma D_{c'}^\sharp)_P =
        D_{c,P}^\sharp+R_Pq^{-1}z .
\]
For clarity, the local formula for the upper lattice before
imposing its conjugate condition is
$N_{c,P}+(P R_P)^{-1}x_0$. The conjugate condition does not change
this localization: the actual neighbor lies inside the preliminary
lattice and contains its multiple by $\bar P$, and some element of
$\bar P$ is a unit at $P$. On this sum the functional sends
$n+tq^{-1}x_0$ to $\bar t$, so its kernel consists precisely of
$N_{c,P}$. This proves the first displayed equality. For the second,
the localized dual of the kernel lattice of any of these decorations
is its underlying localized lattice. The two possible last generators
$q^{-1}x_0$ and $q^{-1}z$ differ by an element of $D_{c,P}^\sharp$,
so their sums with that lattice coincide in both directions.

To prove the functional clause fix $x\in D_{c'}^\sharp$ in the
specified domain. Choose $x_1\in N_{c'}$ and $x_2\in N_c$ representing
$p_{c'}x$ and $p_c(q\gamma x)$ respectively. Compatibility on integral
representatives and the kernel of $p$ give
$x-x_1\in P D_{c'}^\sharp$ and $q\gamma x-x_2\in P D_c^\sharp$.
By the second displayed lattice identity,
$q\gamma D_{c'}^\sharp\subset N_{c,P}$, since $x_0\in N_c$ and
$D_{c,P}^\sharp=N_{c,P}$. Thus
\[
 q\gamma x_1-x_2=(q\gamma x-x_2)-q\gamma(x-x_1)\in P N_{c,P}.
\]
The old functional vanishes on this submodule by residue linearity.
At the integral representative $x_1$, equivariance and the definition
of the new functional give
$f_{c'}([x_1])=f_{P}(q\gamma x_1)=f_P(x_2)=f_c([x_2])$.
Replacing the two quotient classes by their defining projections gives
the asserted equality. See \path{lattices_at_up}
(\path{ExistsNotMemConjMem.lean}:29--118),
\path{loc_sup_le}, \path{deco_up}, \path{sub_mem} and \path{upNatTie}
(\path{InvSmulMem.lean}:40--129, 262--329).

\paragraph{Splicing a move and reading its character.}
Let $b$ be the fixed right factor for the chart and $p$ the
representative of $c$. Write $p_m$ for the actual moved point and
$p'$ for the representative of its class $c'$. By the definition of
the transport to a point in an orbit there is a specified group
element $a=(\gamma,j)$, with generator and twist parts collected in $j$,
such that $a p'=p_m$. Its rational part is unitary for Gram matrix $1$
and is the rational matrix in the two local-lattice identities.

First compare the chart labels before using a common $b$. The upward
move has the same four invariants as the downward move proved below:
its generator-place lattice is unchanged since only the neighbor
primes $P,\bar P$ change, its generator maps are unchanged by
definition, and the underlying frame and marked subset are kept by the
point move. Hence the simultaneous-translation and preferred-orbit
argument below gives the same chart label for $c,c'$ in both directions.
Also $p,p_m$ have the same coupling coset before choosing a new
representative. A purely rational move preserves this coset on any
point. At a twist place this follows because the transported explicit
frame and the vector have the same rational left factor. At a generator
place with rational residue characteristic $q_g$, let $F,F'$ be the
chosen source and target lifts of the residue frame. The matrix
$F'^{-1}\gamma F$ preserves the standard local lattice and reduces
to $1$, since the residue frame has been transported by $\gamma$
with generator part $1$. The penultimate marked-set condition gives
$q_g F^{-1}y$ integral for every marked $y$. As $q_g$ is a uniformizer
at this unramified place, $(F'^{-1}\gamma F-1)F^{-1}y$ is integral.
Thus the two cosets of the standard lattice coincide. Away from the
penultimate stage the marked subset is empty. See \path{up_keeps_chart}
(\path{UpKeepsChart.lean}:16--32,38--56,62--66) and \path{cosetIn_smul_unitary}
(\path{CouplingCosetFramesInvariantRationalUnitary.lean}:11--107).

For the upward move apply the common decorated-point construction to
$p_m$ with the tuple's twist set $Q$. The prime laws and $I.\mathrm{Ok}$
give pairwise disjointness of the level, generator and twist sets,
with their union disjoint from its conjugate. Concretely $q=\pi\bar\pi$
belongs to the primes of the oriented level set, while neither the
generator nor twist primes lie above $q$; $\pi$ and $\bar\pi$ are
coprime, and the generator and twist union is split. These are also
the hypotheses of the all-twists version on $S_Q\supset Q$
(\path{exists_ptOf_twists},
\path{ExistsAdelicPointDecoratedLatticeAllTwists.lean}:11--60).
The construction \path{Splice.exists_splice} uses \path{PtExist.exists_isPtOf_cm} directly
on $Q$ (\path{ValuedInvMulLe.lean}:146--272).
Write $h_0$ for that representative of $p_m$ and $h=y_c b$ for
the old representative, restricted from $S_Q$ to $Q$.
Splice $h_0$ and its inverse at $P,\bar P$, and $h,h^{-1}$ elsewhere.
The resulting matrix $h_*$ is invertible componentwise and unitary
since the two replaced components are conjugate. It represents $p_m$:
at the two changed primes it has the required lattice; elsewhere
the old and new local lattices agree. This gives the global lattice
since full lattices are determined by their localizations. At other
level primes and generator primes clear a denominator outside the
prime to compare the equal residue maps on the common local lattice;
the twist frames stay fixed. Then $\delta=b(h^{-1}h_*)b^{-1}$ is
unitary and supported at the base prime below $P$.

Put $u_*=y_{c'}^{-1}\gamma^{-1}y_c\delta$. The matrix $y_{c'}u_*b$
represents $(1,j)p'= (\gamma^{-1},1)p_m$ because a rational matrix
on the left composes inverse pullbacks with its inverse. This point
has the same coupling coset as $p'$ by rational invariance,
the equality for $p,p_m$, and the common chart label. The level test
at $c'$ therefore gives $u_*\in U_\beta$; its character test gives
$\chi(u_*)$ equal to the embedded character of $a$, since the rational
part contributes no level character. Use $\kappa=u_*^{-1}$ in the
required relation $\gamma y_{c'}=y_c\delta\kappa$ and multiplier
$\chi(\kappa)^{-1}=\chi(u_*)$. See \path{isPt_up}, \path{exists_level}
and \path{upNatTie} (\path{InvSmulMem.lean}:158--207,211--260,281--329).

For the downward move put $h_m=\gamma y_{c'}b$, and let $m_0$ equal
$h_m^{-1}y_c b$ at the base prime under $P$ and $1$ elsewhere.
Then $h_m m_0$ represents $(1,j)^{-1}p$. Indeed at the specified
base prime its component is that of $y_c b$; the transformation
$(1,j)^{-1}$ leaves both the underlying lattice and the level functional
there fixed. Away from it use the representative $h_m$ of
$(\gamma,1)p'=(1,j)^{-1}p_m$: the neighbor has the same lattices and
residue readings there, and its twist frames are the original ones.
Equality of all localizations again identifies the global lattice.
The coupling coset of $(1,j)^{-1}p$ equals that of $(1,j)^{-1}p_m$
by the simultaneous-translation comparison below, then that of $p'$
by rational invariance, and finally that of $p$ by equality of chart
labels. Thus the level and character readings at $c$ applied to
$y_c\kappa b=h_m m_0$ give $\kappa\in U_\beta$ and
$\chi(\kappa)$ equal to the inverse of the embedded character of $a$.
Explicitly
\[
 \eta=y_c^{-1}h_m m_0^{-1}h_m^{-1}y_c,\qquad
 \kappa=y_c^{-1}h_m m_0 b^{-1}.
\]
Multiplication gives $\gamma y_{c'}=y_c\eta\kappa$ and $\eta$ is
unitary, supported at the one base prime. In both directions the frame
identity is an equality on every function. At each embedding $\tau$,
$\mathrm{IdY}$ identifies the two frames with their matrices. The
condition $s_\tau=\gamma_\tau s'_\tau$ implies
$\gamma_\tau^{-1}s_\tau=s'_\tau$. The defining transport evaluates
$h$ on these inverse translates and multiplies by the embedded
character of $a$. In the up case that multiplier is $\chi_\beta(u_*)$
and $\kappa=u_*^{-1}$; in the down case it is $\chi_\beta(\kappa)^{-1}$
by the character reading at $c$. This proves the frame clause with
the very same $\gamma,\eta,\kappa$ that gave the factorization.
For down the rational part in that factorization is explicitly the
one taking the representative to the down point, so it applies in
both \path{dn_locP_lat} and \path{dn_level_functional}. For up the
kernel and quotient calculations above use that same chosen move.
These identifications complete \path{UpChart.upNatTie}
(\path{InvSmulMem.lean}:281--329) and \path{DnMove.dnNatTieG_cols}
(\path{ChartTiedLevelChainHost.lean}:48--58).
See \path{isPtOf_splice}
(\path{IsPointOfDecoratedLatticeSplicedAtPlace.lean}:12--54),
\path{pointGen}, \path{move_down_adelic}
(\path{MoveDownLevelPrimeAdelicFactorisation.lean}:38--59,78--119,129--183)
and \path{dnNatTieG_cols} (\path{ChartTiedLevelChainHost.lean}:24--58).

\paragraph{Why the downward move keeps its chart.}
Use the same hypotheses as in \path{dn_locP_lat}, dropping the choice
of $\gamma$. Write $p$ for the old representative and $p'$ for the
down point before taking its class. The chart reading declares two
classes to have equal labels exactly when the index and the function
\path{cosetIn} of their representatives agree. A down move keeps the
index. It leaves the explicit frames, the marked subset $Y$ and the
generator maps $\mathrm{be}$ unchanged. At each generator prime it
also leaves the localized underlying lattice unchanged, since the
neighbor changes the lattice only at $P,\bar P$ and neither is a
generator prime. For either of the alternative definitions in the
source, a failed decoration test simply gives back the original point.
The four equalities just listed constitute \path{SameG}. They survive
a simultaneous group translation: the two frames have the same left
and right factors, the two subsets are taken to their images by the
same rational map, the two maps $\mathrm{be}$ are transformed by the
same conjugation and change of coordinates, and localization commutes
with a rational linear isomorphism. For this last assertion a
denominator $s\notin P'$ clears $x$ after changing coordinates exactly
when it clears its inverse image before changing coordinates.
Consequently \path{frsF} agrees at every prime (at a generator it
uses $\mathrm{be}$ and the localized lattice, elsewhere the explicit
frame); the equal subsets give equal \path{cosetIn}. This holds after
every simultaneous group translation of $p,p'$.

At the penultimate stage a preferred representative realizes a fixed
choice from the nonempty set of attainable \path{cosetIn} values on its
orbit. This set does not depend on the starting point within that
orbit, since multiplication by a fixed group element permutes the group.
The sets for $p,p'$ are equal by the preceding equality after every
translation, so their preferred representatives realize the same value.
At every other stage $Y$ is empty and all reduction cosets are empty.
Thus the labels agree even after replacing the moved point by the
chosen representative. This proves \path{TC.dn_keeps_chart};
see \path{SameG.dn}, \path{SameG.smul}, \path{cosetIn_eq_of_sameG},
\path{GStep.NormalFor.eq}, \path{TC.cosetIn_rep_cl_dn_eq},
\path{MoveDownLevelPrimeKeepsChart.lean}:17--101, 107--202.

\paragraph{The coefficient-lattice test in the distinguished basis.}
Fix a seen place $w$ of residue characteristic $q$ and $P$ among the
selected oriented primes above $q$. The seen-place identification indexes
the embeddings by those primes; write $\tau_P$ for the embedding whose
pullback of $w$ is $P$. At any class the source defining $\mathrm{Bw}$
uses $M=D^\sharp_{\mathrm{ker},P}$, the sublattice of $D^\sharp$
on which $f_P$ vanishes. At $P$ it has the localization of
$D+PD^\sharp$ by the criterion above.
At another selected prime $P'$ it has the localization of $D^\sharp$:
multiplication by an element of $P\setminus P'$ sends $D^\sharp$ into
the kernel and becomes invertible at $P'$. Thus for the tests above $w$,
a semilocal basis $B'$ of $D+PD^\sharp$ supplies the basis at $\tau_P$,
and a semilocal basis $B$ of $D^\sharp$ supplies all others. Such bases
exist: localization of a Dedekind domain at the finitely many primes
over $q$ is a semilocal PID, and these full lattices are finite and
torsion-free. Extending the resulting bases to a single prime has the
stated effect on their spans.

Apply \path{ValLat.mem_localization_delAtr} through
\path{Formalism.L2.DIV.nabDelLoc} (Section~\ref{r13:nabdelloc})
to the lattices dual to the scalar extensions of $M$ in each block.
They are full lattices since the base-extended lattice is full and the
dual of its localized basis has inverse-transpose matrix. Thus membership
of $F$ in the localized transported $\Delta$-lattice is equivalent to
$F=\mathrm{act}(A)z$ for one $z\in(\Delta_{\mathrm{std}})_w$ with
$A_\tau$ a dual-frame matrix. After evaluating, its transpose gives
$B'^{-1}$ at $\tau_P$ and $B^{-1}$ elsewhere, with the prescribed
determinant shift. The image of $(\Delta_{\mathrm{std}})_w$ under
the complex embedding is exactly the integral span used by the frame
test: expand an element as a finite linear combination of integral
word tensors with coefficients integral at $w$. Conversely clear
denominators outside $w$ in any such finite sum. The evaluations lose
no information. After cancelling the nonzero determinant powers,
a polynomial identity on every invertible frame holds on every
matrix: for fixed matrices add $tI$ in each block; only finitely many
$t$ are excluded, hence the polynomial in $t$ vanishes also at zero.
Injectivity of the coefficient embedding gives equality over the
coefficient field.

For the forward test this construction works for each pair of bases;
for the converse use their existence and apply the same equivalence
to one such pair. At finitely many classes local denominators may
differ; their product lies outside $w$ and clears all classes at once.
The two shifts agree since $t=2k+1$ gives $i(t+1)/2=i(k+1)$,
and $\mathrm{IdY}$ identifies the frame coordinates. This proves the
implication \path{TC.bw_tie} with the $\Delta$ equality of
\path{nablaDeltaLocal_holds} as its localization input. See \path{mem_closure_iff},
\path{mem_localization_delAtr}
(\path{Development/AbelianVarietyTheory/Proofs/GlobalInterpolation/MemClosureIff.lean}:78--125,141--150);
\path{evFr_injective} (\path{Formalism/LocalToGlobal/Columns/PolyFn.lean}:265--287);
\path{mem_localization_ker}, \path{mem_localization_latticeOf},
\path{bw_tie} (\path{IsBasisAtOfIsLocBasis.lean}:59--123,137--153,180--266).

\paragraph{Identifying the value module once per chart.}\phantomsection\label{r13:fv-field}
The implication \path{LevFun.fvTieG} fixes $I,\mathrm{Ok}$, the pins,
$l,Z,u,a,i\le m$, the bad set, chart family, $e_V$ and $L$. It uses
the chart reading, $\mathrm{IdV}$ and $\mathrm{IdL}$ and assumes that
for every chart at every $v\mid q$ the level functional kills each
$t'=\pi t$ in its completed dual lattice. Its conclusion chooses
one additive isomorphism $e_{\beta,P}:M_{\beta,v}\simeq k(P)$.
The first required identity is
\[
 f_c(p_c z)=e_{\beta,P}\bigl(f_T(y_{c,v}^{-1}z_v)\bigr)
\]
for every class in the chart and $z\in D_c^\sharp$ whose displayed
inverse translate lies in that completed dual lattice.
The second concerns any rational vector $x$, last index $j$,
integral scalar $r$ and local vector $t$ in that dual lattice with
$t=x_v$ and $r=\bar\pi x_j$ in $E$. For those data one must obtain
$e_{\beta,P}(f_T(t))=\mathrm{res}_P(r)$ with the same isomorphism.
See \path{FvTieG}
(\path{Formalism/LocalToGlobal/HermitianForms/TiedTwLv.lean}:55--64)
and \path{fvTieG} (\path{ChartTiedLevelChain.lean}:300--330).

Fix a chart $\beta$ and an oriented level prime $P$ over $v$.
Write $n$ for the positive rank, $\pi$ for the level element,
$O_v$ for the integral ring in $E\otimes_F F_v$, and
$\rho_P:O_v\to k$ for reduction in the $P$-factor.
The chart reading identifies its base lattice with the diagonal model,
and the coordinate law gives Gram matrix $1$. Its completed dual
$\Lambda$ has integral coordinates off the last coordinate and
$\bar\pi t_n\in O_v$ in the last. Put
$\psi(t)=\rho_P(\bar\pi t_n)$. If $r\in R$ lifts a given residue,
the vector $(r/\bar\pi)e_n$ lies in $\Lambda$ and maps to that residue.
Also $\ker\rho_P=\pi O_v$: the base-element ideal is the product
of the selected distinct primes, of which exactly $P$ lies over $v$;
hence $\pi$ generates the maximal ideal of the $P$-factor and is a
unit in the other factors. If $\psi(t)=0$, write
$\bar\pi t_n=\pi w$ with $w\in O_v$. Then
$t=\pi(w/\bar\pi)e_n+t_0$, where $t_0$ has zero last coordinate
and integral remaining entries. The chart functional $f$ kills the
first term by the given $\mathrm{FunKilled}$ condition and the second
by its base-lattice vanishing law. Thus $f$ factors through $\psi$
as an additive map $g:k\to M$. The field \path{fT_ne} makes $g$
nonzero. The residue norm law \path{vT_norm} gives $|k|=q$ and the
chart field \path{fM_card} gives $|M|=q$ (\path{exists_equiv_level},
\path{ChartTiedLevelChain.lean}:211--250). The kernel of a nonzero map
from a group of prime order is trivial; equal finite orders then give
surjectivity.
Use the single additive isomorphism $g^{-1}$ for all representatives
of this chart at $P$. In particular if $t=x$ comes from a rational
vector and $r=\bar\pi x_n\in R$, then $g^{-1}f(t)=[r]$.

For a representative $c$ write $y$ for its chart matrix and $b$ for
the fixed right factor in the reading. For the prescribed dual vector
$z$, set $t=y_v^{-1}z\in\Lambda$ and choose $s\notin P$ with
$sz\in N_c$ by the first local lattice equality. The decorated point
for $yb$ reads its functional on $sz$ from the last coordinate of
$s b^{-1}y^{-1}z$. In the $P$-factor every coordinate of $y^{-1}z$
is integral: off the last this is the dual-lattice condition and at
the last $\bar\pi$ is a $P$-unit. The last row of $b^{-1}$ is
congruent modulo $P$ to the row with sole entry $\bar\pi$ in its
last position, by the reading of the fixed right factor. Hence the
functional on $sz$ equals $[s]\psi(t)$. The class projection on
$sz\in N_c$ is the ordinary quotient and its residual functional
agrees there with the decoration. By scalar compatibility of that
same projection, cancel $[s]\ne0$ to obtain the required identity
on $z$. This proves both assertions of \path{LevFun.fvTieG} under
its stated hypotheses (\path{ChartTiedLevelChain.lean}:90--150,
152--250, 281--330; \path{Formalism/LocalToGlobal/Columns/RecLevelLocF.lean}:100--173).

\paragraph{The stage-zero comparison for each tuple.}\phantomsection\label{r13:lv-zero}
Fix $I,\mathrm{Ok}$, the pins, $l,Z,l_0,u,a$ and any
$A:\mathrm{TwAtT}$ for this word and tuple. The hypotheses particular
to \path{TC.StepZero.lvSame_zero} are $A.e_{VT}=\mathrm{id}$ and
$A.e_{Fv}^{-1}(\mathrm{res}_{\mu_0}(r))=\mathrm{res}_a(r)$
for every marked place and every $r\in R$. The conclusion is the
existence of a single \path{LvSame} at stage zero with the prescribed
$e_C=I.eCls0$, identity frame equivalence and coefficient map
$e_L=(I.\mathrm{tw}\,l\,l_0\,Z).e0$. In particular $e_L$ is fixed
before comparing the lattices. After identifying the places,
every $t\in k(P)$ is $[r]$ for some $r\in R$. Applying $A.e_{Fv}$
to the second hypothesis proves that it fixes $t$, and hence fixes
each residue unit as well. See \path{LvSame}
(\path{Formalism/LocalToGlobal/HermitianForms/TiedTwSame.lean}:17--44)
and \path{lvSame_zero} (\path{LvSameZero.lean}:335--371).

At stage zero the index is unique, the generator level is the same
for both tuples, the twist-frame groups are full integral groups
with trivial character, and the auxiliary subset is empty. Classify
a point by the orbit of its underlying decoration under the rational
unitary group and the generator-level group. This forgets exactly the
twist frames: after an element $(\gamma,k)$ identifies two
decorations, the two frames $F_P,F'_P$ at each twist prime are related
by $(F'_P)^{-1}\gamma F_P$, integral invertible because the frames take
the standard lattice to the corresponding localized lattices.
These matrices give the permitted twist-frame component. Conversely
any framed equivalence compares decorations in that orbit.
The preferred representative selects the same canonical decoration
in each such orbit: one can realize the chosen member by acting on the
given representative with $(\gamma,k,1)$ and using the empty auxiliary
subset, and preference at zero selects precisely that member.
Consequently corresponding classes have equal representative
decorations even when their frames differ
(\path{StepZeroClassesBijectiveSameDeco.lean}:12--97, 113--193).

Write $c_a=e_C(c)$ and $d_a=d_0$ for that equality of preferred
decorations. Define $e_{Lv}$ on the twisted $N/PN$ by sending the
class of $x$ to the class of the same vector in the reference quotient;
define $e_{Lk}$ in the same direction on $N_-/PN_-$. Both are
well-defined because the submodules, and hence their multiples by $P$,
are equal; transporting by the reverse equality gives the inverses. They preserve the residue scalars
because the residue-field isomorphism fixes every $[r]$, and $R\to k$
is onto. The descended functional and distinguished class are defined
from that same decoration, hence are carried to themselves. For an
up move use the common quotient representative with value one; for a
down move the two input functionals agree after this identification,
so their extensions to the ambient space agree: if $y\in N_{-,P}$,
choose $s\notin P$ with $sy\in N_-$ and evaluate the extension as
$[s]^{-1}g([sy])$ on both sides. Off this localization both extensions
are zero. This applies in particular to every $g\ne0$ and every $g'$
with $A.e_{Fv}(g x)=g'(e_{Lk}x)$ for all $x$. For up, after transport
of equal quotients the two prescribed choices of lift are the same
function of the same residue class. The level move laws
identify the underlying decorations of the new points with the
respective decorated moves. In the branch where a decorated move is
inadmissible the definition returns the original decoration in both
copies. Dilation acts on the common decoration by the same residue unit.
All three moves therefore produce corresponding classes.

To compare their transports, start with the group element carrying
the preferred representative to the moved point in the second tuple.
Retain its rational unitary matrix $\gamma$ and generator-level part.
In the first tuple choose twist components
$F_{\mathrm{mov},P}^{-1}\gamma F_{\mathrm{rep},P}$ as above. Their
integrality follows from the equal moved and representative decorations
on the two sides. The resulting element sends the first representative
to the first moved point, preserves the combined character since the
twist character is trivial, and is the chosen reference transport there.
Indeed if $g_0$ is the constructed reference element and $m_0$ is
the prescribed one, $g_0^{-1}m_0$ fixes the reference representative
and is the identity by $(l.\mathrm{lev}\,u\,\mu_0).\mathrm{free0}$.
Only this reference freeness is needed. Transport of a
function on frames uses only $\gamma$ and this character. This gives
the three transported identities even if the twist components differ
(\path{Development/AbelianVarietyTheory/Proofs/GlobalInterpolation/ExistsForgetsFrames.lean}:144--172;
\path{LvSameZero.lean}:77--149, 156--273).

Finally the given stage-zero column equivalence intertwines the complex
embeddings after relabelling the classes. Frame evaluation is injective,
so the actual coefficient values at corresponding classes agree.
Their kernel lattices $d.\mathrm{LshKer}(P)$ agree by the equality
of decorations. The integral coefficient-lattice membership condition
at every class is therefore identical; use the inverse class bijection
for the reverse implication. Extending scalars at $w$ preserves this
equivalence by expressing localized elements as finite sums with
coefficients integral at $w$ and applying $e_L$ and its inverse
termwise. Here the intertwining of embeddings is the second
clause for the specified witness of $l_0(u,a)$; injectivity can
be seen by evaluating on invertible matrices after cancelling the
nonzero determinant factors. The resulting polynomial equality on
these frames holds everywhere by replacing each test matrix by
$X+tI$ and using infinitely many admissible $t$. This argument
gives the $\mathrm{bw}$ implication in both directions, in particular
for $w$ of residue characteristic $q$ with $\mathrm{seenP}(w)$ true.
The two quotient equivalences and their scalar laws, $\mathrm{fv}$
and $\mathrm{te}$ give six fields; the three class equalities and the
three equalities on transported functions give six further fields.
Together with $\mathrm{bw}$ these give \path{TC.StepZero.lvSame_zero}
(\path{LvSameZero.lean}:19--69, 185--232, 279--371). For the distinguished
tuple at every stage the class and frame casts are identities;
surjectivity of $R\to k$ forces the residue isomorphism to be the
identity there as well. The comparison maps can all be identities:
the two down functionals coincide by extensionality, and equality of
the images of residue units implies equality of the units. See
\path{lvSameMu0}
(\path{TupleChainsLevelStructureCompatibleAtLevelPrime.lean}:44--90).

\paragraph{Assembly for this $W$.}
At $\mu\ne\mu_0$ and $1\le i\le m$, apply the two projection
constructions to the standard underlying lattice of $W$ and use the
fixed chart reading, its $\mathrm{FunKilled}$ law, and the identity
frame and place coordinates to get the value-module isomorphisms,
upward and downward naturality for those very projections. The
localized coefficient test supplies their basis field. The residue
and lattice identifications are exactly the given standardness of
$W$, so the same choice satisfies $\mathrm{StdLv}$. For the second
conclusion use the preceding stage-zero comparison on each tuple and
the identity comparison on $\mu_0$ at every stage. Together these
are the two fields of $\mathrm{TwLvSameS}(W)$. This proves the
implication \path{tupleChains_levelStructure_compatible_at_levelPrime}
(\path{TupleChainsLevelStructureCompatibleAtLevelPrime.lean}:94--134;
\path{TupleChainsLevelStructureCompatibleAtLevelPrimeStatement.lean}:8--9),
using the common decorated-point construction and the integral test from
Section~\ref{r13:nabdelloc}. The review confirms the supplier under these hypotheses.

% blueprint/ch_P_r12.tex lines 2225 through 3170; r12:s1
\subsection{Supplier 1. Compatible standard forms across all words}\label{r12:s1}
The conclusion asks for finite chart labels $\mathrm{Ch}_{u,i}$ with class
maps, a unique label at $m$, and simultaneous witnesses $F_u$ of
$\mathrm{StdF}$ and $\mathrm{StdLvC}$. Moreover $(F_u).\mathrm{gl}$
must be equal for all $u$, and at the reference word it must satisfy
$\mathrm{CharacterCompatible}$. All these quantifiers occur in
\path{GlobalTheory.T.exists_standardFormComparison_symmetricLevel_members}
(\path{ExistsStandardFormComparisonSymmetricLevelMemberStatement.lean}:8--9).
There are 53,609 inclusive core lines.

\paragraph{One generator datum for all words.}
Assume the original $I.\mathrm{Ok}$, pins, laws $l$ and later data $Z$.
For each of the two generator primes $p=g_j$, the prime and splitting
laws give $[E:\mathbf Q]$ primes above $p$ in the quadratic CM extension
$E/F$ (\path{SplitsIn}, \path{SemilocalFrames/SelfDualAt.lean}:18--19).
In the sum of $e_w f_w$ giving $[E:\mathbf Q]$, every summand is at least
one, so every $e_w=f_w=1$. The tower identities give $e_v=f_v=1$ for
$v\mid p$ in $F$, and $e_{w/v}=f_{w/v}=1$; summing over $w\mid v$ gives
two distinct primes above $v$. In particular the residue fields in question
identify with $\mathbf F_p$ by the integer map. For the local algebra
take an integral $a$ in one of these primes and outside the other.
Its norm is in $v$ by contraction. Its trace is a unit at $v$:
otherwise at the second prime the norm would force $\bar a=0$ in
the residue field, and the trace would then force $a=0$ there too.
Thus $a\notin F$ and its quadratic polynomial reduces to
$X(X-\operatorname{Tr}(a))$ with distinct roots. Both simple roots
lift over the complete DVR (successive corrections by a unit derivative);
evaluation at them identifies $E\otimes_F F_v$ with $F_v\times F_v$,
sends $F_v$ to the diagonal, and interchanges factors under conjugation.
Since $\pi_j\bar\pi_j=p$, at least one of the two components of
$\bar\pi_j$ has absolute value $<1$. Take that component first.
See \path{ncard_base}, \path{exists_two}, \path{nonempty_splitCore},
\path{split_place}, \path{CmFieldSplitCompletelyPlaceStructure.lean}:94--125,
137--177,244--266; \path{exists_oriented},
\path{ExistsGLocIdGChiTied.lean}:41--66.

For each $j$ choose a place $w_j$ in the selected set above $p$:
its cardinal is $f>0$ by the rank and prime data. A selected $w$ over a
$v\mid g_j$ has label $j$: distinct rational primes cannot both belong
to a proper ideal. Transport $\psi_{0,w_j}$ through the identifications
with $\mathbf F_p$ and the coefficient embedding $\sigma_0:K\to\mathbf C$.
It is nontrivial, since both maps are injective (the residue identification
is onto). The supplied agreement on integers shows that this one choice
gives the prescribed value of $\psi_{0,w}$ simultaneously at all $w$
above the indicated $v$. For the multiplicative data use the injective
$\omega_j:\mathbf F_p^\times\to K^\times$ supplied by
$\mathrm{Ch.Ok}$. Extend $\sigma_0\omega_j$ to its multiplicative
character of $\mathbf F_p$, then at index $a$ transport its power
$2(a+1)$ to $k_v$. On a unit $c\in k_w$ represented by an integer $d$,
the hypothesis \path{chi_om} gives exactly
$\sigma_0(\chi_{a,w}(c))$ for $a<m$. Since $\mathbf F_p^\times$ is
cyclic and the character injective on units, its order is $p-1$.
If two of the transported powers with $a,b<m$ agree, comparison over
$w_j$ gives equal original characters and \path{chi_inj} gives $a=b$.
At different $v\mid g_j$ they agree on all integer arguments by
construction. Together with $\pi_G=\mathrm{Pr}.\pi_G$ these are the
identities \path{IdG} and the distinctness, agreement and common
primitive character in \path{CharacterCompatible}. No choice here
depends on $u$. The exact candidate is
\path{GlobalTheory.Mem.TypePrimes.exists_gLoc_idG_chiTied},
\path{ExistsGLocIdGChiTied.lean}:194--211; see also
\path{orderOf_ofUnitHom} and \path{exists_gLoc}, same file, 23--33,
106--188.

\paragraph{The fibre of the chart class map.}
Write $d(q)$ for the underlying decoration of a framed point, $b(q)$ for
its tuple of reduced cosets, and $r(c)$ for the chosen representative.
Fix $c_0$ of index $j$, $C=b(r(c_0))$, and put
$H_C=\prod_w (\Lambda_{j,w}\cap\operatorname{Stab}(C_w))$.
Here $C_w$ is a set of residue vectors: setwise stabilization is required.
The three hypotheses of \path{MemChart.cls_bij} are the section
variable \path{hinj} and the arguments \path{hsur}, \path{hnorm}
(\path{ExistsMemberStepCharts.lean}:301--310). Explicitly,
$(d,b)$ determines a framed point; any decoration can be lifted to
a point with any already occurring $b$; and at each fixed index, equal
$b$'s give equal $b$'s after replacing the points by their chosen class
representatives. The map sends classes $c$ of index $j$ with
$b(r(c))=C$ to the decorated orbit of $d(r(c))$ under
$U_F\times H_C$. It is bijective.

Indeed the reduction of the action of $(\gamma,h,t)$ on a framed point
is $b_w\mapsto h_w b_w$: in the definition of the reduced set substitute
$z'=\gamma z$; the rational action cancels against the inverse in
evaluation. Thus a level element transports a representative without
altering its reduction exactly when each $h_w$ preserves $C_w$.
If two target orbits coincide, lift the transporting element
$(\gamma,h)$ to $(\gamma,h,1)$ in the frame group. This matches
decorations and fixes the reduction. Uniqueness from $(d,b)$ matches
the framed points, hence the original classes. Conversely take any
decoration $d$, choose $q$ above it with $b(q)=C$, and form its class at
$j$. By the representative hypothesis $b(r([q]))=C$. Some $(\gamma,h,t)$
transports $q$ to $r([q])$. Its $h$ thus lies in $H_C$ and its action on
the decoration uses only $(\gamma,h)$, giving the desired target orbit.
This proves \path{GlobalTheory.Mem.Col2.MemChart.cls_bij}
(\path{ExistsMemberStepCharts.lean}:310--348; reduction and the lifted
action at 150--207,285--308).

The same calculation identifies the rational image of the point
stabilizer with the matrices $\gamma$ for which there is $g\in
U_F\times H_C$ with rational matrix $\gamma$ fixing $d(r(c))$.
A framed stabilizer supplies such $g$ by forgetting its last component.
Conversely $(g_1,g_2,1)$ fixes both decoration and reduction and hence
the framed point. This proves \path{GlobalTheory.Mem.Col2.MemChart.stM_range}
(\path{ExistsMemberStepCharts.lean}:355--366) under the stated
injectivity hypothesis; the separate injectivity of the stabilizer map
uses the stage's \path{Ok.stab_inj} at lines 368--370.

For the reference application the normalization hypothesis comes from
\path{bar_norm}. Away from the penultimate stage the coset is constant.
At that stage the preferred representative chooses a coset from its orbit
of possible cosets. The set of choices depends only on the orbit:
replacing $q$ by $gq$ permutes the indexing group by right multiplication.
Two points with the same reduced sets have identical sets of reachable
cosets by the reduction action above and the coset comparison.
The same choice function on equal sets gives identical normalized
reductions. For chart endpoints with all labels equal to $\beta$,
a reading of the chart family gives the singleton reading by taking
the witness $k(\beta)$. Surjectivity provides a class over $\beta$.
Every class has this label, so the representative, character and
stabilizer clauses specialize with the very same witnesses.
See \path{prefRep_normal}, \path{cos_iff}, \path{bar_norm},
\path{ExistsMemberStepCharts.lean}:259--274,390--404,405--423;
\path{readsK_one}, \path{MemFiCComp.lean}:147--159.

\paragraph{Realization of a reduction fibre on the reference tuple.}
The two hypotheses about decorations used by the preceding orbit arguments
can themselves be verified in the reference application.
Take a tuple with $Q_E=\varnothing$, and suppose
$q_w=|\mathcal O_E/w|$ is outside $w^2$ for every generator place $w$.
Here $q_w\in w$ and therefore $q_w$ generates the maximal ideal in the
local DVR. With the extra twist set empty the framed decoration has no
further coordinates over $Q_E$, and its remaining marked set is empty
away from $i+1=m$. Thus at those stages the decoration determines the
point and any decoration supplies a new point with the same empty
reduction. At $i+1=m$ the marked set has the form $y+L$ where $L$
imposes membership in $N_w$ at all generator places and $q_w y\in N_w$.
The reduced set at $w$ is the singleton containing $b_w(q_w y)$,
since adding any $z\in L$ changes that value by
$\overline{q_w}\,b_w(z)=0$. In the chosen integral frame the map $b_w$
is coordinate reduction: one may lift a residue basis to the lattice
and use the unit determinant of the transition matrix. Thus the kernel
on $N_w$ is $wN_w=q_wN_w$. If two lifts of the same decoration have
equal reductions, $q_w(y-y')\in q_wN_w$, so $y-y'\in N_w$ at each
$w$. The difference belongs to $L$ by its defining local tests, and the
two cosets coincide. This proves
\path{GlobalTheory.T.columnPoint_eq_of_deco_and_cosetReduction},
\path{ColumnPointEqOfDecoAndCosetReductionStatement.lean}:6--7;
\path{inv_smul_mem_of_be_eq_zero}, \path{eq_of_deco_bar},
\path{ColumnPointEqOfDecoAndCosetReduction.lean}:10--25,62--100.

For existence over a prescribed new decoration $d$, write $v_w$ for the
old reduced vector and choose $x_w$ in the new global lattice mapping to
$v_w$ via $b_w$ (the supplied surjectivity field of $d$).
For each $w$ choose an integer of $E$, $e_w\equiv1\pmod w$, divisible
at every other generator prime $u$ by a power of $u$ at least
$\operatorname{ord}_u(q_w)$. The relevant prime powers and $w$ are
pairwise coprime, so these choices exist by the Chinese remainder theorem.
Put $y'=\sum_w(e_w/q_w)x_w$. At a fixed generator place $u$ the
$u$-term of $q_u y'$ reduces to $v_u$ because $e_u\equiv1\pmod u$.
In every other term $e_w/q_w$ is integral at $u$ and the extra factor
$q_u$ kills its reduction there. All terms of $q_u y'$ lie in $N_u$.
Consequently $y'$ is an admissible centre, and the singleton calculation
shows that its coset for $d$ has exactly the original reduced sets.
This proves \path{GlobalTheory.T.exists_couplingCoset_with_given_reduction},
\path{ExistsCouplingCosetWithGivenReductionStatement.lean}:6--7;
\path{exists_one_intAt_div}, \path{exists_vec}, \path{coset_be},
\path{exists_ptAt_bar_eq},
\path{ExistsCouplingCosetWithGivenReduction.lean}:16--55,59--102,105--134.

\paragraph{The reference chart and all its readings.}
Fix a word $u$ and a stage $i\le m$, write $n=N_i>0$, $T=T(u)$, $\pi=\pi_T(u)$,
$S_T=\operatorname{divs}(\pi)$ and $S_G$ for the generator places.
The prime laws give
$\pi\bar\pi=T$, coprimality of $\pi,\bar\pi$, and
$(\pi)=\prod_{P\in S_T}P$ with $f$ distinct factors of norm $T$.
The generator primes avoid $T$ and are unramified, so the preceding
reduction arguments apply with $q_w\notin w^2$. For the fixed right
factor put $k=\operatorname{diag}(1,\ldots,1,a)$: at $P\in S_T$
take $a_P=1/\bar\pi$, at $\bar P$ take $a_{\bar P}=\pi$, and take
$1$ elsewhere. Coprimality makes the prescriptions disjoint and all
components units. They are conjugate inverses, so $k$ is unitary and
preserves the full standard lattice. It is the identity off $T$, and
the last row of $k^{-1}$ reduces to $[\bar\pi]e_n^t$ at each $P$.
The specified choice $k_{\rm std}$ has these same properties by
existence (\path{ExistsUnitaryDiagonalBasePointLevelPrime.lean}:51--112,
122--214). For each preferred representative $p_c$, extend its empty
twist decoration to the global twist set as prescribed. The common
all-twists point construction supplies $g$ representing this extension.
Then $gk_{\rm std}^{-1}$ is a witness in the definition of $y_c=I.yStd$;
thus this particular choice satisfies $\mathrm{PtOf}(y_c k_{\rm std},p_c^{\rm ext})$.
Since the factor fixes the standard lattice, both $y_c$ and $y_c k_{\rm std}$
give the same full lattice. At $S_G$ its identity component gives the
same frame; at $P\in S_T$ the row computation in the pullback
$k_{\rm std}^{-1}y_c^{-1}z$ multiplies the last-coordinate functional
by $c_T(P)=[\bar\pi]$. Equality on the global lattice extends to its
localization by clearing a denominator outside $P$; off that localization
both decorated maps vanish. Hence $y_c$ is unitary and
$\mathrm{decoOf}(y_c).\mathrm{dia}(c_T)=d(p_c)$
(\path{DecoratedLatticeDiagonalBasePointMul.lean}:47--134;
\path{ExistsMemberStepCharts.lean}:425--452).

Define the base lattice $L_0=\mathcal O_E^{n-1}\oplus\pi\mathcal O_E$.
It is integral, free and full. Its Hermitian dual has last coordinate
$\bar\pi^{-1}\mathcal O_E$ and integral earlier coordinates. Thus
$TL_0^\vee\subset L_0$ and $[L_0^\vee:L_0]=|\mathcal O_E/(T)|=T^{2f}$,
using the $f$ selected residue factors and their distinct conjugates.
For a chart fibre $C$ as above, take as level $U$ the integral unitaries
whose last row at each $P\in S_T$ is $e_n^t$ modulo $P$ and whose
reductions at $S_G$ belong to $H_C$. This is a product of local levels:
restricting a matrix at one base place and putting the identity at all
others preserves the conditions since $H_C$ is a product. These are
finite congruence conditions in the compact open integral stabilizer,
so $U$ is compact open in the finite adelic unitary group. Away from
the selected and generator places the level is the lattice stabilizer.
At $v\mid T$ the canonical functional on the completed dual is the
quotient by the kernel of $t\mapsto[\bar\pi t_n]_P$, at the unique
selected $P\mid v$. Every value lifts by $(r/\bar\pi)e_n$; the map
kills $L_0$ and $\pi L_0^\vee$ and has $T$ values. This gives both the
prime-order quotient and $\mathrm{FunKilled}$ for the charts chosen here.
For its stabilizer, at the selected factor $L_0^\vee$ is the full integral
lattice, and at the conjugate factor it is $L_0$ that is the full
integral lattice. A unitary preserving $L_0$ preserves its dual.
Fixing the quotient on that dual is then exactly the last-row congruence
at $P$, since $\bar\pi$ is a $P$-unit. Conversely integral unitaries
satisfying that row condition preserve $L_0$, its dual and the quotient,
as do their inverses. This identifies the local level with the required
intersection of stabilizers. Finally use the character
$\chi_U(x)=\sigma(\Lambda.\mathrm{chr}(\mathrm{red}_G x))^{-1}$,
which is trivial when the generator components are trivial. To check
every field of \path{StepLevel} (\path{HermitianForms/IsLocBasis.lean}:58--88)
in \path{stepOf} (\path{RecLevelStep.lean}:20--51), group the data as follows.
\begin{itemize}
\item \path{Lb}, \path{Lb_fin}, \path{Lb_span}, \path{Lb_int},
\path{Lb_dual}, \path{Lb_index} come from the displayed coordinate bases
of $L_0,L_0^\vee$, multiplication by $T$, and their index.
\item \path{Uloc}, \path{Uloc_le}, \path{Uloc_lev}, \path{Uloc_hyp}
are the specified subgroups of local unitaries, their preservation of
$L_0$, and equality to the full stabilizer away from $T$ and the bad set.
\item \path{fM} is the order-$T$ quotient just computed, with its
additive group and scalar action (\path{fM_acg}, \path{fM_mod}).
The quotient map \path{fT} is linear, is nonzero by surjectivity, kills
$T L_0^\vee$ and $L_0$, and takes values in a set of cardinality $T$.
These are \path{fT_ne}, \path{fT_T}, \path{fT_L}, \path{fM_card}.
The last-row calculation in both directions gives \path{UlocAlt}.
\item \path{U}, \path{U_eq}, \path{U_compact}, \path{U_open} follow
from the componentwise test and the finite congruence cut inside the
compact integral unitary group. Finally \path{chi}, \path{chi_triv}
are the displayed homomorphism and its restriction off the bad places.
\end{itemize}
See \path{RecLevelLocF.lean}:90--210 and \path{RecLevelLocK.lean}:359--403.
The confirmed coordinate theorem supplies Gram matrix $1$ and identifies
the class lattice with the kernel lattice of $d(p_c)$.

The same $y_c$ gives the double-coset representatives. Indeed if
$\mathrm{decoOf}(x)=(\gamma,h)\mathrm{decoOf}(y)$ then
$x^{-1}\gamma_f y$ is integral with integral inverse, since it carries
the completed standard lattice to itself. Evaluating the functionals
and frames gives its last row $e_n^t$ at $S_T$ and its reduction $h$
at $S_G$, so this quotient belongs to $U$. The converse follows by
applying these identities to the pullback of any lattice vector.
For this local-to-global step choose a nonzero integer $d$ bounding
the adelic image lattice between $d\widehat{\mathcal O}_E^n$ and
$d^{-1}\widehat{\mathcal O}_E^n$, also for the inverse change of
coordinates. The quotient by the lower lattice has only the factors
above $d$. The Chinese remainder theorem identifies its submodule of
prescribed local classes with the image of their global preimage in
$d^{-1}\mathcal O_E^n$. Thus the completion of that preimage gives
exactly the prescribed lattice at every prime. Conversely a rational
vector satisfying all local tests lies in the preimage, so the
double-coset computation detects equality and preservation in both
directions (\path{levU_y_latD}, \path{RecLevelAsm.lean}:166--175).
Diamond multiplication by the unit tuple $c_T$ commutes with the
orbit action and is invertible. The fibre bijection above therefore
gives exactly one $y_c$-class for each double coset. The kernel of
all last-coordinate tests in the completed lattice is $y_c\widehat L_0$.
The stabilizer range above, applied with $x=y=y_c$, is now precisely
$U(F)\cap y_cUy_c^{-1}$ and the given stage law makes the stabilizer map
injective. See \path{chart_exists}, \path{ExistsMemberStepCharts.lean}:52--120;
\path{RecLevelCls.lean}:59--110,435--463.

Take the labels to be the image of $c\mapsto(j,\mathrm{cosetIn}(p_c))$;
this is a finite surjective image of the finite stage class set. Equality
of these labels is equality of $j$ and of the reductions by the preceding
normalization and coset calculation. Pick a preimage of each label for
the definition of $H_C$. For the level reading, when $y_c x k_{\rm std}$
represents a test point, multiplication by $x\in U$ acts on its decoration
relative to that of $p_c$ as $(1,(\mathrm{red}_G x)^{-1})$. This lifts
to a frame operator with rational component $1$ fixing the coset, since
$\mathrm{red}_G x\in H_C$. In the converse direction such an operator
has residue component $h$ fixing $C$, and the pullback comparison gives
$x\in U$ with $\mathrm{red}_G x=h^{-1}$. Its extra twist component is $1$
because the reference tuple is empty. Substitution in $\chi_U$ gives
$\sigma(\Lambda.\mathrm{chr}(h))$, the required character value.
The base-lattice, diagonal, point and stabilizer readings are the ones
just verified, all with the same representatives. At $0$ and $m$ the
index is unique by the stage law and every auxiliary coset is empty since
$m\ge2$. The class set contains the class over the standard decoration,
so in both cases the image is a singleton. This proves the exact
reference theorem \path{GlobalTheory.Mem.Col2.MemChart.exists_charts}
(\path{ExistsMemberStepCharts.lean}:222--251,458--486,521--566), together
with the canonical functional used in this construction. In Lean
\path{MemChart.exists_charts} returns the reading, $y_c=I.yStd$, and
the endpoint uniqueness. At the next call \path{TC.charts_funKilled}
rebuilds each level with \path{Reside.swap}. Its new functional has the
local quotient just computed. The old and new stabilizers agree:
linearity and prime order force the old functional to factor through
one residual line, and unitarity identifies the stabilizer condition on
the paired split lines (\path{Reside.stab_eq}, \path{Reside.swapLaws},
\path{InclMulInvMul.lean}:312--359,362--399). The swap retains
$L_b,U_{\rm loc},U$, the character evaluated on each underlying
element, every $y_c$ and the stabilizer map. Thus it retains the
witness for $\mathrm{ReadsK}$ and $y_c=I.yStd$, and now gives
$\mathrm{FunKilled}$ for the same coordinates
(\path{InclMulInvMul.lean}:422--439;
\path{formsTiedSym_std}, \path{ExistsStandardFormComparisonSymmetricLevelMember.lean}:63--77).
This construction uses supplier~6 and the common point construction.

\paragraph{The reference primed chart.}
The primed stage uses rank $N_m-1$, Gram matrix $1$ (supplier~6),
$S_T=Q=\varnothing$ and the primed residue level $\Lambda$ at $S_G$.
Use the standard full lattice and $U=\mathrm{red}_G^{-1}(\Lambda)$ in
the integral unitary group, with inverse embedded residue character as
above. This is \path{stepOf} with $\pi=T=1$: the dual is the lattice
and its index is $1$, and all the clauses indexed by a place over $1$
are vacuous. There is one stage index, no auxiliary coset and no twist
frames, so the underlying decoration identifies the points and permits
lifting the action of $(\gamma,h)$ by $(\gamma,h,1)$. Apply the common
decorated-point construction to each preferred point with $S_T=Q=\varnothing$;
its representative $y'_c$ realizes that very point. The orbit/double-coset
argument above, now without a diamond correction, gives the class bijection,
the completed lattice identity and the stabilizer range. In this case
the stabilizer projection is also injective directly: when its rational
part is $1$, fixing a decoration forces $h=1$ by surjectivity of its
residue frames. For $y'_c x$ representing a test decoration, the same
calculation gives $x\in U$ exactly when a pure level operator carries
the preferred decoration to the test; its residue is the inverse of that
operator's component. Its character is therefore the prescribed one.
This also gives the stabilizer character when $y'_c=\gamma_f y'_c x$.
These are all clauses of the reference reading $\mathrm{ReadsP}$,
including $L^\sharp=N$ by self-duality since $S_T=\varnothing$.
See \path{UfGLRatHom.lean}:131--150,155--284.

For this reference $y'_c$ use the representative-value construction of
\hyperref[r13:s9-values]{values and metrics} with the identity on row functions.
For a complex value, the equivariant extension is well defined exactly when
the two presentations of $\gamma_f y'_c x$ give the same stabilizer action.
For rational values apply the coefficient embedding to each class
polynomial. For lattice membership use the frame test on the dual class lattice.
These are \path{RowTie.mem_FiC_iff} (\path{UfGLRatHom.lean}:43--108)
and the \path{RowTie} comparisons \path{mem_Fi_iff}, \path{rat_iff},
\path{lat_iff} (\path{ExistsEmbFIff.lean}:29--55,59--141).
For the metric the same representative indexes the summand with weight
$|\mathrm{Stab}(c)|^{-1}$ in \path{petersson_apply};
\path{MetricTie.norm_step} computes its majorant in an orthonormal frame
(\path{ExistsEmbFIff.lean}:146--151,260--293;
\path{UnitaryOfIsON.lean}:111--155). Here the class
lattice is $N$, so $N=N^\vee$ also supplies the dual used in the
integral basis tests. Oddness $t=2k'+1$ gives
$m(t+1)/2=m(k'+1)$ and makes the two determinant shifts identical.
The row entries split at the first $k-1$ entries; afterwards write the
index offset as $a(t+1)+b$, $0\le b\le t$, and use the $b=0$ or
$b>0$ formula in both weight definitions. The strict upper index bound
gives $a<m$. Thus the point, rational and integral function spaces,
Petersson form and norm belong to this very primed chart. The identity
on row functions gives equal values on frames with equal underlying
matrices. This proves \path{GlobalTheory.Mem.Col2.TC.clsP_tie}
(\path{ExistsEmbFIff.lean}:59--114,144--199,209--293), using the
common decorated-point argument and the transported integral lattice
proved in this round. No previously existing primed reading is needed.

\paragraph{Values and metrics on the reference stages.}
More generally, for any tuple and a given read chart with $I.\mathrm{IdY}$,
Gram matrix $1$ and the equality of its class lattice with the decoration's
$L$ come from supplier~6. For $t=2k'+1$ the stage shift is
$i(t+1)/2=i(k'+1)$; thus the two weights and evaluation shifts agree.
For a complex function on a fibre its representative values extend by
$\gamma_f y_c x$ using the character of $x$. In a stabilizing relation
$y_c=\gamma_f y_c x$ pull back the represented point by $\gamma^{-1}$;
its resulting action has rational component $1$, so the given reading
identifies the character of $x$ with the stabilizer character. Exactly
as in the primed case, this is independence of the extension on the
whole double coset. For rationality choose polynomials above the given
rational values class by class. Evaluation on invertible frames is
injective: in a fixed frame it is evaluation of a polynomial on the
dense open of invertible matrix tuples with an invertible determinant
factor. Thus the class invariance conditions descend to the chosen
rational polynomials. The converse is evaluation followed by scalar
extension. For the integral test the Hermitian dual equals $L^\sharp$
in these coordinates. The stabilizer comparison just used restricts the
point reading from $I.SQ$ to $I.Q(a)$: \path{LevTwist.levOk_twists}
gives the separation from the generators, and \path{IsPtOf.of_extTo}
restricts the frame equations to the smaller set while keeping the
lattice and reductions (\path{LevOkTwists.lean}:48--80;
\path{ReferenceCharactersFactorTame.lean}:275--278).
For the integral equivalence at $w\mid\ell$, finite generation and
full rank give a basis $B$ of the class dual lattice at all primes
over $\ell$ (\path{exists_isLocBasis}, \path{PBwLocGen.lean}:248--289).
It is a basis at each prime pulled back from $w$ by an embedding.
Apply the $\Delta$ localization formula in Section~\ref{r12:weyl}:
on the dual frame $(B^{-1})^{\mathsf t}$ it identifies the localized
transported lattice with the action on the localized standard one.
After frame evaluation the transpose gives $B^{-1}$.
Lean uses \path{mem_localization_delAtr_fn}, then
\path{mem_localization_delAtr}, \path{nabDelLoc} and
\path{nablaDeltaLocal_holds}
(\path{IsBasisAtOfIsLocBasis.lean}:26--78;
\path{MemClosureIff.lean}:141--150;
\path{CoefficientsModP.lean}:24;
\path{NablaIndependenceHolds.lean}:578--580).
If $B'$ is another such basis, each transition matrix and its
inverse preserve the corresponding local lattice. Applying the
localization equality to both frames identifies both images with the
same intrinsic localized lattice. Frame evaluation is injective on
these polynomials and the coefficient embedding identifies the
localized standard span with the required integral polynomial span.
This gives the test for every $B'$; conversely one simultaneous $B$
per class suffices. Multiply the finitely many denominators outside
$w$ to clear every class at once
(\path{mem_localization_latticeOf}, \path{IsBasisAtOfIsLocBasis.lean}:137--153;
\path{HermitianChartComplexRationalIntegralForms.lean}:456--510).
This proves \path{GlobalTheory.Mem.Col2.TC.val_tied}
(\path{HermitianChartComplexRationalIntegralForms.lean}:394--510).
For the metric comparison \path{ipV_norm_tie}, the original hypotheses
$\mathrm{Ok}$, pins, laws and later data supply the stage and field laws;
$\mathrm{IdY}$ is the only frame comparison it requires. At each class
apply the orthonormal-frame computation in
\hyperref[r13:s9-values]{the class-value comparison}: with partition fixed under the shift, the unitary
change of orthonormal frame preserves the normalized Fock form and the
absolute value of each determinant is $1$. This gives
$\mathrm{nfAll}(Q_c,P_c)$ for the local inner product.
For the norm, let $K_0=\mathrm{cmMax}(K)$, the subfield on which
the given involution agrees with complex conjugation at every embedding.
The coefficient matrices of the stabilizer actions lie in $K_0$:
the matrices are obtained from embedded source-field entries and integral
structure constants, and the character factor is a root of unity.
Conjugation inverts every such root at every embedding, as it does under
the given involution by the distinguished embedding. Choose a basis of
$K/K_0$. Applying its coordinate projections entrywise commutes with
all stabilizer equations. Thus every invariant form is a $K$-linear sum
of invariant forms with coefficients in $K_0$, and one can choose a
basis $b_i$ of the form space among these. This is
\path{LatGood.hasRatBasis} (\path{WeightLawGoodOfFieldEmbedding.lean}:151--181;
\path{hasRatBasis_of_act}, \path{HasRatBasisOfAct.lean}:58--91).

Write $V_c$ for evaluation of this basis in the coefficient space at $c$
and $H_0$ for its value Gram matrix at an embedding $s$. By the rational
basis property, conjugation on the first slot of the embedded pairing
is ordinary complex conjugation. Its form Gram matrix is therefore
$H_s=\sum_c |\mathrm{Stab}(c)|^{-1} V_c^*H_0 V_c$. The value form
$H_0$ in identity Gram coordinates is positive definite by the normalized
Fock pairing. Stacking the $V_c$ gives an injection: the values at every
class determine the form, and a full-rank minor stays nonzero under $s$.
Each nonzero vector hence has some $V_c x\ne0$, where the term is
strictly positive; all other terms are nonnegative. This proves
\path{LatGood.step_gram_posDef}
(\path{WeightLawGoodOfFieldEmbedding.lean}:69--109,183--193).

For a positive definite $H_s$, the metric in the spectral frame of its
majorant equals $H_s$. Expanding $x=\sum_i a_i b_i$ and its value at
each class in the above Gram identity shows that the majorant is the sum
of the normalized Fock values of $s(x_c)$ with weights
$|\mathrm{Stab}(c)|^{-1}$. This proves \path{majAt_classSum} on this
datum (\path{UnitaryOfIsON.lean}:15--44,95--132). For a prescribed
polynomial family $P_c$, equality on every frame and $\mathrm{IdY}$ imply
$P_c=x_c$ by injectivity. At any allowed orthonormal frame the determinant
factors have absolute value one and the unitary coordinate action preserves
the Fock pairing. The two conjugate representatives of the same infinite
place have identical real diagonal values. This is \path{norm_step} and,
together with \path{ipVal_step}, both fields of
\path{GlobalTheory.Mem.Col2.MetricTie.ipV_norm_tie}
(\path{UnitaryOfIsON.lean}:111--155,159--169,195--205).

\paragraph{The diamond on a fixed chart.}
Take a read chart, a place $v$ corresponding to $P\in S_T$, a unit $s$
and an integer $z$ reducing to $s$. By the factorization in supplier~8,
its nonzero order-$T$ functional identifies with
$t\mapsto[\bar\pi t_n]_P$ whenever $\mathrm{FunKilled}$ holds. For
$\delta$ satisfying $\mathrm{IsDia}(v,z)$ the defining equation
$z f(\delta_v t)=f(t)$ gives $f(\delta_v^{-1}t)=s f(t)$ under this
identification. The element preserves the base lattice and its dual.
At the $P$-factor the latter is the full integral lattice; at $\bar P$
the former is the full integral lattice. Thus $\delta$ and its inverse
are integral, and testing on basis vectors gives the last row of
$\delta^{-1}$ congruent to $s e_n^t$ at $P$. Off $v$ it is the identity.
In the pullback $k_\beta^{-1}\delta^{-1}y_c^{-1}$, the original last
row of $k_\beta^{-1}$ at $P$ is $[\bar\pi]e_n^t$; hence the new
functional is $s$ times the original. At the other selected primes
and at the generator and twist places the pullback is unchanged since
these places do not lie over $v$. Integrality of $\delta^{\pm1}$ and
preservation of the full standard lattice by $k_\beta$ give the original
underlying lattice. This proves that $y_c\delta k_\beta$ represents the
diamond-translated point. Source checks: \path{IsLocBasis.lean}:137--143,
\path{DiamondElementIntegralSupportLastRow.lean}:376--463,469--514,
\path{IsPointOfLastCoordinateDilationOnePrime.lean}:113--210.

The diamond affects neither the marked coset nor the index. Its action
commutes with the stage action, so at the penultimate stage the set of
cosets reachable from the translated point is the original reachable
set. Normalization chooses the same coset; away from that stage it is
empty. Thus the class $c'$ of the translated point has the same chart
label as $c$. Let $(\gamma,\eta)$ be its prescribed mover from the
preferred representative of $c'$, and put
$x=y_{c'}^{-1}\gamma_f^{-1}y_c\delta$. Then $y_{c'}xk_\beta$
represents the action of $(1,\eta)$ on that representative.
Rational transport leaves the framed coset in its own frame unchanged
by applying the inverse frame to $\gamma$ times each vector
(\path{cosetIn_smul_unitary}, \path{CouplingCosetFramesInvariantRationalUnitary.lean}:11--107). The level
and character readings therefore give $x\in U_\beta$ and
$\chi_\beta(x)=\sigma(\mathrm{chr}(\gamma,\eta))$. The identity
$y_c\delta=\gamma_f y_{c'}x$ proves the asserted double-coset equality.
Evaluating a test form using it gives exactly the framed dilation action:
its argument is transported by $\gamma^{-1}$, and $I.\mathrm{IdY}$
identifies the underlying matrices of the two frames. All three clauses
of \path{GlobalTheory.Mem.Col2.DiaChart.dilTied} follow
(\path{DiamondOperatorsOnCharts.lean}:19--52,74--129,131--231).

\paragraph{Tilting the distinguished endpoint.}
We prove \path{TC.til_ties} for the actual top and primed readings.
Fix $u$ at the distinguished tuple $\mu_0$. The premises are $I.\mathrm{Ok}$,
the pins, $l,Z$, the given $A,K,K',L,W$, and these five identities or
readings: $A.\mathrm{res}=I.\mathrm{resT}$, the top reading
$\mathrm{Reads}$ for a common $k$, the primed reading $\mathrm{ReadsP}$,
$L.\mathrm{pr}(x)=[x]$ on the ambient decorated lattice $N$, and
$W.e_Y$ preserving the matrix of each frame.
The inclusion pin sends the $j$th primed basis vector to its namesake in
the top frame. By linearity its coordinate inclusion is $i(x)=(x,0)$.
The cofactor pin is $c_u\pi_u=\pi_{u_0}$; the present construction
uses only one $u$. The standard frame and the identification of its
class lattice with the decorated kernel are the confirmed supplier~6.
See \path{LatAtLattice.lean}:125--128,256--276,292--305 and
\path{TrueColumnSpacesStandardInFrameCoords.lean}:90.

Write $n=N_m-1$, $\pi=\pi_u$ and $z_0=\pi e_{n+1}$.
Then $h(z_0,x)=\bar\pi x_{n+1}$ and $h(z_0,z_0)=T$ by the
norm condition. The given anisotropic perpendicular vector has all
first $n$ coordinates zero by testing against the coordinate vectors,
and its last entry is nonzero, so $z_0$ lies on its line.
For a primed representative $p_c$ the underlying tilted self-dual
lattice is $N_c'\oplus\mathcal O_E$. At a level prime $P$ the
functional on it is reduction of $\bar\pi x_{n+1}$. The factor
$\bar\pi$ is a unit at each selected $P$; their product is $(\pi)$
by the divisibility condition on $\pi$. The kernel lattice of the
tilt is therefore $i(N_c')\oplus\mathcal O_E z_0$. The move
$m$ from the chosen top representative to the tilted point has the
rational unitary component $\gamma$ given by the stage action. The
residue-group components change the frames and preserve the common
kernel; thus $\gamma$ carries its class lattice to this direct sum.
For the base lattice the very same decomposition follows from
the two base readings: the first $n$ coordinates are integral and
the last belongs to $\pi\mathcal O_E$.

For a vector $x$ in the dual of that top class lattice put
$s=\bar\pi$. Then $sx\in N$: the defining kernel lattice contains
$\pi N$ at each representative, and pairing $x$ against $\pi N$
places $\bar\pi x$ in $N^\vee=N$ (conjugate symmetry gives either
ordering of the test). Suppose $r=h(z_0,\gamma x)$ is integral as
in the functional condition. Apply the projection reading to $sx$.
Transport by the move and use the tilted functional above; its value
is $\mathrm{red}_P(sr)$ since $\gamma(sx)=s\gamma x$.
Both the projection and induced functional are residue-linear, and
$\mathrm{red}_P(s)\ne0$. Cancellation proves the required value
$\mathrm{red}_P(r)$ on $x$, including $x$ outside $N$.

For the adelic condition take the primed representative $y'_c$
in its given reading. The block matrix $D=\mathrm{diag}(y'_c,1)$
carries the standard completed lattice to the completed tilted
self-dual lattice by the direct decomposition. The top reading
says $k$ is integral unitary on the standard lattice, is $1$ off
the places over $T$ (there are no twist primes at $\mu_0$), and
the last row of $k^{-1}$ reduces at each selected $P$ to
$\bar\pi e_{n+1}^{t}$. Thus $Dk$ has the same lattice and its
inverse gives the stated level functional. At the generator
primes $T\ne g_j$ gives $k=1$ and $\pi\equiv1$; the first $n$
residue-frame entries are those of $p_c$ and the last one is
$\mathrm{red}(x_{n+1}/\pi)=\mathrm{red}(x_{n+1})$. These are
exactly the tilted decoration fields. Write $y_t$ for the top reading's
representative of the tilted class. Rational equivariance shows that
$\gamma_f^{-1}Dk$ represents the residue-group part of $m$ acting on
that representative: inverse rational transport applies $\gamma$ to
the tested vector and removes the rational component of $m$.
The residual coset tests agree because the permitted top coset set
is empty for every point. Apply the level criterion of $\mathrm{Reads}$
to $b=y_t^{-1}\gamma_f^{-1}D$, for which $y_t b k=\gamma_f^{-1}Dk$.
It yields $b\in K.\mathrm{lev}.U$ and the character reading gives
$\chi(b)=\iota(\mathrm{chr}(m))$; the stage character depends on
the two residue parts. Put $\kappa=b^{-1}$. Then
\[
 \gamma_f y_t=D\kappa,\qquad
 \chi(\kappa)^{-1}=\iota(\mathrm{chr}(m)).
\]
Any candidate extension agrees with $D$ on the first $n$ columns
and fixes $\pi e_{n+1}$. The scalar $\pi$ is a unit in the adelic
algebra as the image of a nonzero field element, so the last column
is $e_{n+1}$ too. Hence this same $\kappa$ works for every allowed
extension. Finally under each complex embedding the frame identity
$s=\gamma s'$ gives $\gamma^{-1}s=s'$. Using the character
identity in the definition of transported frame evaluation proves
$\mathrm{EvzTie}$ with the identity coefficient-space equivalence.
This proves the tilted endpoint implication with all data chosen from
the two stated readings. The primed class and its metric reading are
a separate construction in \path{clsP_tie}
(\path{ExistsEmbFIff.lean}:209--293); the argument here uses its
specified representative, lattice and base-lattice equalities.
See \path{HermitianFormLastCoordinateUniformizerVectorSelf.lean} in
\path{Proofs/IwasawaDecomposition/}, lines 123--151,357--445,467--536,
549--612,623--652,692--702,766--781,835--855,864--942,953--1001,
1025--1032,1038--1080 and 1086--1143.

\paragraph{One common penultimate level on the reference charts.}
Take the chart family $K_\beta$ at $i=m-1$ on the reference tuple
$\mu_0$, with its specified reading and the supplied $\mathrm{IdG}$
for $G$. Its twist set is empty. The index is unique by the final-index
law and $m\ge2$. Use the standard decoration and the auxiliary coset
with centre $0$ to make a point, take its class $c_0$, and put
$J=(K_{\beta_0}).\mathrm{lev}$, $\beta_0=\mathrm{chOf}(c_0)$.
The chosen representative of $c_0$ still has zero in its auxiliary set:
group transport applies a linear map to that set. For any such point
the coupling coset at a marked prime $P$ is precisely
$q_P^{d_P}\mathcal O_{E,P}^n$, where $d_P=0$ at generator places and
$d_P=M$ at the other marked places. Indeed subtract $0$ from a member
of the same auxiliary coset and pass through the inverse local frame;
the difference belongs to this standard sublattice. Conversely use
the representative vector $0$ in the definition of the coupling coset.

Let $J_0$ be the unitary integral group whose reductions at level
primes fix the last row vector $e_n^t$. The fixed right chart factor
$b_\beta$ stabilizes the standard lattice by its standard-point
reading, is the identity at generator places, and has
$e_n^t\bar b_\beta^{-1}=c e_n^t$ at level primes. The scalar $c$ is
a unit because the reduced matrix is invertible. Conjugation by
$b_\beta$ therefore preserves $J_0$ in both directions, and acts
trivially on its generator reductions $r(x)$. At a class $c$, writing
$y_c$ for its chart representative, $y_c x b_\beta$ for $x\in J_0$
reads the old decoration transformed solely by $r(x)^{-1}$ at the
generator primes: conjugate $x$ by $b_\beta$, use that its last row
fixes the level readings, and reduce the inverse matrix on the
generator readings. Conversely, if an element $x$ reads such a pure
generator translate, the two unitary points have equal underlying
lattices. In relative standard coordinates the change of basis is
integral invertible at every prime. Equality of the level functionals
on a local basis forces its reduced last row; equality of generator
coordinates forces its entire generator reduction. Conjugating back
gives $x\in J_0$ and the generator translate $r(x)^{-1}$.
Every unitary matrix yields its underlying lattice and residual
readings by pullback; here there are no twist frames to lift.
The chart membership test at $c$ thus says $x$ belongs to $J_0$,
$r(x)^{-1}$ belongs to the generator index subgroup, and that
translate preserves its coupling coset. Its character is the embedded
generator character on $r(x)^{-1}$.

At $c_0$ the last test is automatic, since pure generator translations
preserve the occurrence of zero. It follows that $J.U$ is exactly the
preimage of the generator subgroup in $J_0$. All chart groups are
contained in $J.U$: each reads the same index and hence the same
generator subgroup. Their characters are restrictions of $J.\chi$,
including its inverse-reduction convention, and their base lattices
coincide because each base reading gives the same coordinate
divisibility by $\pi_T$ in the last coordinate.
Moreover $\ker(r:J.U\to\prod_w\mathrm{GL}_n(k_w))$ is contained
in every $K_\beta.U$, since the resulting generator translate is
the identity on any coupling coset. Its index is finite by finiteness
of the target. The chart labels are a quotient of the finite class
set by the surjective chart reading. Now take
\[
 Y_b=\coprod_\beta J.U/K_\beta.U,\qquad y_b(\beta)=1K_\beta.U
 \quad\hbox{in its $\beta$-summand}.
\]
This is a finite left $J.U$-set; each point belongs to the orbit of
a unique label, and the stabilizer of its marked point is exactly
$K_\beta.U$. These checks prove every orbit and restriction clause of
\path{stepBefore_oneLevel}
(\path{ChartTiedTopLevelChain.lean}:30--167,173--215,229--266,369--484).

It remains to identify the local group and its character in the
specified identity frame. At a generator prime $w$ over $v\mid g_a$,
the standard base lattice has no level modification, because
$v\nmid T$. The split projection prescribed by $G$ reduces at the
selected $w$: of the two conjugate factors, the vanishing of
$\overline{\pi_{G,a}}$ selects this one by $\mathrm{IdG}$. It sends
integral elements to integral elements and identifies their residues
by $e:k_w\simeq k_v$ (the scalar extension on the base field is the
identity). Thus in frame $1$ the local lattice is standard and the
local level condition is entrywise reduction into the generator group.
For a local element satisfying the ambient lattice condition, insert
it at $v$ with identity elsewhere to test the global preimage
description above; restriction to $v$ gives the converse. Distinct
pairs $(a,v)$ enumerate the marked generator primes once, by the two
distinct rational primes and the choice of one of each conjugate pair.

Here is the local matrix test. Put $r=k+2r_S+1$, and for $x<r$ put
$\ell(x)=p(x)$, where $p$ sends $x<k$ to $r_S+x$, sends
$k\le x<k+r_S$ to $x-k$, and fixes the remaining indices.
For $x\ge r$ put $\ell(x)=r+\lfloor(x-r)/t\rfloor$.
The last-index sharp law has zero extra coordinates. Its flags are
coordinate lattices with a single uniformizer on the coordinates
$D_s=\{x:p(x)\ge d_s\}$, where $d_s=s$ except
$d_{k+r_S}=r_S$, and $s<r$ or $s=r+ht$. For $s<r$ the marked
basis vector is $e_{p^{-1}(s)}$. Preservation of the flags forces the
reductions of the matrix and its inverse to vanish in rows of higher
$\ell$ than the column: in the tail choose the block boundary,
and for an initial column use its marked vector, which fixes the
column modulo the flag on $D_s$. The latter also gives diagonal
entries $1$ before $r$, and entry $0$ for $x<k$, $y=k+r_S$ from the
exceptional $d_s=r_S$. Conversely lift a residue matrix and its
inverse by lifting the entries of one matrix to the local ring;
its determinant is a unit. For $x\in D_s$, $y\notin D_s$ one has
$\ell(y)<\ell(x)$ (on the tail the cut is a block boundary), so
the cross entries of both lifts are divisible by the uniformizer.
Multiplying a vector shows preservation of each flag in both directions.
For the marked-vector congruence the requisite components of the
column are its diagonal $1$ and zeros strictly below in this order;
at the exceptional cut the additional rows $x<k$ vanish by the
extra column condition. This constructs the integral lift in the
sharp stabilizer, with no extension of coordinates at $i=m-1$.
Thus the reduced group is exactly $\mathrm{JSymG}(k_v,n,k,r_S,t,
\mathrm{false})$. See \path{SymmetricLevelBoundLe.lean}:30--91,95--129
and \path{SymmetricFlagBasisVectorEqSingleOfNoExtra.lean}:37--88,
152--221,241--275.

Define $\lambda_J$ on this group by transporting the generator
character through $e^{-1}$ and the fixed coefficient embedding. In the
sharp integral model the extra-coordinate additive term is empty.
Applying the reduced matrix to each standard basis vector identifies
its columns with the reductions of the lift. Its remaining additive
sum is exactly $\sum_{x+1=y<r,\ y\ne k} M_{xy}$: the ranges
$0\le x<k-1$ and $k\le x<r-1$ partition it. The multiplicative
terms are the determinants of the consecutive $t$-blocks beginning
at $r+ht$, $0\le h<i$. Reduction commutes with determinants, and
each residue is a natural-number class since $|k_w|=g_a$ is prime.
Thus the $\mathrm{IdG}$ identities for the additive and
multiplicative characters give precisely
$\lambda_J(M)=\mathrm{lamSym}(G.\psi_0,G.\chi_B,k,r_S,t,i,M)$.
For a global element, reduce at each $w$ and transport by $e$ to
define the reductions for $V$. Their inverse characters multiply to
$J.\chi$ by the inverse-reduction reading; reindexing by $(a,v)$
preserves the product. Each factor has norm one since it is the
image of an element of a finite group. This gives the two remaining
character and normalization fields of $\mathrm{GLevH}$ and
$V.\mathrm{fr}=1$ at all generator places
(\path{ChrLast.lean}:20--50,56--117,123--160;
\path{GloRhoZeroOfCompat.lean}:49--135;
\path{ChartTiedTopLevelChain.lean}:278--363). This proves
\path{GlobalTheory.Mem.Col2.OneLevel.stepBefore_oneLevel} for
the charts in question.

\paragraph{The reference generator places.}
Fix the reading at $\mu_0$. Its twist-prime set is empty. A place $v\mid g_j$
of $F$ has a selected place $w$ of $E$, the one cut out by
$\overline{\pi_{g_j}}$ in the splitting defining $G.\mathrm{phi}$.
Indeed the kernel on $\mathcal O_E$ of reduction after this projection
is a prime above $v$, and the splitting identity for $\pi_{g_j}$ and
$\mathrm{IdG}.\mathrm{piG}$ put $\overline{\pi_{g_j}}$ in it.
The split factors are distinguished by this divisor. Distinct $g_j$ are
coprime, so the selected $w$ reindex the pairs $(j,v)$ without repetitions.
The resulting residue comparison respects integral reductions. Every element
of $\mathcal O_E/w$ is a natural-number residue, since its order is the
prime $g_j$. Thus the character comparisons in $\mathrm{IdG}$ on natural
numbers give comparisons on all residues, and on units for the multiplicative
characters. At these places $T=T(u)$ is a unit and
$\pi_T\overline{\pi_T}=T$. The base lattice is integral and contains
$Te_a$ for every coordinate $a$, so its localisation is the whole
integral lattice. Under the first split projection its image is exactly
$\mathcal O_{F_v}^n$: lift an integral vector $z$ by $(z,0)$ in split
coordinates. The reference Gram matrix is the identity in the coordinates
of \path{spaces_in_coordinates}.
See \path{LevelTwistPrimesStepZeroCharacterEqOne.lean:158--208},
\path{LevelsTopRowSym.lean:91--115} and \path{ValLtOfRedM.lean:245--282}.

\paragraph{Stage zero.}
The assertion is \path{glev_zero}.
At this stage $N(0)=k$, the twist character is an empty product and the
extra coset restriction is empty since $m\ge2$. For $x_v$ preserving
the local Hermitian lattice, put $x=x_v$ at $v$ and $1$ elsewhere.
The chart translator $k_0$ is $1$ at $v$. The matrix
$A=k_0^{-1}x^{-1}k_0$ and its inverse are integral, are the identity
above $T$, and $y x k_0=y k_0 A^{-1}$. Its residues act on the generator
frames; all but the one above $v$ are identity. The level reading tests
precisely membership of these residues in the stage-zero column group:
a witnessing rational factor acts by the identity matrix and hence is $1$;
two remaining residue factors acting in the same way on a basis coincide.
Conversely the specified residue tuple together with identity rational and
twist factors is a witness. Here the group is upper unitriangular:
each of the $k$ diagonal blocks is a singleton with entry one and there
are no further blocks. This condition is preserved by passage to the
inverse and detected by the injective residue comparison. It is exactly
\path{InJ} with parameters $(k,t,N(0),\mathrm{false})$, in frame $1$.

For arbitrary $x\in U$ the reading of its character gives the product
of the additive characters evaluated on the selected superdiagonal sum
of the inverse residues (the generator components of $k_0$ are identity).
For upper unitriangular matrices every first-superdiagonal entry of a
product is the sum of the two entries. Consequently inversion negates
the selected sum, so each factor is the inverse of $\lambda_J$ on the
residue of $x_v$. There are zero determinant blocks at this stage.
Reindex by $(j,v)$ as above. The values have norm one because the
additive groups of the residue fields are finite. Together with the
local lattice calculation this proves every field of \path{GLev} with
all frames $1$. See \path{LevelTwistPrimesStepZeroCharacterEqOne.lean:217--350},
\path{LevelTwistPrimesStepZeroCharacterEqOne.lean:352--386} and
\path{LevelTwistPrimesStepZeroCharacterEqOne.lean:482--573}.

\paragraph{Full and primed top rows.}
The assertion is \path{levels_top_row_sym}.
Choose one class on each side by surjectivity of its class bijection.
At stage $m$ the full chart has empty coset condition; on the primed side
the translator is $1$. For a unitary matrix supported at the generator
places and integral with integral inverse, multiplication on the right
changes only the generator frames, by its inverse residues. As in stage
zero, freeness on bases gives membership exactly when each residue lies
in the prescribed subgroup; the character is the product of its inverse
local characters. This uses the top column reading on the full side and
the row reading on the primed side. For arbitrary $x\in U$ discard its
components away from the generator places: the truncation belongs to $U$
by its local membership criterion and its quotient with $x$ has character
one by the defining triviality in a step level. An element supported at
one $v$ similarly gives the local membership equivalence.

To identify the two subgroups put $r=k+2r_S+1$, $t>0$,
$n'=r+mt$, $n=n'+1$. Below $r$, the ordering of scalar levels moves the
$k$ initial positions past the next $r_S$ positions, leaving the rest
in order. Under this order, entries below the levels vanish, all the
diagonal entries are $1$, and entries $(a,k+r_S)$ with $a<k$ vanish.
The remaining coarse blocks $s<m$ are precisely
$[r+st,r+(s+1)t)$ and are block upper triangular. These are the row
subgroup conditions in rank $n'$, read entry by entry. In rank $n$ the
extra terminal column is the last basis vector; its row is also the
last basis covector since its coarse block follows all the others.
Conversely these entry conditions recover the top parabolic, its
leading-block conditions, and the fixed final column. Applied also to
inverses, this identifies both subgroups with \path{HSym}, with flags
$\mathrm{true}$ and $\mathrm{false}$ respectively. In the leading
block the additive sum takes entries $(a,a+1)$ with $a+1<r$ and
$a+1\ne k$. In each interval above the determinant factor is the
determinant of that $t$-by-$t$ diagonal block. Residue comparison
preserves finite sums and determinants, so the two transported local
characters have exactly the prescribed \path{lamSym} values. They
are homomorphisms by transport of the original characters, and their
finite domains give norm one. See \path{MemHJsymIffJSymG.lean:12--16},
\path{CoarseLast.lean:37--104} and \path{LamJsymRead.lean:55--88}.

Take the primed frame to be $1$ and the full frame to be
$\operatorname{diag}(1,\ldots,1,\phi_v(\pi_T))$. The last scalar
is a local unit. The diagonal frame preserves the integral lattice and
its reduction centralises the full subgroup by the terminal row and
column just calculated. Thus the local membership and character formulas
also hold in these frames, with the reduction map on $U$ given by its
framed components. The standard inclusion of the primed space appends
zero, so the first $n'$ columns agree under inclusion while the last is
the split image of $\pi_T e_{n-1}$. These are the two requirements of
\path{FrExtFr}. See \path{IsPtOfOfDecoRaw.lean:68--116},
\path{LevelsTopRowSym.lean:52--89} and \path{LevelsTopRowSym.lean:117--224}.

\paragraph{Arbitrary chart at the reference tuple.}
The assertion is \path{chart_gLevH}.
For $i\le m$ and a chart label $\beta$, choose $c$ above $\beta$ by
surjectivity in the reading. For $v\mid g_j$ the chart translator is
$1$ at $v$. In the standard frame, selected split reduction defines a
homomorphism $r_v$ on the unitary lattice stabiliser (integrality of a
matrix and its inverse gives invertibility after reduction). If
$r_v(a)=1$ then both split factors reduce to identity: with Gram matrix
$1$ the other factor is its inverse transpose. Insert $a$ at $v$ and
identity elsewhere to obtain $\widetilde a$. This commutes with the
translator, preserves the integral adelic lattice in both directions,
and fixes the represented decorated point. Indeed on each integral
pullback vector the difference after applying $\widetilde a^{-1}$ is
a sum of entries congruent to zero times integral coordinates; its
residue is zero at every place. Thus both the functionals and generator
bases are unchanged, and there are no twist-prime decorations. The
identity witness in the level and character readings gives
$\widetilde a\in U$ and $\chi(\widetilde a)=1$.
In particular $a\in U_v$ and the local character is trivial on this
kernel. Put
\[
 H_{j,v}=r_v(U_v),\qquad
 \lambda_{j,v}(r_v(a))=\chi(\widetilde a)^{-1}\quad(a\in U_v).
\]
The kernel calculation makes the formula well defined and multiplicative.
A lattice stabiliser whose residue is in $H_{j,v}$ differs from a lift
in $U_v$ by an element of the kernel and thus belongs to $U_v$. For a
global $x\in U$ take the reductions of its components. Its quotient by
the product of the finitely many isolated generator components has all
bad components $1$, hence character one. This gives the required product
of inverse $\lambda_{j,v}$ and norm one since each image is finite.
The identity frame gives the asserted lattice by the first paragraph.
This proves existence of the chart groups and characters for each label.
See \path{ValLtOfRedM.lean:102--124}, \path{ValLtOfRedM.lean:133--241},
\path{ValLtOfRedM.lean:284--405} and \path{ValLtOfRedM.lean:431--488}.

\paragraph{The common level-prime datum and simultaneous assembly.}
For $A_u$ take the actual residue maps $\mathcal O_E\to k(P)$,
$\pi_u=X.\pi_T(u)$ and contraction from $P\in S_T$ to $v\mid T$.
Every $v\mid T$ has a prime of $E$ above it; $T=\pi_u\bar\pi_u$
places that prime or its conjugate in $S_T$. Above the same $v$ two
primes are equal or conjugate, and coprimality excludes two conjugates
in $S_T$. Thus contraction is bijective. Each residue map is onto,
with kernel $P$, and $[\bar\pi_u]_P$ is a unit. At a coefficient prime
$w\mid T$, Galois transitivity on primes of $E$ over $T$ and the CM
representatives for embeddings into $K$ give, for each $P$, an embedding
index $\tau$ whose pullback of $w$ is $P$ or $\bar P$. The index cannot
serve two different $P$'s; equal cardinalities give the required
bijection. In either case the pullback on the fixed field is $v$,
so the valuation test on $F$ is the required residue test. The images
of $\pi_u$ in $K$ under the chosen embeddings agree by the prime datum.
Thus if $w$ does not see all those images, it sees none and the
product $\pi_u\bar\pi_u=T$ forces all the conjugate images into $w$.
The congruence and $K$-value of $\pi_u$ are the original data.
See \path{UnderTBijective.lean}:18--46,61--98,
\path{HermitianCodeWordWeilElementMemQuadraticField.lean}:110--153,
\path{LiesOverSpanOfNatCastMemHeightOne.lean}:17--58.

For $0\le i\le m$ choose the constructed chart families and put the
unique endpoint chart into the endpoint records. The selected fibre
representative at every stage is exactly $I.yStd$, since these were
chosen in the proof of the chart. For $i>m$ take the singleton label;
the chart fields have empty range of stage hypotheses there. Take the
single $G$ already chosen, every $A_u$ above, and the frame $e_Y$ of
supplier~6 throughout. In the determinant comparison use $C=1$:
its Gram matrix is $1$, its integral lattice is the given one in
coordinates and its determinant contributes the factor $1$. Use the
reference primed chart just constructed simultaneously for its class,
value, norm and reading fields; its identity equivalence is the one used
in the tilted evaluation. For each chart retain the projections and
residue-value identifications for those representatives from the direct
local calculations of supplier~8. This supplies both moves, the localized
basis, their endpoint restrictions and the mid-stage projection fields.
The ordinary class, form, norm and diamond data, the common penultimate
level and its transitive finite sets, and the bottom, top and primed
residue levels use exactly these charts and $G$ as in the preceding
comparisons. Thus the generator-level restrictions in the mid-stage
records and the character formulas at both endpoints hold on the same
$F_u$. For the tilted functional the common last-coordinate comparison
gives $[\bar\pi_u x_n]$ because the distinguished vector is
$\pi_u e_n$ and the Gram matrix is $1$. The extension identity for
the generator frames uses this same vector. Use the conjugate-linear
period embedding constructed in \path{ch_P_r11.tex}, lines 961--1001,
for the image, surjectivity onto the row point space and period clauses.
These choices give all fields in \path{FormsTiedSym}; the place and
frame identifications, equality to $yStd$, the two projection tests,
the primed reading and the tilted evaluation are precisely
$\mathrm{StdF}$. The mid-stage projection tests and the identity generator frames for the penultimate level give $\mathrm{StdLvC}$.
Every word uses $(F_u).\mathrm{gl}=G$, so these are equal literally,
and at the reference word $G$ has the character-compatibility certificate
proved above. Together with finiteness at all stages and the singleton
top label this is the full implication
\path{GlobalTheory.T.exists_standardFormComparison_symmetricLevel_members}
(\path{ExistsStandardFormComparisonSymmetricLevelMember.lean}:52--243).

% blueprint/ch_P_r12.tex lines 3172 through 4285; r12:s7, r13:interior-charts, r13:stage-pack
\subsection{Supplier 7. The twisted chain over the chosen columns}
\label{r12:s7}
Fix $I$ with its correctness and pin data, the laws $l,l_0$, later data
$Z$, and the chosen columns $H$ and forms $F$ with
$\mathrm{CarH}$, $\mathrm{StdF}$, $\mathrm{StdLvC}$. The prescribed
predicate at the base primes must agree with $I.\mathrm{UnderQ}$.
The output is one triple $T,J,W$ satisfying \path{TwChMu0} $T$,
$I.\mathrm{StdW}(W)$ and $\mathrm{Nonempty}(\mathrm{SymmetricTwistedJointForms}(T,F))$.
The first assertion identifies the frame comparison at each reference
stage with the canonical rank cast (\path{eY}) and the bottom class
comparison with the canonical class cast (\path{eCls0}). The nine fields of
\path{StdW} are the bottom class and frame identifications
$\mathrm{eCls0},\mathrm{eY}$, the lattice $\mathrm{lat}$,
the place and residue comparisons $\mathrm{eVT},\mathrm{eFv}$,
the chart readings $\mathrm{reads}$, the top-level test $\mathrm{top}$,
the common top frame $\mathrm{topk}$, and the vanishing of the
uniformizer on the chart functional $\mathrm{killed}$.
The joint forms witness uses the same $T,F$ for the local generator
tests and the two projections of the finite coupling set. See
\path{TiedQCJoint.lean}:34--36, \path{StdLvPr.lean}:100--120 and
\path{TiedTwJSym.lean}:25--45. For a nonreference tuple, every
word and $1\le m$,
one family of top base points both reads the data of $W$ and sends every
allowed block extension of a primed point to its tilted point. Projecting
the latter assertion from $\mathrm{StdW}.\mathrm{topk}$ is valid only for
the $W$ just constructed. See \path{tw_ties_with_points}
(\path{TrueColumnChoiceExistsChainDataWithBasePoints.lean:13--23}) and
\path{trueColumnChoice_exists_chainData_with_basePoints}
(\path{TrueColumnChoiceExistsChainDataWithBasePointsStatement.lean:10--18}).
In the following, $Q_E$ denotes selected primes of $E$ and $Q_F$ their
contractions to $F$. In the application the prime-separation theorem
\path{twistPlaces_prime_to_levelAndTypePrimes}
(\path{TwistPlacesPrimeToLevelAndTypePrimesStatement.lean}:7--8)
excludes $T$ and both generator primes over every $v\in Q_F$.
The choice of one split prime over a twist place excludes conjugate
pairs (\path{good_of_mem}, \path{levOk_twists},
\path{LevOkTwists.lean}:21--29,48--70). At the reference tuple,
\path{not_twistPlace_referenceTuple} and
\path{ColIn.twistPrimes_reference_tuple_eq_empty} give $Q_F=Q_E=\varnothing$
(\path{NotTwistPlaceReferenceTupleStatement.lean}:7--8;
\path{TrueColumnChoicesHeckeCharReferenceEqOne.lean}:273--279).

\paragraph{Constructing the diagonal at the selected primes.}
Let $E/F$ be the specified CM extension, $n>0$, and assume that $Q_E$
is disjoint from its conjugate, and that both avoid $S_T,S_G$. Let
$\pi\in\mathcal O_E$ be a unit at every $P\in Q_E$. There exists
$t\in K_1$ with $t_P=\operatorname{diag}(1,\ldots,1,\pi)$ on $Q_E$
and $t_w=1$ off $Q_E\cup\bar Q_E$. For empty $Q_E$ take $t=1$, including
when $\pi=0$. Otherwise $\pi\ne0$. For the last diagonal entry take the
idele $s$ with components $\pi$ at $P\in Q_E$, $\bar\pi^{-1}$ at
$\bar P$, and $1$ elsewhere. One can write $s$ in the specified coefficient
algebra as
\[
 1+e(\pi-1)+\bar e(\bar\pi^{-1}-1),
 \qquad e=\sum_{P\in Q_E} e_P .
\]
Here $e_P$ has component one at $P$ and zero elsewhere: subtract one from
the idele with component two at $P$ and one elsewhere. The summands are
orthogonal idempotents and $e\bar e=0$, by componentwise evaluation.
It follows that $s\bar s=1$. Every component of $s$ and its inverse
is a unit in the appropriate local integer ring, so $t$ and $t^{-1}$
preserve the standard integral lattice. At $S_T$ the identity components
satisfy the mirabolic condition. At $S_G$ the reduction is the identity.
These are precisely the additional conditions defining $K_1$ inside the
unitary group. This proves \path{exists_unitaryAdelic_twist_diagonal_at_primes}
(\path{ExistsUnitaryAdelicTwistDiagonalAtPrimesStatement.lean:8--9},
\path{TwUnit.exists_tw}, \path{ExistsUnitaryAdelicTwistDiagonalAtPrimes.lean:31--113}).
The component tests, conjugation of idempotents and the level definition
are \path{ad_ext}, \path{adAt_idem}, \path{adAt_sig_idem}, \path{adAt_glue}
(\path{ValuedLtOfMem.lean:18--35,71--101}) and \path{J0}, \path{redG},
\path{K1} (\path{Formalism/LocalToGlobal/Columns/Adelic.lean:143--158}).

\paragraph{The top chart and its tilted block.}
Use $k_0$ from the reference reading, and set $k=k_0t$. Its prescribed
last row at $S_T$ is unchanged. For each class choose an adelic
representative $g_c$ realizing the complete decorated lattice and put
$y_c=g_c k^{-1}$. Thus one common $k$ works for the representative
equations. In the identification of the action with triples
$(\gamma,h,\theta)$ the stabilizer image consists of the $\gamma$ for
which some $h,\theta$ fix that representative with $\gamma$. Projection
is injective by the stabilizer axiom of the supplied step. At the top,
the index is unique and the coupled coset is empty, so the chart label
can be a singleton. A block extension fixing the distinguished last
vector has the form $\operatorname{diag}(g',1)$: its first block is
prescribed by the inclusion, and the fixed vector is a nonzero scalar
times the last basis vector. At $P\in Q_E$, $k_0$ is the identity, and
multiplying by $t$ replaces the last diagonal entry by $\pi$. At all such $P$ the upper left
block is the rational frame of the given primed point $d$. The two
off-diagonal blocks vanish, so the frame is exactly
$\operatorname{diag}(\mathrm{fr}_P(d),\pi)$. More generally this
calculation works for any $t\in K_1$ with the prescribed components,
and any reference $k_0$ with the three hypotheses of
\path{GlobalTheory.T.pointOverX_has_adelicPoint_blockExtend_mul}.
Forgetting the twist frames of $d$ gives the following other three
point equations. The direct sum block carries the lattice to
$N(d)\oplus\mathcal O_E$; both $k_0$ and $t$ preserve the standard
lattice. The last row of $k_0^{-1}$ at the level primes reads
$\bar\pi$ times the last coordinate, and the inverse of $t$ has the
mirabolic last row there. The generator vectors read those of $d$ on
the first block; at each generator $w$, $\pi\equiv1\pmod w$ by
\path{piT_cong}, $k_0$ is one since $T\ne g_j$, and $t$ reduces to
one. This yields the generator equation for the tilted last coordinate.
All factors are unitary. These are the lattice, functional, reduction
and frame fields of the assertion for every extension $g$ of $g'$.
See \path{PointOverXHasAdelicPointBlockExtendMul.lean:73--206} and the
untwisted direct-sum computation in the reference tilted endpoint proof.

For the top chart itself, first recover the given old level in the
$k_0$ frame. At the reference tuple the extra frame set and coupled
cosets are empty by \path{ColIn.twistPrimes_reference_tuple_eq_empty}
(\path{TrueColumnChoicesHeckeCharReferenceEqOne.lean}:273--279).
Any new unitary adele over a fixed representative
therefore has a represented unframed decoration, from its lattice
and pullback reductions. Applying the given level reading to it and
to an element $x$ gives precisely $k_0^{-1}xk_0\in J_0$ with inverse
generator reduction in the top type group. Comparing characters on
the same move gives the old residue character. The base lattice
condition follows by canceling the common representative in its two
completed-span equalities. These are the three facts used from
\path{kread_member} (\path{TrueColumnMemberLevelReadAtBasePoint.lean}:34--119).
The top twist
congruence condition and its character are trivial on the next deeper
principal subgroup by the definition of the top level and its character
(\path{LastStepTwistLevelContainsPrincipalCongruence.lean:47--54}).
The local character thus factors through a finite quotient of the
integral matrix group.

The shared decorated-lattice construction supplies $g_c$ on the full
twist set via \path{exists_adelicPoint_decoratedLattice_allTwists}.
Apply \hyperref[r13:s9-twopoints]{the two-point level test} to $k_0,t$ and these
type and twist groups. Its \path{mem_framedLevel_iff_two_adelic_points}
identifies pure residue moves with membership in the cut level.
The proof of \hyperref[r13:s9-congruence]{prescribed congruences}, \path{exists_K1_point_congruent},
realizes a point at every prescribed finite depth, as required by
\path{topU_dense}. The proof of \hyperref[r13:s9-cosets]{decorated classes and double cosets},
\path{decoratedPointClasses_bijective_doubleCosets}, then identifies each
represented orbit with one double coset
(\path{MemFramedLevelIffTwoAdelicPoints.lean}:40--124;
\path{ExistsK1PointCongruent.lean}:11--52;
\path{DecoratedPointClassesBijectiveDoubleCosets.lean}:34--83;
\path{ExistsFramedChartFromAdelicDictionary.lean}:63--92,98--158). Its orbit classes
are the classes at the unique top index and its rational stabilizers
are the triple stabilizers projected above. Restrict the old level
at the marked primes by the conjugated twist conditions, and multiply
the old character by the inverse of the extra local character product.
The two-point equivalence above gives the level reading on this same family.
For a move $(1,h,\theta)$ its twist transition has image
$\theta_P^{-1}$, so the inverse product contributes
$\sigma_0(\Theta.\mathrm{chr}(\theta))$ and gives the character
reading too. The character has finite exponent: its old part
factors through the finite type reductions and the extra local
factors have finite quotients. It agrees with the old character on
elements trivial at $Q_F$, and the local levels agree off $Q_F$.
If $y_c(k_0t)$ and $g(k_0t)$ both represent the prescribed tilted
decoration, the same two-point test with the identity residue pair
gives $g^{-1}y_c$ in the chart level. This proves all the assertions of
\path{GlobalTheory.T.chart_lastStep_twistBasePoint}, including the
common $k_0t$ and the extra base-point test on the chosen $y_c$.
See \path{ChartLastStepTwistBasePoint.lean:124--184,195--336}.

\paragraph{A supported lift of all the type-prime reductions.}
For $I,l,u$ as above, any $n$ and arbitrary
$x_w\in\mathrm{GL}_n(\mathcal O_E/w)$, one for every type prime
$w\in S_G$, we can supply the matrix lift required in the interior
matching argument. A selected $w$ divides $\bar\pi_{G,j}$ for some $j$.
Its conjugate cannot also be selected: for the same $j$ this would
put both coprime elements $\pi_{G,j},\bar\pi_{G,j}$ in $w$; for
different indices it would put the two distinct rational primes
$g_0,g_1$ in $w$, by their norm identities. A prime containing the
level prime $T$ cannot lie under one of these places since $T$ is a
different rational prime. These are precisely the disjointness
conditions used in \path{exists_unitary_lift}
(\path{ExistsUnitaryAdelicLiftTypePrimeReductions.lean:179--248}).

Choose integral entrywise lifts $A_w$ of the matrices $x_w$.
Their determinants have nonzero residue, hence their inverses have
$w$-integral entries by the adjugate formula. Take the component $A_w$
at $w$, the conjugate inverse transpose at $\bar w$, and the identity
at every other prime. Entrywise the finite sum with the idempotents
$e_w,e_{\bar w}$ constructed above is an adelic matrix, and the same
construction with the inverse matrices gives its inverse. Conjugate
transpose interchanges the two components of each pair and yields
the inverse at each place. Both matrices are integral everywhere, so
the resulting unitary element preserves the standard lattice; with
all decoration sets empty this is exactly a representative of the
standard point. Outside the places over the type primes every component
is the identity, including all those over $T$. Projection onto the
factor at each selected $w$ reduces to $x_w$ as prescribed. For rank
zero all matrices are the unique empty identity matrix and the same
claims hold. This gives
\path{exists_unitaryAdelic_lift_typePrime_reductions}
(\path{ExistsUnitaryAdelicLiftTypePrimeReductionsStatement.lean:7--13};
\path{ExistsUnitaryAdelicLiftTypePrimeReductions.lean:101--121,142--155}).

\paragraph{Matching the coupled coset with the reference tuple.}
Fix $1\le i<m$, a tuple class $c$ and a prescribed reference index $b_0$.
If $i+1\ne m$, take the class at $b_0$ of the base decoration with its
empty coset and set $y=1$. Every move preserves the empty framed cosets.
Suppose $i+1=m$, and write $r$ for the preferred representative of $c$.
Choose a centre $z$ of its defining coset. At each generator prime $w$,
$q_wz\in N_w$, where $q_w=|\mathcal O_E/w|$ and $N_w$ is the
localized lattice. Forget the twist frames of $r$. For this reference
decoration prescribe the coset of $z$ modulo the vectors integral in
$N_w$ at every generator prime. This is admissible by the denominator
conditions just given, since the reference twist set is empty.
Call the resulting point $p_0$. Its reduced sets $B_w$ are precisely those
of $r$: both are the singleton consisting of the reduction of $q_w z$ in
the same generator frame. This explicitly realizes the matching fibre
for a twisted source, using the reference-fibre calculation of supplier 1.

Let $c_0$ be the class of $p_0$ at $b_0$. Choose a move $s$ sending its
preferred representative $r_0$ to $p_0$, and take $y=s_G^{-1}$. This
lies in the type group at $b_0$. Reduction intertwines the action with
$s_G$, since moving a vector and its rational frame together cancels
the rational factor. Thus the reduced set of $r_0$ at $w$ is $y_w B_w$.
At a generator prime, two framed cosets agree exactly when their
reduced sets agree: in standard coordinates the coset is
$u+\mathcal O_{E,w}^n$, with $q_w u$ integral. The source prime law
$q_w\notin w^2$ makes $q_w$ a local uniformizer. Equality after reducing
$q_w u$ is then equivalent to integral difference of the centres.
All remaining coordinates of the framed coset are the twist coordinates.
A move $(1,h,1)$ leaves each of them fixed: the frame and its vectors
undergo the same rational change and the twist factor is the identity.
The reference has no such coordinates. Therefore coset preservation is
equivalent to $h_w B_w=B_w$ at all generator primes. Under the required
conjugacy $h_{0,w}=y_wh_wy_w^{-1}$ this is equivalent to preservation
of $y_w B_w$ by $h_{0,w}$, giving both implications of
\path{GlobalTheory.T.exists_reference_class_matching_coupling_coset}
(\path{ExistsReferenceClassMatchingCouplingCoset.lean:41--139},
10,738 core lines). The centre, its reduced singleton, and the
twist-coordinate action are explicit in
\path{CouplingCosetsInFramesEqIffReductions.lean:19--75,77--130}
and \path{CouplingCosetMovesWithFrameChange.lean:14--45}.

\paragraph{Matching the chart index at generator primes.}
For $1\le i<m$, every tuple index $b$ has a reference index $b_0$
and residue conjugators $x_w$ identifying the generator-level subgroups
and characters; $x_w=1$ when $i+1=m$. This is
\path{GlobalTheory.T.exists_conj_referenceChartIndex}
(\path{ExistsConjReferenceChartIndexStatement.lean}:7--16).
The source first proves the identity-double-coset assertion
\path{cosetOne} via \path{twSurjG} and existence \path{idx_nonempty}
via \path{twNonVanG} (\path{ExistsConjReferenceChartIndex.lean}:34--79,96--130;
\path{ChainTwOk.lean}:18--34).
The two identity-coset and existence declarations jointly dominate
5,045 core lines. The proof here replaces that classification with a
comparison for the given index. The alternative \path{Idx} constructor
$\mathrm{none}$ is impossible at $i<m$: its membership condition
would rule out $i<m$ (\path{idx_match},
\path{ExistsConjReferenceChartIndex.lean}:97--106).
Set $m_2=m-2$ and write $h_{b,w}$ for its local representative,
$J_j$ for the sharp local level, and $R_{1,j}$ for the group preserving
the prescribed source flag and acting trivially on its successive marked
directions. A source double coset means $R_{1,j}h_{b,w}J_j$.
The claim we induct on is that one reference index has the same
double coset at every generator prime. It is enough to prove this
comparison; we do not require a classification of all nonzero sums.

\paragraph{Lemma. Tuple independence of the sharp generator sum.}\phantomsection\label{r13:tuple-independence}
Fix a generator prime $P$, the common ranks and residue characters,
a stage $j<m_2$, and matrices $a,h$ stabilizing the source and target
lattices. Compare the given tuple and the reference tuple under their
index laws, in particular \path{ChIdx.Ok.psi0}; both sharp flags are
true by \path{ColIn.Ok.sharp}
(\path{TrueColumnChoicesHeckeCharReferenceEqOne.lean}:242).
For every \emph{same} pair of depths $A',B'$ in both evaluations
and every $X$, the two unfolded sums divided by their respective
Weil prefactors agree. Consequently the condition $V_j(a,h)$ that
the unfolded function be nonzero for every sufficiently large pair
of depths is tuple-independent (the witness $X$ may depend on the
depths). This is a hand lemma comparing two tuples.

Indeed \path{PhiAt} selects precisely $\Phi_{GS}$ at $P$ on both
sides. In \path{Loc.om} its argument is $h^{-1}Ya$. The determinant
of each stabilizer has $P$-valuation zero, since its integral-basis
matrix and inverse are integral. Thus the exponents of the unramified
factors $\mathrm{muq}$ and the square-root factors $\mathrm{sq}$ vanish.
The residual unit-character factors are nonzero and independent of $Y$.
See \path{omPhi}, \path{E1at}
(\path{FiniteSumNonzeroLargeDepth.lean}:927--944) and
\path{Loc.om}, \path{Loc.E1}
(\path{Formalism/PeriodIntegrals/TwistedEisensteinNonvanishing.lean}:181--184,202--207).
On its support $Y$ maps the scaled source lattice into the scaled
target lattice, since $a,h$ preserve them. The phase is the negative
of the sum of the entries in the adjacent marked coordinates. For
each such entry the source exponent exceeds the target exponent by
one, so $q\,\mathrm{chArg}(Y)$ is integral. The index law
\path{psi0} thus makes both additive characters equal there to
$\psi_0(\mathrm{red}(-q\,\mathrm{chArg}(Y)))$. Off support both
terms vanish. The finite box and its normalized weight at the common
$A',B'$ coincide, so the two sums differ only by the nonzero scalar
ratio just identified. In particular, dividing each sum by its own fixed nonzero prefactor gives the same sum on both tuples for every $X,A',B'$. Consequently $V_j(a,h)$ is tuple-independent
for these fixed local matrices. This compares each depth, including
the quantifier that the test argument may depend on $A,B$.
See \path{FiniteSumNonzeroLargeDepth.lean}:925--945,988--996,1012--1017;
\path{UnfoldedDatumNeZeroIffDoubleCoset.lean}:124--159,185--238.

\paragraph{Passing from the sum to a chart index.} For fixed $h$ the truth of $V_j(a,h)$ is unchanged upon
replacing the source by $a'=r a\kappa$, with $r\in R_{1,j}$,
$\kappa\in J_j$ and $a,a',h$ lattice stabilizers.
Indeed $r=a'\kappa^{-1}a^{-1}$ stabilizes the source lattice. In
the affine unfolding matrices $Y=\mathrm{xx1}(X,t)$, composition
$Y\mapsto Yr$ preserves the flag and the prescribed unit diagonals
and gives the new initial block $Xg$, where $g$ is the restriction
of $r$ to that block. Its inverse is composition with $r^{-1}$.
In the sum defining the phase, the next-to-diagonal entries of $Yr$
sum to $\mathrm{chArg}(Y)+\mathrm{n1Arg}(r)$: the marked
diagonals are one and the entries below each allowed column vanish.
The right $J_j$ action on the kernel contributes its row character
divided by the determinant character. In the defining residue kernel
this is the source transformation law: on the allowed blocks the
leading determinants multiply, the additive contribution translates
by its value on $\kappa$, and the inverse action preserves the
support. All multipliers are nonzero. To justify changing the finite
sum, let $H$ be the lattice of maps carrying the source lattice into
the target and let $M$ be the differences of free parameters whose
linear matrix lies in $qH$. The integrand is constant on $M$: the
kernel reads reduction modulo $qH$, and the added phase is integral.
Composition by $r$ preserves $H$ in both directions and its induced
affine bijection identifies the cosets modulo $M$. For a common bound
$d$ clearing the denominators of $H$ and its reverse inclusion, the
support lies in $q^{-A}\mathcal O_P^{\rm free}$ for $A\ge d$, and
$q^B\mathcal O_P^{\rm free}\subset M$ for $B\ge d+1$. The finite
box is a complete system of representatives modulo this finer lattice;
each contributing $M$-class is counted the same number of times.
This proves the change of variable even when the map does not
preserve the particular chosen representatives. The inverse change
gives the converse. See \path{ExistsMoveRight.lean}:18--50,
\path{PhiMapsAndOmegaLaw.lean}:819--836, and
\path{unfold_change₂}, \path{unfoldedDatum_ne_zero_of_doubleCoset_source}
(\path{UnfoldedDatumNeZeroIffDoubleCoset.lean}:20--118,171--296).

At stage zero both sets of representatives are $1$ by
\path{zero_uniq} and \path{zero_h}. Given a tuple index $b$ at
$j+1$ with $j<m_2$, for each reference marked prime choose
\emph{its own} predecessor $a_P$ for $b$ from the tuple's
\path{nonvan} law. Every reference marked prime is a generator
because $Q_{\mu_0}=\varnothing$. By induction $a_P$ has a reference
predecessor matching its source double coset at $P$. The two
invariances give the reference nonvanishing test with this reference
predecessor and the target matrix $h_{b,P}$. Apply the reference
\path{surj} once to the family $(h_{b,P})_P$: all targets stabilize
the required lattices by \path{stabG}, and its nongenerator clause
has empty domain. This gives one $b_0$ matching at all $P$.
At $j+1=m_2+1$, choose the unique last reference index. Both
parameters equal one there by \path{last_h}. This proves the
induction using only the nonvanishing already supplied for $b$.
See \path{FiniteSumNonzeroLargeDepth.lean}:1018--1039; the source
implementation uses the stronger identity-coset argument at
\path{ExistsConjReferenceChartIndex.lean}:38--81,87--135.

At a generator prime write $h_b=r h_{b_0}\kappa$ with
$r\in R_{1,i}$ and $\kappa\in J_i$. All three of $h_b,h_{b_0}$
and $\kappa$ stabilize the local lattice, hence so does $r$.
Its restriction $g$ to the initial coordinate space preserves the
standard sublattice there, as does $g^{-1}$. Set
$x_w=\mathrm{red}_w(g^{-1})$. Conjugation by $r^{-1}$ identifies
$R_{1,i}\cap h_bJ_i h_b^{-1}$ with
$R_{1,i}\cap h_{b_0}J_i h_{b_0}^{-1}$: in the second condition
$\kappa$ conjugates within $J_i$. Restrict to the initial coordinates
and reduce; the defining group law \path{grpG} gives the desired
two inclusions, using the inverse conjugation for the reverse one.
For a lift $s$ in this intersection, \path{chrG} writes the
character as the row character at $h_b^{-1}s h_b$ divided by the
phase at $\mathrm{n1Arg}(s)$. The row character on $J_i$ is
multiplicative and tuple-independent at $w$, hence invariant under
$\kappa$-conjugation. The flag condition implies that the adjacent
marked entries add under multiplication in $R_{1,i}$, so their
sum is invariant under $r$-conjugation. Multiplying this sum by
$q$ is integral by the same adjacent-exponent calculation; hence
the values of its phase on the two tuples coincide by \path{psi0}.
This proves local character compatibility. Take the product over
$w$ for the character of the full group, whose membership is tested
componentwise. In the last branch take $r=\kappa=1$, giving the
identity conjugator at every $w$. See
\path{DoubleCosetRepsConjugateLevelGroups.lean}:14--55,65--187 and
\path{pi_conj}, \path{idx_match}
(\path{ExistsConjReferenceChartIndex.lean}:14--26,87--135).

\paragraph{The middle chart on a fixed twisted coset.}
Fix $i<m$, a selected class $c_0$ with stage index $b$, and the given
adelic frame $k$. Write $C_w=\mathrm{bar}(r(c_0))_w$ at a generator
prime and $D_P=\mathrm{cqOf}(r(c_0))_P$ at a twist prime.
Take exactly the classes whose selected representatives have index $b$
and invariants $C,D$. This is the fibre in \path{ChExM.chart_of};
its decorated orbit group is
\[
 U_F\times\prod_w(\Lambda_{b,w}\cap\operatorname{Stab}(C_w))
       \times\prod_P(\Theta_{b,P}.H\cap\operatorname{Stab}(D_P)).
\]
The stabilizers preserve the indicated sets in both directions. The
rational action cancels in frame coordinates: under $(\gamma,h,t)$,
$C_w$ goes to $h_w C_w$, and a centre $f_P^{-1}y$ goes to
$(\gamma f_P t_P^{-1})^{-1}\gamma y=t_P f_P^{-1}y$.
Thus $D_P$ goes to $t_P D_P$.

Here is why these coordinates suffice even for the penultimate stage.
Write $q_P=|\mathcal O_E/P|$ and $M=\mathrm{cst}.M$. When
$i+1=m$ a point consists of a framed decoration and a translate
with centre $y$. At a generator $w$ its invariant is the singleton
$\{\beta_w(q_w y)\}$; changing the centre by $N_w$ leaves this
unchanged. Since $q_w\in w\setminus w^2$, two allowed centres
in $q_w^{-1}N_w$ yield the same singleton exactly when they differ
by $N_w$. Indeed the reduced-frame kernel is $wN_w=q_w N_w$.
At $P\in Q$, $D_P=f_P^{-1}y+q_P^M\mathcal O_{E,P}^{n}$.
For a fixed decoration, equality of these cosets is precisely that
the two centres differ by $q_P^M N_P$. These are all the defining
conditions on the difference of centres for equality of the global
translates. Even for distinct decorations the equality of the
full tuple of framed local cosets is exactly equality of $C,D$:
at $w$ use the lifted residue frames to identify the singleton with
$q_w$ times the centre modulo $w$, and at $P$ the two expressions
are literally the twist component. When $i+1\ne m$ every auxiliary
set is empty. See \path{CouplingCosetsInFramesEqIffReductions.lean}:19--129,
\path{StepPointDeterminedByLatticeInvariants.lean}:15--68 and
\path{CouplingCosetMovesWithFrameChange.lean}:14--47.

For any prescribed framed decoration $d$ with lattice $N'$ construct
a translate with the same $C,D$. Lift $\beta_w(q_w y)$ to a vector
$x_w\in N'$ by surjectivity of its prescribed residue map, and take
$a_w=q_w^{-1}x_w$. At a twist prime use
$a_P=f'_P f_P^{-1}y$, using the prescribed frame of $d$.
Produce one rational $y'$ with $y'-a_w\in N'_w$ and
$y'-a_P\in q_P^M N'_P$ at all marked primes. Indeed each $a_j$
is a scalar multiple of a global lattice vector, by clearing its
local denominator. Choose CRT integers $e_j$ congruent to
$\delta_{jv}$ modulo a sufficiently high power of each distinct
marked $v$. The powers can be chosen so that every
$(e_j-\delta_{jv})a_j$ satisfies the required lattice congruence
at $v$, including the $q_v^M$ factor. Then $y'=\sum_j e_j a_j$ works.
It is an admissible centre because $q_w a_w\in N'_w$ and
$q_P^M a_P\in N'_P$. Its generator reductions and twist cosets
have the required values by the preceding computations.
At the other stages use the empty set over $d$.
See \path{ExistsStepPointGivenReductionsCosets.lean}:18--94,136--246.

The chosen representative of an orbit preserves the normalization
needed for surjectivity. At $i+1=m$ the preferred coset value is
selected from the set of all coset values in that orbit. Two points
with the same $C,D$ at the same index have identical sets of possible
values, by the two coordinate action formulas above. They therefore
give the same preferred value. At every other middle stage it is
the constant empty tuple. Now lift any decoration $d$ with $C,D$
as just constructed, form its stage class, and move the point to
its selected representative. That representative has the preferred
value of $r(c_0)$ by normalization and therefore lies in the fibre.
The move from the lift to it preserves both invariants, and hence
its components belong to the two smaller stabilizers. Thus every
decoration orbit has a representative in the fibre. Conversely if
two fibre decorations are in the same restricted orbit, their
transporter preserves $C,D$. By the determination of the full point
proved above this transporter also identifies their stage points,
so the two classes coincide. The same argument with one class fixed
identifies the rational image of its stabilizer with all rational
parts of restricted triples fixing its decoration. Injectivity of
its rational matrix map follows from the supplied step law applied
to the difference of two stabilizer elements. See
\path{ChartClassesBijectiveOrbitsFramedLattices.lean}:52--98,
\path{ChartClassStabiliserRationalParts.lean} and
\path{ExistsFiniteOrderChartsMiddleSteps.lean}:268--297.

It remains to realize the restricted groups and characters rather
than merely count their orbits. At a twist prime, the index group's
congruence depth can be chosen larger than the sum of the denominator
bounds of the local defining tableau and its inverse, plus a depth
which makes the conjugated perturbation zero in the required
reduction. Extending $g-1$ by zero and conjugating multiplies the
entry bounds by those two bounds. Hence the extended matrix and its
inverse meet the prescribed local block conditions, with identity
reduction and trivial character (\path{LevelAtTwistPrimesIsOpen.lean}:16--70).
Cutting by $D_P=z+c\mathcal O_{E,P}^n$ retains a principal subgroup:
increase the depth until $(g-1)c^{-1}z$ is integral for both $g$
and $g^{-1}$. Then $c^{-1}(gx-z)=g(c^{-1}(x-z))+(g-1)c^{-1}z$
lies in the standard lattice precisely when $c^{-1}(x-z)$ does.
The empty cut imposes no condition. Complete this finite-depth
subgroup using its reduction classes at the selected split factor
and take the same character on their rational lifts; independence
follows because the deeper principal group is in the character kernel.
This is the local record $q$ in the statement. See
\path{OpenLevelCutByCosetStabiliser.lean}:71--102 and
\path{Formalism/LocalToGlobal/Columns/TwQBridge.lean}:96--109.

Use the unframed mirabolic level at $S_T$, the restricted generator
groups and the newly completed twist groups. Its global subgroup
$U_*$ is their simultaneous inverse image in the unitary group,
with character the generator type times the inverse product of the
twist characters. The open cut at finitely many split places gives
an open closed subgroup of the old compact level. Its local images
determine the global membership test as in
\hyperref[r13:s9-shrink]{the restricted level}, \path{exists_stepLevel_shrink}
(\path{ExistsStepLevelSmallerAtFinitePlaces.lean}:38--132).
A unitary whose components lie in the indicated local groups first belongs
to the uncut level by its product test. Each of the finitely many
twist conditions depends only on the component at its own base place,
so these components place it in $U_*$. The converse is restriction.
Off the marked places, insertion at one place with identity elsewhere
satisfies every additional cut, giving the old local group. For the
middle fibre this is applied to the cut level at $C,D$ in
\path{chartQ_exists} (\path{ExistsConjugatedChartDataLocalLevelAtQ.lean}:48--124). The base lattice
and level functionals persist because these places are away from $T$.
Each character has finite exponent: the generator quotient is finite
and the local character factors through the finite reduction at its
chosen positive depth. Take the product of their exponents.
The supplied $g_c$ realizes each decoration at the enlarged marked
set; forgetting extra frames gives its exact frame at all required
primes. Use \path{mem_framedLevel_iff_two_adelic_points} on the
stabilizer cuts in place of the original groups
(\hyperref[r13:s9-twopoints]{the two-point level test}; \path{MemFramedLevelIffTwoAdelicPoints.lean}:40--124).
Their finite congruence depth is the one in the present $q$.
For a new adelic class, \path{exists_K1_point_congruent}
(\hyperref[r13:s9-congruence]{prescribed congruences}; \path{ExistsK1PointCongruent.lean}:11--52)
puts a represented point in its level coset. Equality of two such classes
is the orbit test just proved. Thus \path{decoratedPointClasses_bijective_doubleCosets}
(\hyperref[r13:s9-cosets]{decorated classes and double cosets}; \path{DecoratedPointClassesBijectiveDoubleCosets.lean}:34--83)
uses exactly these $g_c$. The middle-fibre implementation makes this
identification with \path{recClsFy} and transports it by \path{chartQ_exists}. The
stabilizer range is the rational intersection with $g_cU_*g_c^{-1}$.
Equivalently these are \path{recClsFy} at \path{TwQExact.lean}:64
and the conjugated construction at
\path{ExistsConjugatedChartDataLocalLevelAtQ.lean}:48--124.

Finally replace $U_*$ by $kU_*k^{-1}$ and put $Y_c=g_c k^{-1}$.
The standard-lattice point hypothesis puts $k,k^{-1}$ in the
standard lattice stabilizer. The last-row condition at $P\mid\pi_T$
puts $e_n^{\mathsf t}k_P^{-1}=\overline{\pi_T}\,e_n^{\mathsf t}$
in the residue field. This scalar is a unit, so the mirabolic
functional stabilizer is normalized and the divisibility condition
on the last coordinate is preserved in both directions. Off the
places over $T$, the twist places, and (when $a\ne\mu_0$) the generator
places, the support hypothesis gives $k_v=1$ (\path{chart_of},
\path{ExistsFiniteOrderChartsMiddleSteps.lean}:314--321). The
local lattice and conjugacy tests then give the same base lattice
for the conjugated level. The chart character is $\chi_*(k^{-1}xk)$,
and $c\mapsto Y_c$ retains its double-coset bijection since right
multiplication by $k^{-1}$ identifies the two quotients. Its lattice
test follows from $Y_c^{-1}=k g_c^{-1}$ and preservation of the base
lattice by $k$. Its rational stabilizer remains the same because
$Y_c(kU_*k^{-1})Y_c^{-1}=g_c U_*g_c^{-1}$. This proves each
field of \path{ChExM.chart_of} on the specified fibre, with its
required $K.y(c)=g_c k^{-1}$, stabilizer and finite character exponent.
See \path{AdmissibleBasePointUnitaryNormalisesLevel.lean}:78--176 and
\path{ExistsFiniteOrderChartsMiddleSteps.lean}:314--400.

For the family, a chart label is by definition an occurring pair
$(b,\mathrm{cosetIn}(r(c)))$. Choose a preimage $c_0$ for each;
its fibre is the one just proved. Choose $g_c$ once for every stage
class by the shared decorated-point realization on the enlarged
marked set $I.SQ$. Then $Y_c x k=g_c(k^{-1}xk)$. If both sides
of the two-point test represent the specified points, the residue
and frame actions imply precisely a move with rational component
one preserving $C,D$, hence the original coupling coset. Conversely
such a move places $k^{-1}xk$ in $U_*$. For that move $(1,h,t)$
the generator reduction is $h^{-1}$ and the twist components are
$t_P^{-1}$ in the chosen frames; applying the inverse conventions in
both character factors gives $\sigma(\Lambda.\mathrm{chr}(h)
\Theta.\mathrm{chr}(t))$. The same $Y_c$ and $k$ satisfy the point,
support, last-row and stabilizer clauses. This gives
\path{GlobalTheory.T.exists_finiteOrder_charts_middle_steps} on
its supplied admissible $k$, with all the original finite-order
chart readings. The registered implementation is \path{ChExM.exists_charts}
(\path{ExistsFiniteOrderChartsMiddleSteps.lean}:450--499).

\paragraph{Comparison at an interior stage.}\phantomsection\label{r13:interior-charts}
To state \path{TC.ChTw.charts_tuple_trivQ}, fix the usual $I$ with
$\mathrm{Ok}$ and pins, the laws $l$, the later data $Z$, a word $u$,
tuple $a$, and reference charts $K^0$ with class map $c^0$ at step $i$.
Write $Q(a)=\{v:\mathrm{UnderQ}(a,v)\}$. For $Q$ finite, its assertion is
\[
\begin{aligned}
 &1\le i<m,\quad Q=Q(a),\quad \mathrm{ReadsK}_{\mu_0}(c^0,K^0)\\
 &\quad\Longrightarrow\ \exists\ \mathrm{Ch}\text{ finite},\ c:\mathrm{Cl}_a(i)\to\mathrm{Ch},\ K,\ 
                    g:\mathrm{Ch}\to\mathrm{Ch}^0,\ k,\\
 &\qquad \mathrm{Reads}_a(c,K,k),\qquad k_{\beta,v}=1\quad(v\in Q),\\
 &\qquad U_{\beta,v}=U^0_{g\beta,v}\quad(v\notin Q),\qquad
                 \forall\beta\ \exists n_\beta>0:\ \chi_\beta^{n_\beta}=1,\\
 &\qquad \chi_\beta(x)=\chi^0_{g\beta}(x)\quad
             (x\in U_\beta\cap U^0_{g\beta},\ x|_Q=1).
\end{aligned}
\]
Here $K$ is a family of level charts at the prescribed bad places,
and both evaluations on $x$ use the same underlying adelic element
(\path{ChartsTupleTrivQ.lean}:86--101).
For $0<i<m$, label the charts by the image of the index and reduction
coset. For the reference tuple, the identity conjugator compares a chart
with itself. For a twist, the matching-index theorem supplies residue
conjugators $x'_w$ and the reference-coset theorem supplies $c_0,y$
with $y$ in its type group. Put $x_w=y_w x'_w$. Conjugation by $x$
still identifies the groups, since $y$ conjugates the reference group
to itself. It leaves its character unchanged, since the character has
abelian target. At $i=m-1$ the $x'_w$ are one, so conjugation by $x$
matches the prescribed reduction coset by the defining property of $y$.
Lift the $x_w$ to a unitary adele $t_i$ preserving the standard lattice.
The admissible point of the new chart is the product of the reference
point with this lift by \path{admissibleBasePoint_chart_mul_twist}
(\path{AdmissibleBasePointChartMulTwist.lean:20--64}).
At a place under the twist, \path{underQ_prime_to_bad} excludes both
the level prime and the generator primes, so both factors there are one
(\path{TwistPlacesPrimeToLevelAndTypePrimes.lean:20--23};
\path{ChartsTupleTrivQ.lean:23--40}).
For a test element trivial at $Q_F$ the twist factor of a local move
must be one; the reference element then has type
$(1,xhx^{-1},1)$, and reversing the conjugation gives the inverse
membership test. Consequently the reading on the new chart and that on
the reference chart give equal characters on such elements.
For $v\notin Q_F$, insert a single local element at $v$ with identity
at all other places. This identifies their local subgroups by
\path{stepLevel_uloc_eq_of_agree}
(\path{StepLevelLocalGroupEqOfAgreeOffPlaces.lean:40--46};
\path{ChartsTupleTrivQ.lean:92--101}).
For each class the correctness field \path{Ok.exists_isPtOf} supplies
an adelic representative $g$ of its chosen decoration
(\path{StepZeroTwistReferenceTupleEqId.lean:217--228};
\path{LocalMoveClassRepIffLocalConditions.lean:325--337}). For a represented point
$(g,p)$ with rational component of the move one, equality of the
lattices of $p$ and $(1,h,\theta)p$ gives preservation of the standard
completed lattice by $s$ after cancellation of $g$. The pullback
coordinates for $gs$ are $s^{-1}$ times those for $g$. Choose rational
lattice vectors lifting a full residue basis (scale a local basis by
denominators prime to that place). Testing on them determines the last
row of $s^{-1}$ at the level primes and the entire reduction of $s^{-1}$
at the generators, namely $h_w$. The twist frame is exact, so
$s_P=\iota_P(\theta_P^{-1})$. Conversely these equations give all
five fields of a represented moved decoration for every representing
$g$. Include the original coset preservation separately in both
directions. The criterion is used in both directions by \path{match_idx}
(\path{MatchTupleChartMemberChart.lean:61--108}).
By \path{T.localConditions_conj_unitary_keepsStdLattice}, under
$s\mapsto t_i^{-1}s t_i$ unitarity and preservation
of the standard lattice are equivalent in both directions; the test
at $T$ is unchanged because $t_i=1$ there. At $w$ the inverse matrix
reduces by conjugation by $x_w$, in precisely the convention
$h_{0,w}=x_wh_wx_w^{-1}$. When $s$ is one over $Q_F$, the twist frame
test is equivalent to $\theta=1$ by
\path{T.twistLocalCondition_conj_iff_frameChange_trivial} and faithfulness of extension
$E\hookrightarrow E_P$. The forward and inverse conjugations thus
produce the local moves with the same character. For membership inferred
from the reading, a move is obtained by choosing a represented new
decoration: $s=1$ at the twist places, so the original rational twist
frames still represent the altered adelic point there; away from them
take its lattice and pullback residues. The reverse reading applied
to a known move likewise tests membership and the character on the
same point. These two implications are \path{locMv_of_mem} and
\path{reads_of_locMv} (\path{ChartLevelTwistTrivialLocalMove.lean:86--115};
\path{ReadsOfLocMv.lean:18--32}; \path{ChartsTupleTrivQ.lean:58--82}).
See \path{LocalMoveClassRepIffLocalConditions.lean:218--312,314--339},
\path{LocalConditionsConjUnitaryKeepsStdLattice.lean:186--257},
\path{TwistLocalConditionConjIffFrameChangeTrivial.lean:35--79}.
This establishes both implications of
\path{GlobalTheory.T.localMove_classRep_iff_localConditions}
(\path{LocalMoveClassRepIffLocalConditionsStatement.lean}:7--12).
The comparison uses the initial index conjugator and the twisted
middle-chart construction described separately.
See \path{match_idx} (\path{MatchTupleChartMemberChart.lean:17--110}),
\path{charts_tuple_trivQ} (\path{ChartsTupleTrivQ.lean:23--40,44--82,86--101}),
\path{exists_reference_class_matching_coupling_coset}
(\path{ExistsReferenceClassMatchingCouplingCosetStatement.lean:7--14};
10,738 lines), \path{exists_unitaryAdelic_lift_typePrime_reductions}
(\path{ExistsUnitaryAdelicLiftTypePrimeReductionsStatement.lean:7--13};
10,095 lines), and \path{ChExM.exists_charts}
(\path{ExistsFiniteOrderChartsMiddleSteps.lean:450--501}).
The nonreference branch of \path{twStep_mid} applies this comparison
to the reference interior reading in $\mathrm{StdF}$
(\path{TrueColumnChoiceTwistedChainsFormsClasses.lean:170--179}).

\paragraph{A single character on a given penultimate level.}
Fix $i=m-1\ge1$, finite charts $K_\beta$, their labels and comparison
map to the reference charts, and their tuple reading. Assume the
reference forms standard at every word with identity generator frames,
equality of local chart groups outside $Q_F$, and agreement of the two
chart characters on common elements trivial at $Q_F$. Take
\emph{any} coupled level $J$ for this data. It comes with an initial
character $\eta$ that agrees on each chart with the chart character on
the common elements. We construct a character on $J.U$ which has the
same agreement and also agrees with the reference level character
$\chi_0$ on every common element trivial at $Q_F$.
All comparisons below use the specified rank reindexing.

First $J$ and the reference level have the same base lattice. There
is a penultimate class: choose the index provided by final-index
uniqueness and take the base decorated point with its zero residue coset.
Using its chart $\beta$, the two readings test membership of every vector
against the same prescribed lattice. Now let $\rho(x)$ replace the
components of $x$ at $Q_F$ by $1$. This is a homomorphism on $J.U$:
it acts componentwise, preserves the unitary equation, and at a replaced
place $1$ belongs to the local subgroup. Reindex the same matrix to obtain
$\psi(x)$ in the reference level. At $Q_F$ this follows from the identity
component. Outside $Q_F$ and the generator places it follows from
$J.\mathrm{J\_off}$. At a generator place, the common lattice gives the
same lattice-level condition after reindexing, and $J.\mathrm{J\_g}$ and
the reference generator criterion select the same reduced matrix in the
prescribed symmetric subgroup. This proves $\psi$ is defined on all of
$J.U$. Use the character
\[
 \chi(x)=\chi_0(\psi(x))\,\eta(x)\,\eta(\rho(x))^{-1}.
\]
For an element $x$ also in chart $\beta$, $\rho(x)$ lies in the same
chart and, after reindexing, in its reference chart. This is true locally
because the two groups coincide off $Q_F$. At $Q_F$ both contain one.
The references and the assumed chart-character agreement give
$\chi_\beta(\rho(x))=\chi_0(\psi(x))$. The element $\rho(x)^{-1}x$
lies in both the coupled level and the chart, so $\chi_\beta$ on it
equals $\eta$. Multiplication yields $\chi_\beta(x)=\chi(x)$.
When $x$ is trivial at $Q_F$, $\rho(x)=x$, and $\psi(x)$ is the
reference copy of $x$; the two $\eta$ factors cancel. This proves
\path{exists_levelJ_character_eq_member_off_places}
(\path{ExistsLevelJCharacterEqMemberOffPlacesStatement.lean:7--8}).
Its witnesses are explicit in \path{cutL_prV}, \path{cutL_mem},
\path{twLevJ_chi_off_of_lb}, \path{twLevJ_chi_off},
\path{stepBefore_cl_nonempty_tup}, \path{twLevJ_exists_chi_off}
(\path{ExistsLevelJCharacterEqMemberOffPlaces.lean:19--75,141--236,240--273}).

\paragraph{Lemma. A direct calculation of the supported character.}\phantomsection\label{r13:near-character}
For a marked place $v\in Q_F$ write $e_v(x)$ for the element with
component $x$ at $v$ and identity elsewhere. Under the reading,
identity-frame and finite-exponent hypotheses below, every
$x\in K_\beta.U_v$ satisfies
$\chi_\beta(e_v(x))=\mathrm{val}_v(x)^{-1}$.
The source constructs \path{AdT.lev_mid} and uses its level and character
tests (\path{TrueColumnKernelDataCompatibleStepCharts.lean}:71--98;
\path{FurtherPlacePenultimateChartCharacterInverseValu.lean}:77--95,110--126).
Its intermediate level comparisons include
\path{exists_stepGroupElement_of_chartLevel},
\path{stepNearLifts_component_mem_localChartGroup}, and
\path{penultimateChart_localGroup_mem_nearLevelSet}
(\path{ExistsStepGroupElementOfChartLevelStatement.lean}:7--8;
\path{StepNearLiftsComponentMemLocalChartGroupStatement.lean}:7--8;
\path{PenultimateChartLocalGroupMemNearLevelSetStatement.lean}:12--13).
The following hand lemma calculates the character directly by rational
approximation at the single split place. These four source declarations
jointly remove 2,848 core lines in a cut at their declared statements.
Cutting \path{QChi.mem_nearSet_chart}
(\path{PenultimateChartLocalGroupMemNearLevelSet.lean}:134--158)
instead of its two-line wrapper leaves that wrapper reachable and
gives the 2,846 lines of review B.
Use the exact hypotheses of
\path{GlobalTheory.T.furtherPlace_penultimateChart_character_inverseValue}:
in addition to the chart and reference data above, the reading frame is one
over $Q_F$ and each chart character has finite exponent on elements
trivial at the generator places. Fix $v\in Q_F$, a chart $\beta$, and
$x\in K_\beta.U_v$. The reading is onto, so choose a class $c$ in this
chart with representative $r$. Write $Y$ for its adelic point, $k$ for
its reading frame, and $e_v(x)$ for extension by identity.
There is exactly one marked prime $P$ over $v$, and it is a twist
prime: in the quadratic CM extension the other prime is its conjugate,
and the full set of marked primes is disjoint from its conjugate
(\path{good_of_mem}, \path{levOk_twists}, \path{LevOkTwists.lean}:24--29,48--70;
\path{PlaceUnder.main}, \path{ExistsStarNe.lean}:21--35).
Both places over $v$ avoid $T$ and the generator primes. The denominator
$\pi_T$ is a unit there by $\pi_T\bar\pi_T=T$. Hence the local base
lattice read as $\mathrm{baseL}(\pi_T)$ is standard there.
The unitary local subgroup of the chart preserves it by
\path{StepLevel.Uloc_lev}. Its split projection $X=x_P$ and $X^{-1}$
are therefore integral.

Let $D$ be the congruence depth at $P$ and $e>0$ the above exponent
for this chart. Projection to the selected split factor identifies
the local unitary group with $\mathrm{GL}_n(E_P)$ in the source's
identity Gram coordinates; the other component is determined by
conjugate inverse transpose. Since the chart level is open, the
elements supported at $v$ and sufficiently near the identity belong
to it. Normalize $v_p(p)=1$ on $E_P$, and use the minimum entry
valuation on matrices. Choose a rational $s>1/(p-1)$ so large that
$v_p(A-1)>s$ puts the split unitary extension of $A$ in the chart.
Take $B=1+C$ so close to $1$ that it has the prescribed depth $D$
and $v_p(C)>s+v_p(e)$. The logarithm converges for $v_p(C)>0$ and
the exponential for $v_p(Z)>1/(p-1)$: the bounds
$v_p(k)\le(k-1)/(p-1)$ and $v_p(k!)\le(k-1)/(p-1)$ show that,
on the smaller ball used here, the terms after the linear one have
strictly higher valuation. Hence
$v_p(e^{-1}\log B)>s$ and $A=\exp(e^{-1}\log B)$ lies in the
chosen chart neighbourhood. The one-variable power-series identities
in the convergent algebra generated by $C$ give $A^e=B$.
This includes residue characteristic two and $p\mid e$.
Density of
$\mathrm{GL}_n(E)$ in $\mathrm{GL}_n(E_P)$ gives a rational $d$ with
$X^{-1}d=B=A^e$ in this neighbourhood. Both $d$ and its inverse
are $P$-integral. Extend $A$ unitarily at $v$ and by one outside $v$.
This extension belongs to the chart and is trivial at every generator
place. Its $e$th power has chart character one. Thus $e_v(x)$ has the
same character as its product $b$ with this $e$th power.
The $P$-component of $b$ is exactly $d$.

The reading gives a represented decoration for $Yk$, and $k_v=1$.
The adele $Ybk$ has the same lattice since $k^{-1}bk$ preserves the
standard local lattices everywhere. Its generator reductions and
level-prime functionals are unchanged since $b$ is one there. The
only changed marked twist frame is $\mathrm{fr}_P d$:
at $P$ the matrix equality follows from $k_v=1$, and this is still
a rational frame of the same lattice since $d\in\mathrm{GL}_n(\mathcal
O_{E,P})$. These data define a decoration by the defining frame axiom.
Apply the level reading to $b$ and this decoration. It supplies
a move with rational component one. Its type components must be
one by surjectivity of the unchanged reduction maps. At each twist
prime the moved frame is $\mathrm{fr}_{P'}\theta_{P'}^{-1}$, whence
$\theta_{P'}=1$ off $P$ and $\theta_P=d^{-1}$. The character reading
now gives $\chi_\beta(b)=\sigma_0(\lambda_P(d^{-1}))$, where
$\lambda_P$ is the local level character.

Since $\theta_P$ lies in the twist level group so does $d$.
The penultimate index has local parameter one by \path{last_h};
\path{grpQ} identifies this group with the projected local
automorphisms. A lift of $d$ in the defining middle-level relation
then has value $\lambda_P(d)$ by \path{chrQ}. Also $X^{-1}d$ has
depth $D$, so this is a near pair for $x$. If $d'$ is any other such
lift, $d^{-1}d'$ lies in the depth-$D$ principal group: write it as
$(X^{-1}d)^{-1}(X^{-1}d')$. The ideal of depth $D>0$ is stable under
products and inversion of these matrices. The index law makes
$\lambda_P$ trivial on this deep subgroup, so all near-pair values agree.
All other factors in $\mathrm{val}_v(x)$ are one by definition. Thus
\[
     \chi_\beta(e_v(x))\,\mathrm{val}_v(x)
       =\sigma_0\!\left(\lambda_P(d^{-1})\lambda_P(d)\right)=1 .
\]
This proves the claimed local identity (25,883 core lines). Source checks:
\path{FurtherPlacePenultimateChartCharacterInverseValu.lean:71--126},
\path{ExistsStarNe.lean:21--35}, \path{ExistsLevelElt.lean:42--100,106--145}
(two-sided rational approximation and the finite-exponent correction),
\path{ReferenceCharactersFactorTame.lean:262--278} (the five fields of
a represented point), \path{LatAtLattice.lean:256--280}
(the level and character readings), \path{CongAtOfNearSame.lean:43--104}
and \path{StepGroupCharacterInverseTypeAtLifts.lean:101--131}
(depth invariance and the sign).

\paragraph{The character and finite indices on the generated level.}
At $v\in Q_F$ let $W_v$ be the subgroup generated by all local chart
groups, after the specified rank reindexing. For the unique $P$ above
$v$, the integral matrices congruent at depth $D$ to some rational
$g$ in the projected local twist level form a subgroup. Indeed the
depth subgroup is normal in $\mathrm{GL}_n(\mathcal O_{E_P})$, so
products and inverses are represented by $gg'$ and $g^{-1}$, respectively.
The assignment $g\mapsto\lambda_P(g)$ factors through this reduction
group by depth invariance just proved. Every chart group lies in its
inverse image by the construction of $d$. Therefore
$\sigma_0(\lambda_P(g))^{-1}$ defines one character on $W_v$
extending all local chart characters. Take the product over $v\in Q_F$.
Any chart element supported in $Q_F$ factors as the product of its
single-place extensions, all in the chart by its componentwise level
test. Applying the local result to each factor proves
\path{GlobalTheory.T.exists_character_allCharts_furtherPlaces}
(\path{ExistsCharacterAllChartsFurtherPlaces.lean:19--71}).

\paragraph{The finite local index in the penultimate coupling.}
Fix the charts, readings, $v\in Q$, and the common full base lattice
in \path{FinAtQ.chartLoc_finiteIndex}
(\path{ChartLocalGroupFiniteIndexAtFurtherPlace.lean}:195--209).
Let $L$ be its completion at $v$, with the rank reindexing understood.
The base-lattice test from the reference chart agrees with the one
from every supplied chart: apply both readings to each vector, using a
class at the penultimate stage to identify the reference base.
Local chart groups are therefore unitary stabilizers of $L$ or subgroups
of that stabilizer, by the chart-level inclusion. Their join $W$ also
stabilizes $L$. See \path{chart_Uloc_le_levelVL_base}
(\path{ChartLocalGroupLeBaseLatticeLevel.lean}:20--34).
For one chart $H$, the inverse image of its open adelic subgroup under
the continuous embedding of local unitaries supported at $v$ is the local
level: at all other places the component is $1$. In coordinates some
positive integer $M$ thus tests membership by entrywise congruence
to $1$ (\path{exists_princ_loc_step}, \path{ExistsPrincLocStep.lean}:10--35).
Indeed the powers of the maximal ideal form a neighbourhood basis at zero
in each of the finitely many matrix entries, and a sufficiently high
power of the rational residue characteristic is contained in all the
specified ideals.
Choose a nonzero positive integer $d$ clearing the denominators in
both inclusions between $L$ and the standard coordinate lattice $L_0$,
so $dL_0\subset L$ and $dL\subset L_0$. For $N=Md^2$, triviality
on $L/NL$ implies $(g-1)dL_0\subset Md^2L\subset MdL_0$.
This gives the required entrywise congruence. The additive group $L/NL$
is finite: $L$ is finitely generated over the local integer ring of
$F_v$, whose quotient by the nonzero integer $N$ is finite (and is
zero if $N$ is a unit). Thus $W$ acts on a finite set with the
kernel contained in $H$. The index $[W:H]$ is finite, without a
normality assumption on $H$. See \path{princ_in_chart_gen},
\path{chart_contains_princ} (\path{ChartLocalGroupContainsPrincipalCongruence.lean}:13--55),
\path{finiteIndex_of_trivOn}, \path{finite_quot} and
\path{finiteIndex_of_princ} (\path{ChartLocalGroupFiniteIndexAtFurtherPlace.lean}:70--90,100--173,178--189).

For a nonempty family, the principal congruence subgroup of any chart is also contained in $W_v$.
For the empty family use the full unitary lattice level and the trivial
character, as the asserted chart conditions are vacuous. This gives
\path{GlobalTheory.T.furtherPlaces_group_contains_chart_levels}
(\path{FurtherPlacesGroupContainsChartLevels.lean:18--32},
\path{ChartLocalGroupContainsPrincipalCongruence.lean},
\path{ChartLocalGroupFiniteIndexAtFurtherPlace.lean}).

\paragraph{The penultimate coupling on these charts.}
Use the supplied $F$ and its $\mathrm{StdF}$ and identity generator-frame
hypotheses throughout. Its penultimate joint level, in the reindexed
coordinates, gives a Hermitian level $H_0$. Indeed its base lattice
is read on any penultimate reference class and equals
$\mathcal O_E^{n-1}\oplus\pi_T\mathcal O_E$. The standard-lattice and
level-prime functional tests for this lattice, and the generator
principal level, identify its local groups under reindexing: these are
the lattice and functional arguments of supplier 2 repeated for this new level
(\path{StepBeforeClNonempty.lean:28--60}). Set $J_v=W_v$ at
$v\in Q_F$ and $J_v=(H_0)_v$ elsewhere. Take the unitary adelic
elements with components in every $J_v$. The base lattice, its dual
and index come from $H_0$. Every additional local group preserves
this lattice and contains the positive principal level proved above.
The generator and level-prime conditions are those of $H_0$ since
$Q_F$ avoids those places. This constructs the required level
(\path{ReGLMemLevelVL.lean:84--141}).

For each reference chart $\alpha$ choose a point by surjectivity of
its supplied reading. Transport the singleton reduction of its
penultimate coset to the base residue field along the split factor
at each generator $v$, obtaining a vector $w_{\alpha,v}$. Within
the reference joint group, membership in the chart group at $v$
is precisely the fixing of this vector. To see it, extend a local
element by identity at all other places. Its lattice and level-prime
tests hold since $T\ne g_j$. The given joint-level generator test
requires that its reduction belong to the symmetric group prescribed
by $\mathrm{HSym}$, and the last index has parameter one and no
additional coordinates. The defining symmetric flag conditions say
that the lower entries for each flag, for the matrix and its inverse,
vanish and the designated leading entries have the prescribed identity
rows. Lift the reduced matrix to an integral invertible matrix.
These entry tests and their inverse tests imply preservation of the flag
lattices and the specified vectors modulo the local uniformizer;
hence the lift belongs to the defining sharp generator level.
This is the inclusion needed to apply the class reading. Its remaining
condition is invariance of the framed coset. At the selected generator
it is equivalent to fixing $w_{\alpha,v}$ by the singleton/uniformizer
calculation above, and at all other generators it holds for the identity.
Off the generators the chart level equals the joint level: both are
the full lattice level outside $T$ and the common functional condition
at $T$ gives the same local kernel there.
See \path{LevelMemChartLevelIffReductionFixesTypePrime.lean:16--132},
\path{SymmetricFlagBasisVectorEqSingleOfNoExtra.lean:241--274} and
\path{ChartLevOffType.lean:16--99}. This uses a supplied reference
reading and level in $F$, independently of their construction.

Set $G=J.U$ and let $H_\beta$ be the inverse image in $G$ of the
chart group after reindexing. Each chart group is contained in the
reindexed $G$, by construction at $Q_F$ and the reference comparison
elsewhere. It shares its base lattice by the two readings. At every
generator, reduction of the integral matrix in the given frame
is multiplicative, as is reduction of its inverse. Thus $G$ acts on
the product $V$ of all the finite residue vector spaces.
For $s\in G$, membership in $H_\beta$ requires precisely
$s_v\in K_{\beta,v}$ over $Q_F$ and
$s\,w_{g(\beta)}=w_{g(\beta)}$ in $V$. Outside $Q_F$ this uses the
off-place comparison hypothesized for the charts. Put
$Y_b=\coprod_\beta G/H_\beta$ with the identity coset in each
summand. Map the coset of $s$ to its vector $s w_{g(\beta)}$ and to
$(\beta,(s_vK_{\beta,v})_{v\in Q_F})$. Both maps are well-defined and
jointly injective: equality in their targets gives
$s^{-1}s'\in H_\beta$ by the just-established criterion. Those targets
are finite by the preceding index proof. Left multiplication gives
the required action, the unique orbit per label and its stabilizer;
the vector projection is equivariant and the labelled local cosets
are fixed by any element trivial at $Q_F$.

Finally use $\chi_0(\rho(s))\prod_{v\in Q_F}\psi_v(s_v)$ for the
initial character, where $\rho$ is the off-place restriction in the
reference joint level, with the rank reindexing. For a chart element
the factors agree with its character on the pieces trivial on
$Q_F$ and supported there, respectively, by the supplied off-place
agreement and the product-character proof above. The local
membership and action fields of the joint record therefore hold
on the same $Y_b,J$. Specifically, $J$ gives \path{levJ};
the finite disjoint union with its action and distinguished cosets gives
\path{Yb}, \path{YbFin}, \path{YbAct}, \path{yb}, \path{yb_orbit}.
The common lattice, inclusion and coset test give \path{lb}, \path{le},
\path{stab}. The character and its restriction give \path{chiJ} and
\path{chiJ_eq}. The computed local levels give \path{J_g} and
\path{J_off}. The generator vector and its transformation give
\path{yG}, \path{yG_smul}; the labelled product of further-place
cosets and its projection give \path{YQ}, \path{yQ}, \path{y_inj},
\path{yQ_off} (\path{TwLevJ}, \path{ChartWiden.lean}:46--86).
This proves
\path{GlobalTheory.T.exists_penultimateLevel_couplingCosets}
(\path{ExistsLevJCosetsOfFurther.lean:19--47,92--142,161--198,244--256};
\path{ReGLMemLevelVL.lean:160--230};
\path{ExistsPenultimateLevelCouplingCosets.lean:16--19}).

\paragraph{Assembling the given reference data with the twisted stages.}\phantomsection\label{r13:stage-pack}
For $1\le i<m$ use the initial index match
\path{exists_conj_referenceChartIndex}, followed by the coupled-coset
match and its supported adelic lift proved above. The common frame
is the reference frame times this lift. It is the identity over $Q_F$
since neither factor is supported there, preserves the standard
lattice and has the reference last row at $T$. Now apply the fixed-fibre
middle-chart construction \path{ChExM.chart_of} with this frame to all
occurring labels. Its finite-order reading, its representatives and the
character comparison already proved are for the same charts. The label
set is the finite image of the stage class set. For $a=\mu_0$, take the
supplied reference charts themselves, allowing the enlarged set of bad
places in their type. At $m$ take the top chart proved above, with its
specified common frame; at $\mu_0$ use the supplied singleton top chart.

Each local quotient functional can be chosen killed by $\pi_T$
without changing these readings. In the diagonal base lattice model,
use the canonical quotient on the dual given by the residue of
$\bar\pi_T x_n$ at the selected prime. It has cardinality $T$, and
multiplication by $\pi_T$ kills it. A nonzero old functional with
target of prime cardinality $T$ and annihilating the base lattice
factors through one of the two simple residue quotients of the
completed discriminant module: the complementary idempotents in the
split algebra act on a one-dimensional additive group. In the first
case the induced isomorphism of the quotient identifies its stabilizer
with that of the chosen functional. In the conjugate case unitarity
identifies the action on the complementary line with its conjugate
inverse, so pointwise stabilization of either line is equivalent.
This explains why replacing the functional preserves the local
subgroup in \path{swapLaws} and \path{killK}
(\path{InclMulInvMul.lean:227--242,312--401};
\path{TrueColumnChoiceTwistedChainsFormsClasses.lean:83--139}).
Transfer to the Hermitian chart via the rank reindexing: lattice
membership, left multiplication, double cosets and rational stabilizers
are transported by the same bijection on coordinates. At $Q_F$ the
chart contains a deep principal lattice congruence by
\path{T.open_congruence_subgroup_preserves_principal_lattice}
(\path{OpenCongruenceSubgroupPreservesPrincipalLatticeStatement.lean:7--8};
\path{StepQOfReindex.lean:41--53}). Off
$Q_F$ the comparison with $H$ follows from the previous one with $F$
and the given $\mathrm{CarH}$. This supplies the same-carrier and
local/base-lattice fields (\path{StepQOfReindex.lean:30--63}).

At $i=0$ use the given class equivalence $e_0$ to transport the reference
label. For $i>m$ forget the twist frame and take the class of the same
decoration with the empty marked coset at the unique reference top index.
This function on classes witnesses
\path{GlobalTheory.T.nonempty_classMap_above_lastStep}
(\path{NonemptyClassMapAboveLastStepStatement.lean}:7--8).
These choices satisfy the label requirements outside the chart range
(\path{NonemptyClassMapAboveLastStep.lean:7--16}).
Choose the family with the distinguished top entry. Its $m-1$ entry
supplies the precise inputs to the penultimate construction. The required
finite-exponent hypothesis there follows for nonreference tuples from
the middle finite-order character, and for the reference tuple from
character triviality off the generator places since $Q_F=\varnothing$.
The corrected character on that coupled level retains all its point and
stabilizer data, with the off-place reference agreement. This yields
one stage package for every $u,a$.

Take the image of $Q_E$ under contraction as the finite prescribed set
for each tuple. In $T$ retain the stage data, their lattice and
stabilizer identifications, and the finite penultimate orbit record.
At zero both stabilizers are trivial by the supplied freeness law,
so $e_0$ gives the class map with its unique stabilizer isomorphism.
In $J$ use the comparison of labels, local levels outside $Q_F$ and
base lattices just proved. For the place and residue comparisons of $W$
use the identity on the selected level primes and on their residue
fields. The diamond scales only the chosen residue functional, so it
commutes with the identification by the underlying decoration at stage
zero. If two moves agree on that decoration, their discrepancy in the
twist frames can be removed by an integral twist factor: at zero the
twist group is the full group and has trivial character. At this stage
\path{ClsZero.ys0} makes the marked set empty. Its image under the
rational action is again empty, so the corrected move fixes the whole
decorated point. Freeness makes its stabilizer trivial
(\path{StepZeroClassesBijectiveSameDeco.lean:110--111};
\path{DecoStabTrivial.lean:57--73}). Thus the rational and generator
components of these transports coincide, giving equality on functions
of frames; the stage-zero equivariant spaces reduce to the classwise
polynomial condition and transport along $e_0$. This gives the bottom
diamond and form-space fields of \path{twAtT_exists}, together with
the prescribed residue map (\path{DecoStabTrivial.lean:23--73,88--174}).

With $I$, its correctness and pins, the laws $l,l_0$, later data $Z$,
a predicate $Q$, and the reference columns and forms $H,F$ fixed,
\path{TC.tw_chains} states (\path{TrueColumnChoiceTwistedChainsFormsClasses.lean}:464--484)
\[
\begin{aligned}
 &Q(a,v)\iff\mathrm{UnderQ}(a,v)\quad(\forall a,v),\qquad \mathrm{CarH}(H,F),\\
 &\mathrm{StdF}(F_u)\land\mathrm{StdLvC}(F_u)\quad(\forall u),\qquad
                      (F_u).\mathrm{glevJ.fr}_{j,v,h}=1\quad(\forall u,j,v,h)\\
 &\quad\Longrightarrow\ \exists\ T,J,W,\quad \mathrm{TwChMu0}(T)\land\mathrm{StdW}(W)
              \land\mathrm{Nonempty}(\mathrm{SymmetricTwistedJointForms}(T,F)).
\end{aligned}
\]
Here $T$ lies over the given $Q$, $J$ joins $T$ to $H$, and the twisted form
record $W$ lies over $T,J,F$ for the given $l_0$.
From the displayed conjunction \path{tw_chains} uses the
$\mathrm{StdF}$ projection. Its generator-frame identity is supplied
separately. \path{tw_chains_std} obtains that identity from
$\mathrm{StdLvC}$ and applies \path{tw_chains}
(\path{TrueColumnChoiceTwistedChainsFormsClasses.lean:464--472,488--489}).
For every active stage take the valuation, rational and integral
form-space, metric and diamond comparisons on its actual chart and frame
identification. The proofs for arbitrary tuples under supplier 1
(\path{TC.val_tied}, \path{MetricTie.ipV_norm_tie},
\path{DiaChart.dilTied}) use the identity-on-matrices frame
identification. The metric comparison needs only that identification
under its ambient hypotheses. The valuation and diamond comparisons
also use the reading; the diamond uses the killed functional and place
identification as well. The endpoint or interior reference
frame in $F$ gives that identity-on-matrices property by $\mathrm{StdF}$.
Consequently these applications require no existence assertion for $F$.
Together with the local character comparisons and their finite exponents
these are the fields of $W$. The symmetric joint fields are the
generator action, membership, injectivity and $Q_F$ projections on the
same finite set $Y_b$. The distinguished reference class equivalence is
the identity on already selected decoration classes by
\path{T.eCls0_mu0_eq_self} (\path{StepZeroClassEquivReferenceTupleEqSelfStatement.lean}:7--8).
For the field \path{twisted_e0_at_mu0_eq_eMu0} of $W$,
\path{T.stepZeroFormsIdentification_referenceTuple_eq_self} applies
the injective stage-zero form embedding. After embedding, the chosen
operator acts by reindexing through that same class equivalence, hence
fixes every value (\path{StepZeroFormsIdentificationReferenceTupleEqSelf.lean:9--21};
\path{TrueColumnChoiceTwistedChainsFormsClasses.lean:435--445}).
The direct statement wrapper of \path{tw_chains} is
\path{T.trueColumnChoice_twistedChains_formsClasses}
(\path{TrueColumnChoiceTwistedChainsFormsClassesStatement.lean:8--9}).
Finally retrieve
the single top frame in each nonreference package: its reading and
extension property are exactly the two clauses required for this $W$.
This proves the original supplier
\path{GlobalTheory.T.trueColumnChoice_exists_chainData_with_basePoints}
under its stated hypotheses.
See \path{twPack_ex}, \path{twT}, \path{twJ}, \path{twW}, \path{twJS},
\path{tw_chains} (\path{TrueColumnChoiceTwistedChainsFormsClasses.lean:272--319,325--485})
and \path{tw_ties_with_points}
(\path{TrueColumnChoiceExistsChainDataWithBasePoints.lean:13--23}).

% blueprint/ch_P_r12.tex lines 4287 through 4373; r12:assembly
\subsection{Simultaneous use of the two comparison records}\label{r12:assembly}
Fix the ambient $I,l,Z$, pins and hypotheses of Z4.4. We now assemble the suppliers on the witnesses constructed above. Supplier 1 supplies
one chart family and forms $F_u$ with $\mathrm{StdF}$ and $\mathrm{StdLvC}$
simultaneously, one generator datum for all words and the reference
character. Supplier 2 takes those very charts and forms. At each word
supplier 3 replaces its zero level, preserving the base lattice, local
and global subgroups, representatives, stabilizers, splitting and
generator frame. Substitution in the carrier comparisons $\mathrm{car0},
\mathrm{sp},\mathrm{fr0}$ therefore gives $\mathrm{CarH}(H,F)$ for the
new collection. Supplier 4 acts on this $H,F$; write $A$ for its witness.
Its equality of uniformizers turns supplier 3's annihilation into the
precise input for supplier 5. The latter gives the cofactor, base and
residue compatibilities and $X_0$ with common generator; supplier 6
gives $\mathrm{StdTied}$. The two word-independence identities for the
frames follow by writing them in the given standard coordinates.
These choices give all fields of a single $C$ while $C.\mathrm{forms}=F$,
so both original standardness assertions survive. This is the joint
construction in \path{colsTiedSym_std}
(\path{ExistsStandardColumnComparisonsTrueColumns.lean}:20--72).

With $t_p.\xi=I.\xi$, membership in $t_p.\mathrm{InQ}'(a,v)$ is
equivalent to the ramification of $I.\xi(a)$ at a place over $v$,
hence to $I.\mathrm{UnderQ}(a,v)$. Supplier 7 takes this equivalence
and the carrier and form data extracted from that $C$; write $T,J,W$
for its witnesses, including the common top point. Supplier 8 takes
this $T,J,W$ and supplies $L_v$ and the reference comparison. Supplier 9
uses this $L_v$ and the same top point to give the point, evaluation
and compatibility records. Assemble them without reselecting a chart,
a level or a representative. Only on that common record does the
class sum provide a single period value with the stated bound at
every allowed place. See \path{underQ_iff}, \path{carH_of_cols},
\path{colsTwPTiedS_std} (\path{TrueColumnChoiceExistsStandardChainData.lean}:19--56)
and \path{twistedPeriodAdmissible_of_rowInt}
(\path{AdmissiblePeriodTwistOfIntegralRows.lean}:14--21).

\paragraph{One field element for all allowed places.}
Fix a used nonreference tuple $\mu$, a divisor datum, an integral point
$x$ and a permitted form. At each embedding, positive definiteness of
the primed space gives a frame with Gram matrix $1$. Taking these frames
together gives \path{exists_orthonormal_frame_twisted} by \path{D7.X'_posDef}
(\path{TwistedPeriodDecompositionOfClasswiseTwistedPair.lean}:22--54;
\path{PosOfIsomEmb.lean}:176--177).
In this frame \path{eW_rep_o_tw} represents the value at each class by
a polynomial in the complex weight space. Changing frame acts by block
row substitution and a scalar. Stability under substitution follows from
\path{blkRow_mem} and \path{weyl_eq_dual}: over $\mathbb C$ the two
spaces have the same span, by scalar extension of the full rational
lattices in Section~\ref{r12:weyl}. See
\path{TwistedPeriodDecompositionOfClasswiseTwistedPair.lean}:68--86,
\path{TiedBrRep.lean}:33--46 and
\path{BranchingNormalisedOfAdjoint.lean}:119--121,229--233.
Restricting the resulting finite family to any one class keeps that
polynomial there and gives zero in the same weight space at each other
class. These restrictions are complementary coordinate projections. In the weighted inner
product all products between distinct components vanish. Therefore
the period of the embedded point is the sum of the same-class
contributions. At each $c$ the row-integrality hypothesis supplies
the localized row value and the integral point supplies the opposite
coefficient lattice. Test these in the two dual frames above. The
normalized contraction is integral at every allowed $w$, or zero;
multiplication by $h_c=|\mathrm{Stab}(c)|$ cancels the class weight.
For the field-valued numerators $b_c$ supplied by the row and point
comparison, this yields the single element
\[
                 b=\sum_c h_c^{-1} b_c .
\]
Under the fixed complex embedding its image is the required period,
by the coordinate sum. Each $h_c$ divides $D_{\Gamma,u}$ under the
denominator hypothesis of admissibility by \path{card_stabilizer_dvd₂}.
At a permitted place $w\mid p$, \path{vAdd_dvd_le} therefore bounds
$v_w(h_c)$ above by $e_w v_p(D_{\Gamma,u})$
(\path{AdmPerTw.lean}:14--40,42--47). The nonzero numerator has
$v_w(b_c)\ge0$. Applying \path{OrbitPlaces.vAdd_mul} and
\path{OrbitPlaces.vAdd_inv} gives the lower bound
$-e_w v_p(D_{\Gamma,u})$ for each nonzero summand
(\path{Formalism/SharpRestriction/Divisibility/OrbitPlaces.lean}:545--550,1377--1381).
Induction on the finite class set with \path{WeightDegreeSizes.pair_add_ge}
gives the same bound when $b\ne0$, with zero as the empty sum
(\path{WeightDegreeSizes.lean}:737--763).
The representatives and $b_c$ were fixed before choosing $w$.
See \path{pair_integral}
(\path{ClassPairingIntegralFormRowIsIntegral.lean}:49--69),
\path{perDecTwT_of_std} (\path{PerDecTwTOfStd.lean}:11--44),
\path{petersson_eq_sum_restrict_twist}, \path{perDecTw_of_at}
(\path{TwistedPeriodDecompositionOfClasswiseTwistedPair.lean}:96--132,141--168),
and \path{perTwAt_of_perDecTw}
(\path{AdmissiblePeriodTwistsReductionPlaceTuples.lean}:12--37).

% blueprint/ch_P_r12.tex lines 4414 through 4457; r13:recheck, r14:z44-approved
\subsection{Round 13 corrections for rechecking Z4.4}\label{r13:recheck}
The two assembly records and their conditional implication for the
original root have been checked. The independent recheck recorded the following
seven corrections as fixed in \path{out/R13_recheck.md}.
\begin{enumerate}
\item \hyperref[r13:weyl-dual]{Transported Hermitian duality}.
Verify both inclusions in both dual-lattice equalities, including the
conjugate place in the first slot and the specialization in the period numerator.
\item \hyperref[r13:nabdelloc]{The two uses of localization}.
The coefficient tests use \path{ValLat.mem_localization_delAtr} via
\path{Formalism.L2.DIV.nabDelLoc}. The transported pairing has a separate
caller in the root cone.
\item \hyperref[r12:z12-stage]{The seven local hypotheses in Z1.2}.
The constructor of \path{HypothesesH'} passes its literature record to
the caller. \path{locPrints_of_litN} projects the seven fields of that
record.
\item \hyperref[r13:move-fields]{Upward and downward moves},
\hyperref[r13:fv-field]{the value-module isomorphism} and
\hyperref[r13:lv-zero]{the step-zero comparison}. Check the common
rational and adelic witnesses and both residue domains against the
formulas for the moves.
\item \hyperref[r13:stage-unit]{The stage and base law block}.
Check the isolated use of the unitary equation in \path{step_gam_mem_UJ}.
\item \hyperref[r13:tuple-independence]{Tuple independence at a generator}
and \hyperref[r13:near-character]{the character on an element supported at a
twist place}. The ledger records the two alternative arguments, their
hypotheses, and the Lean results they bypass.
\item For the primed chart, the top chart and the fixed middle fibre, check
the use of \hyperref[r13:s9-values]{the value and metric argument},
\hyperref[r13:s9-twopoints]{the two-point test},
\hyperref[r13:s9-congruence]{the prescribed-congruence construction},
\hyperref[r13:s9-cosets]{the orbit and double-coset argument} and
\hyperref[r13:s9-shrink]{the smaller level} at the four revised passages.
\end{enumerate}

\paragraph{Round 14 approval of Z4.4.}\phantomsection\label{r14:z44-approved}
The ten newly registered declarations are one transported-duality result
and nine chart, chain or period results. Their individual names, Lean
locations and proof sections are in \path{round13.new_subtheorems} in
\path{status.json}. The recheck confirmed the proof of \path{TransDual.nabDelDual}.
The other nine matched the source statements and await proof-level review.
The three elementary literature-class leaves appear separately in the
trust ledger. Following these edits the original assignment Z4.4 has
approved grade \emph{complete}.

